\documentclass[12pt,oneside]{book}

\newcommand{\universidad}{Universidad de Buenos Aires (UBA)}
\newcommand{\materia}{Derecho Procesal Civil}
\newcommand{\titulo}{Derecho Procesal Civil}
\newcommand{\subtitulo}{Manual de estudio integrado}
\newcommand{\autor}{Isaac Edgar Camacho Ocampo}
\newcommand{\anio}{2026}

\usepackage[utf8]{inputenc}
\usepackage[T1]{fontenc}
\usepackage[spanish,es-nodecimaldot]{babel}
\usepackage[a4paper,top=2.5cm,bottom=2.5cm,left=3cm,right=2.5cm,headheight=15pt]{geometry}
\usepackage{amsmath}
\usepackage{amssymb}
\usepackage{mathtools}
\usepackage{graphicx}
\usepackage{float}
\usepackage{caption}
\usepackage{subcaption}
\usepackage{booktabs}
\usepackage{array}
\usepackage{longtable}
\usepackage{tabularx}
\usepackage{enumitem}
\usepackage{xcolor}
\usepackage[most]{tcolorbox}
\usepackage{fancyhdr}
\usepackage{ragged2e}
\usepackage[colorlinks=true,linkcolor=blue!50!black,citecolor=green!40!black,urlcolor=blue!60!black,bookmarks=true,pdftitle={Derecho Procesal Civil: manual de estudio integrado},pdfauthor={Isaac Edgar Camacho Ocampo},pdfsubject={Derecho Procesal Civil}]{hyperref}

\setcounter{tocdepth}{2}
\setcounter{secnumdepth}{2}
\setlength{\parindent}{0pt}
\setlength{\parskip}{0.5em}
\clubpenalty=10000
\widowpenalty=10000
\setlist[itemize]{topsep=0.3em,itemsep=0.2em}
\setlist[enumerate]{topsep=0.3em,itemsep=0.2em}
\setlist[description]{topsep=0.3em,itemsep=0.3em}
\renewcommand{\arraystretch}{1.15}
\sloppy

\newcommand{\abs}[1]{\left|#1\right|}
\newcommand{\blanco}[1][1cm]{\underline{\hspace{#1}}}

\newtcolorbox{definicion}[1]{enhanced,breakable,title={Definición: #1},fonttitle=\bfseries,colback=blue!3,colframe=blue!50!black,boxrule=0.8pt,arc=2mm}
\newtcolorbox{importante}{enhanced,breakable,title={Importante},fonttitle=\bfseries,colback=yellow!8,colframe=orange!70!black,boxrule=0.8pt,arc=2mm}
\newtcolorbox{ejemplo}[1]{enhanced,breakable,title={Ejemplo: #1},fonttitle=\bfseries,colback=green!4,colframe=green!40!black,boxrule=0.8pt,arc=2mm}
\newtcolorbox{regla}[1]{enhanced,breakable,title={Norma: #1},fonttitle=\bfseries,colback=purple!4,colframe=purple!50!black,boxrule=0.8pt,arc=2mm}
\newtcolorbox{nota}{enhanced,breakable,title={Nota},fonttitle=\bfseries,colback=white,colframe=gray!50,boxrule=0.6pt,arc=2mm}

% Cuadro Cornell: tabla real de dos columnas con fila final de resumen
\newenvironment{cornell}[1]{%
  \par\medskip\noindent{\small\bfseries Cuadro Cornell: #1}\par\smallskip
  \begin{longtable}{|>{\RaggedRight\arraybackslash}p{0.27\textwidth}|>{\RaggedRight\arraybackslash}p{0.66\textwidth}|}
  \hline
  \textbf{\small Preguntas / palabras clave} & \textbf{\small Notas / respuestas} \\ \hline\endhead
}{\end{longtable}\par\medskip}
\newcommand{\cqa}[2]{\textbf{#1} & #2 \\ \hline}
\newcommand{\cornellsummary}[1]{\multicolumn{2}{|p{0.95\textwidth}|}{\textbf{Resumen:} #1} \\ \hline}

\pagestyle{fancy}
\fancyhf{}
\fancyhead[L]{\nouppercase{\leftmark}}
\fancyhead[R]{\materia}
\fancyfoot[C]{\thepage}
\renewcommand{\headrulewidth}{0.4pt}
\fancypagestyle{plain}{\fancyhf{}\fancyfoot[C]{\thepage}\renewcommand{\headrulewidth}{0pt}}

\begin{document}

\begin{titlepage}
\thispagestyle{empty}
\begin{center}
\vspace*{1cm}
{\large \textsc{\universidad}}\\[3cm]
{\Huge \textbf{\titulo}}\\[0.8cm]
{\LARGE \subtitulo}\\[1.2cm]
\rule{0.70\textwidth}{0.5pt}\\[0.4cm]
{\large Apuntes de estudio}\\[0.4cm]
\rule{0.70\textwidth}{0.5pt}\\
\vfill
{\large \autor}\\[0.3cm]
{\large \anio}
\end{center}
\vfill
\begin{flushleft}
\textit{Este es un modesto aporte a los estudiantes de «\materia»; no pretende sustituir el material oficial de la materia.}
\end{flushleft}
\end{titlepage}

\frontmatter
\tableofcontents
\mainmatter

\chapter*{Introducción general}
\addcontentsline{toc}{chapter}{Introducción general}
\markboth{Introducción general}{Introducción general}

Este libro reúne en un único desarrollo el contenido de las clases de Derecho Procesal Civil. Parte del conflicto que da origen al proceso y de las nociones que sostienen la disciplina (Parte I); continúa con la organización judicial, la competencia, el debido proceso, los sujetos del proceso y los modos de resolver conflictos (Parte II); estudia los actos procesales, sus plazos, notificaciones y escritos (Parte III); sigue el curso del proceso ordinario desde las diligencias preliminares hasta la contestación de la demanda (Parte IV); y concluye con la audiencia preliminar y la prueba (Parte V). Cada capítulo cierra con un cuadro Cornell para el estudio activo.

\part{Fundamentos del Derecho Procesal Civil}

\chapter{El conflicto jurídico y el Derecho Procesal Civil}
\label{cap:conflicto}

Este capítulo responde a una pregunta previa a cualquier otra: ¿por qué existe el proceso? Parte de la distinción entre el derecho sustancial y el derecho procesal, describe la relación jurídica nueva que nace del conflicto, delimita el carácter autónomo y unitario de la materia y, por último, repasa las fuentes normativas y la evolución histórica que explican el sistema procesal argentino. El capítulo siguiente (\ref{cap:pedestal}) desarrolla las tres nociones que sostienen toda la disciplina.

\section{Concepto e importancia práctica}

El derecho procesal civil es la materia que permite al abogado desempeñarse en la práctica cotidiana: es la herramienta que hace actuar la voluntad de la ley cuando el derecho de fondo dejó de ser operativo entre las partes. Su estudio debe encararse procurando \textbf{comprender} sus nociones y no simplemente memorizarlas.

\begin{nota}
Como material de estudio se mencionan el libro de la cátedra \textit{Nociones Básicas del Derecho Procesal Civil} (en prensa) y, mientras tanto, las obras de Palacio, Falcón y Arazi, junto con el Código Procesal Civil y Comercial en su edición con leyes complementarias. Esa edición incluye, en su parte posterior, normas como la Ley 22.172 sobre exhortos a otras jurisdicciones del país. Las normas procesales se citan a lo largo del libro por su número de artículo; CPCCN designa al Código Procesal Civil y Comercial de la Nación.
\end{nota}

\section{El conflicto como origen del proceso}

\subsection{Del derecho sustancial al derecho procesal}

Debe partirse de una distinción fundamental: el \textbf{derecho sustancial} (o derecho de fondo) regula la vida ordinaria en sociedad, mientras que el \textbf{derecho procesal} entra en escena cuando surge el conflicto y las normas de fondo dejan de ser operativas entre las partes.

\begin{ejemplo}{El colectivo}
Una persona sube a un colectivo y celebra un contrato de transporte. El colectivo choca, el pasajero se golpea la cabeza y requiere una semana de internación. La empresa no responde. En ese momento la relación sustancial entre el pasajero y la empresa de transporte se agotó, y surge una nueva relación jurídica en la que aparece un tercer actor: el juez. Este nuevo vínculo se rige por el Código Procesal Civil y Comercial.
\end{ejemplo}

\begin{ejemplo}{La heladera}
Una persona compra una heladera que no funciona, la garantía no responde y el local cerró. La relación de consumo quedó insatisfecha. Para resolver el conflicto debe recurrirse al Tribunal de Consumo (Diagonal Sur 651, 1.º piso) o a la justicia ordinaria a través del proceso.
\end{ejemplo}

\begin{ejemplo}{El infortunio laboral}
Un trabajador manipula una máquina, hay una pieza rota, cae y sufre una conmoción cerebral. La relación laboral era sustancial; el accidente genera el conflicto que debe resolverse por la vía procesal correspondiente.
\end{ejemplo}

\subsection{La nueva relación jurídica procesal}

El juez, como \textbf{tercero imparcial}, dicta una sentencia que constituye una \textbf{norma individual}: una «ley a medida» que resuelve el conflicto entre las partes. La sentencia es el momento en que el derecho sustancial vuelve a ser operativo, pero ahora por mandato judicial.

\begin{importante}
El Código Civil y Comercial de la Nación rige la relación sustancial. El Código Procesal Civil y Comercial rige el desarrollo del proceso una vez surgido el conflicto.
\end{importante}

\section{Autonomía y unidad del Derecho Procesal}

El derecho procesal es una materia \textbf{autónoma}: no solo tiene código propio, sino también doctrina propia y autores especializados. Al mismo tiempo es \textbf{uno solo}, aunque existan especialidades, porque todas sus ramas se apoyan en el mismo \textbf{pedestal básico}: jurisdicción, acción/pretensión y proceso (capítulo~\ref{cap:pedestal}).

\begin{table}[H]
\centering
\renewcommand{\arraystretch}{1.25}
\begin{tabularx}{\textwidth}{>{\RaggedRight\arraybackslash}p{0.27\textwidth}>{\RaggedRight\arraybackslash}X}
\toprule
\textbf{Rama especializada} & \textbf{Rasgos señalados en las fuentes} \\
\midrule
Derecho Procesal Penal & Incluye figuras como el ministerio público en el rol de demandante, los recursos de nulidad, de casación e inconstitucionalidad, y medidas cautelares como la prisión preventiva. \\
\addlinespace
Derecho Procesal Laboral & Incorpora el \textbf{principio protectorio}: en caso de duda se resuelve a favor del trabajador, buscando equilibrar la asimetría de poder entre empleado y empleador. \\
\addlinespace
Derecho Procesal Constitucional & Estudia el amparo, la acción declarativa de inconstitucionalidad y el recurso extraordinario. \\
\addlinespace
Derecho Procesal Administrativo & Se vincula con el control jurisdiccional de los actos administrativos. \\
\bottomrule
\end{tabularx}
\end{table}

\section{Fuentes del Derecho Procesal}

Las fuentes del derecho procesal son las siguientes:

\begin{enumerate}
\item \textbf{Constitución Nacional:} especialmente el art.~18 (juez natural) y la parte orgánica sobre el Poder Judicial. El art.~31 CN se cita en relación con la supremacía constitucional.
\item \textbf{Tratados internacionales de derechos humanos:} el art.~75 incs.~22, 23 y 24 CN los equipara en jerarquía a la Constitución. La CSJN, en el caso \textit{Ekmekdjian c/ Sofovich} (anterior a la reforma de 1994), reconoció que los derechos consagrados en tratados internacionales no requieren ley reglamentaria para ser operativos.
\item \textbf{Leyes nacionales y locales} (provinciales).
\item \textbf{Decretos,} incluidos los Decretos de Necesidad y Urgencia (DNU). La Ley 26.122 establece el mecanismo de control parlamentario de los DNU mediante la Comisión Bicameral; si el Congreso no los deroga en determinado plazo, quedan aprobados con rango de ley.
\item \textbf{Jurisprudencia:} fallos de los tribunales, en particular de la Corte Suprema de Justicia de la Nación.
\item \textbf{Doctrina:} obras de autores especializados.
\end{enumerate}

\begin{ejemplo}{DNU durante la pandemia}
Los DNU dictados durante la pandemia de COVID-19 suspendieron lanzamientos y ejecutorios en el proceso laboral, incidiendo directamente en el derecho procesal.
\end{ejemplo}

\subsection{Dinámica de la norma procesal}

La norma procesal posee una dinámica especial: ante una misma situación puede aplicarse de una u otra manera.

\begin{ejemplo}{Notificación de la demanda}
Al notificarse la demanda (traba de la litis), el demandado puede comparecer y contestar, comparecer sin contestar, o no comparecer. Cada posibilidad genera consecuencias procesales distintas, todas reguladas por el Código (véanse los capítulos~\ref{cap:notificacion} y~\ref{cap:contestacion}).
\end{ejemplo}

\section{Breve historia del Derecho Procesal}

\subsection{El proceso romano y el derecho germánico}

El derecho procesal tiene raíces en el \textbf{proceso oral romano}, desarrollado en dos etapas: ante un magistrado (\textit{in iure}) y ante un juez o árbitro elegido por las partes (\textit{apud iudicem}).

Con las invasiones bárbaras, el derecho romano fue influenciado por el derecho germánico, que introducía el elemento de la divinidad en la resolución de los conflictos: las \textbf{pruebas de ordalías}. Un imputado podía ser sometido a caminar sobre brasas ardientes; si a los dos días no mostraba lesiones, se lo consideraba inocente. También se lo podía arrojar encadenado al agua: si sobrevivía, era inocente.

\subsection{El derecho común y su llegada a la Argentina}

En la Edad Media, la fuerte influencia de la Iglesia generó el llamado \textbf{derecho común}, que llegó a la Argentina principalmente a través de España, con las Leyes de Enjuiciamiento Civil de 1855 y 1881.

\subsection{Sistemas de justicia comparados}

\begin{itemize}
\item \textbf{Tribu Azande (norte de África):} se inocula doble veneno a un pollo; según muera o sobreviva, se determina la culpabilidad.
\item \textbf{Sistema anglosajón (\textit{common law}):} el jurado y el valor de la declaración del testigo son centrales; la prueba testimonial tiene mayor peso que la documental.
\item \textbf{Sistema continental (\textit{civil law}):} es el sistema al que pertenece la Argentina, donde la prueba documental y el expediente escrito tienen preponderancia. El testigo tiene un valor diferente al del sistema anglosajón.
\end{itemize}

\subsection{El Código Procesal de 1967}

El Código Procesal Civil y Comercial argentino fue sancionado en 1967, con fuerte influencia del derecho procesal italiano, y desde entonces ha sufrido numerosas modificaciones. Es un código escrito, con procedimiento formal y predominio del expediente en papel, aunque en proceso de incorporación de la oralidad y la digitalización.

\begin{cornell}{conflicto, autonomía, fuentes e historia}
\cqa{¿Qué es el Derecho Procesal Civil y para qué sirve?}{Es la herramienta que permite hacer actuar la voluntad de la ley cuando el derecho de fondo dejó de ser operativo entre las partes, es decir, cuando surgió un conflicto. Es la materia que permite trabajar al abogado.}
\cqa{¿Cuándo aparece el proceso y qué es la sentencia?}{Cuando la relación sustancial se agota por el conflicto y se necesita un tercero imparcial, el juez. La sentencia es la norma individual («ley a medida») con la que el derecho sustancial vuelve a ser operativo por mandato judicial.}
\cqa{¿Qué código rige cada relación?}{El Código Civil y Comercial rige la relación sustancial; el Código Procesal Civil y Comercial rige el proceso una vez surgido el conflicto.}
\cqa{¿Es uno o múltiple? ¿Es autónomo?}{Es uno solo (todas las ramas comparten jurisdicción, acción/pretensión y proceso) y es autónomo (código, doctrina y autores propios). Ejemplos de especialidad: prisión preventiva en lo penal; principio protectorio en lo laboral.}
\cqa{¿Cuáles son las fuentes?}{Constitución Nacional (art.~18; art.~75 incs.~22-24 para los tratados), leyes nacionales y provinciales, decretos incluidos los DNU (control: Ley 26.122), jurisprudencia y doctrina.}
\cqa{¿Qué significa que la norma procesal tiene una dinámica especial?}{Que ante la misma situación puede aplicarse de modos distintos; por ejemplo, ante la notificación de la demanda el demandado puede contestar, comparecer sin contestar o no comparecer, con consecuencias diversas.}
\cqa{¿Cómo se distinguen los sistemas anglosajón y continental, y a cuál pertenece la Argentina?}{En el anglosajón pesan el jurado y la declaración del testigo; en el continental, al que pertenece la Argentina, la prueba documental y el expediente escrito. El CPCCN de 1967 recibió influencia italiana.}
\cornellsummary{El derecho procesal civil hace actuar la voluntad de la ley cuando el conflicto torna inoperante el derecho de fondo; el juez, tercero imparcial, dicta una norma individual. Es una materia autónoma y única, con especialidades. Sus fuentes son la Constitución, los tratados, las leyes, los decretos, la jurisprudencia y la doctrina, y su historia va del proceso romano y las ordalías germánicas al derecho común y al Código de 1967.}
\end{cornell}

\chapter{El pedestal básico: jurisdicción, acción y proceso}
\label{cap:pedestal}

Tres nociones están presentes en todas las ramas del derecho procesal y conforman su \textbf{pedestal básico}: la jurisdicción, la acción (concretada en la pretensión) y el proceso. Este capítulo las desarrolla en ese orden y concluye situando la función jurisdiccional dentro de los tres poderes del Estado.

\section{Jurisdicción}

\begin{definicion}{Jurisdicción}
Término que proviene del latín \textit{iuris dictio}: «decir el derecho». Es la potestad que tiene el Poder Judicial de resolver los conflictos mediante una norma individual (sentencia).
\end{definicion}

\subsection{Jurisdicción y competencia}

La jurisdicción no debe confundirse con la competencia. La \textbf{jurisdicción} es una función constitucional del Poder Judicial. La \textbf{competencia} es la organización del trabajo de los tribunales, establecida por el legislador (no por el constituyente); delimita el ámbito de actuación de la jurisdicción y puede determinarse por territorio, materia, turno y grado. Por ejemplo, la competencia originaria y exclusiva de la CSJN surge del art.~117 CN y su competencia apelada, del art.~116 CN. El capítulo~\ref{cap:competencia} desarrolla estos criterios.

\subsection{Los momentos de la jurisdicción}

La doctrina distingue cinco momentos dentro de la jurisdicción:

\begin{enumerate}
\item \textbf{\textit{Notio}:} facultad de conocer el problema; se ejerce cuando el juez toma conocimiento de la demanda y de los hechos que la originan.
\item \textbf{\textit{Vocatio}:} facultad de convocar a las partes al proceso, citándolas para que comparezcan.
\item \textbf{\textit{Coertio}:} poder de coerción, que permite al juez aplicar sanciones para asegurar el desarrollo del proceso. Por ejemplo, si un testigo es convocado y no comparece, el juez puede ordenar su comparecencia por la fuerza pública.
\item \textbf{\textit{Iudicium}:} momento máximo, la sentencia: la norma individual que dice el derecho en el caso concreto, haciendo operativa la voluntad de la ley sustancial.
\item \textbf{\textit{Executio}:} potestad de ejecutar lo decidido. Si el condenado no cumple voluntariamente la sentencia, el juez puede usar la fuerza para hacerla efectiva. Por ejemplo, si se condena a escriturar un inmueble y el demandado no lo hace, el juez puede suscribir la escritura en su nombre.
\end{enumerate}

\subsection{Imperio y jurisdicción}

Es frecuente la confusión que identifica la jurisdicción con el imperio.

\begin{definicion}{Imperio}
Potestad de usar la fuerza pública para hacer cumplir las decisiones.
\end{definicion}

El imperio \textbf{no es exclusivo} del Poder Judicial: también lo tienen el Poder Ejecutivo (que puede detener personas) y el Poder Legislativo (que puede sancionar a sus miembros). En cambio, la potestad de «decir el derecho» sí es privativa del Poder Judicial, aunque con matices (procedimientos administrativos con revisión judicial).

\subsection{Manifestaciones de la jurisdicción fuera del Poder Judicial}

\begin{table}[H]
\centering
\renewcommand{\arraystretch}{1.25}
\begin{tabularx}{\textwidth}{>{\RaggedRight\arraybackslash}p{0.24\textwidth}>{\RaggedRight\arraybackslash}X}
\toprule
\textbf{Ámbito} & \textbf{Alcance y límites} \\
\midrule
Jurisdicción administrativa & La AFIP puede intimar, tramitar procedimientos y aplicar sanciones. Esa actuación es válida solo si el afectado puede acceder al control del Poder Judicial; la CSJN lo estableció en \textit{Fernández Arias c/ Poggio} (Fallos 247:646), que exige que todo procedimiento administrativo habilite la revisión judicial. \\
\addlinespace
Jurisdicción eclesiástica & La Sacra Rota Romana puede tramitar la nulidad de un matrimonio canónico: una jurisdicción acotada al ámbito eclesiástico. \\
\addlinespace
Entes reguladores & Los entes creados con la privatización de empresas públicas ejercen una jurisdicción limitada («jurisdicción primaria» según la doctrina americana). No pueden dictar sentencias de condena con indemnización de daños; la CSJN lo estableció en el caso \textit{Ángel Estrada y Cía.} \\
\addlinespace
Jurisdicción arbitral & Los árbitros son jueces privados elegidos voluntariamente por las partes que desean sustraerse de los jueces del Estado. Tienen jurisdicción pero \textbf{no tienen imperio}: para ejecutar sus laudos necesitan recurrir a la jurisdicción oficial. \\
\bottomrule
\end{tabularx}
\end{table}

% TODO: verificar consistencia entre fuentes originales: el resumen Cornell de la clase indicaba que la jurisdicción administrativa tampoco tiene imperio, pero el desarrollo no lo explicita.

\section{Acción y pretensión}

\begin{definicion}{Acción (sentido amplio)}
Derecho constitucional de petición ante cualquier autoridad (art.~14 CN). Es un concepto abstracto.
\end{definicion}

\subsection{La evolución científica del concepto}

La \textbf{polémica Windscheid--Muther (1856)} es el hito fundacional del derecho procesal como ciencia autónoma. Hasta entonces, la \textit{actio} del derecho romano confundía el derecho sustancial con el procesal: solo podía accionar quien era acreedor. A partir de ese debate se separó la acción del derecho de fondo: se puede accionar incluso cuando uno cree ---aunque no sea cierto--- que se le debe algo.

\textbf{Bülow (1868)} estableció que el proceso es una nueva relación jurídica, distinta de la que le dio origen. Esta tesis, completada por \textbf{Chiovenda (1913, Universidad de Bolonia)}, consolidó la autonomía científica del derecho procesal.

\subsection{La pretensión y el triángulo procesal}

\begin{definicion}{Pretensión}
Es la acción ejercida en concreto ante la jurisdicción: la manifestación de voluntad dirigida al juez contra una parte distinta, con el objeto de que se cumpla determinada conducta.
\end{definicion}

La pretensión genera el \textbf{triángulo procesal}: el actor pretende ante el juez, tercero imparcial, y el demandado resiste. La relación jurídica procesal es bilateral: el actor reclama ante la jurisdicción y el demandado tiene la posibilidad de contestar.

\begin{definicion}{Demanda}
Acto de petición formal a través del cual se inicia un proceso. Contiene la pretensión.
\end{definicion}

Son ejemplos de pretensiones contenidas en una demanda: cobrar una indemnización por los daños del accidente de colectivo, obtener el divorcio, cobrar la indemnización por despido, exigir la rendición de cuentas y obtener el desalojo de un inmueble.

\section{Proceso}

\begin{definicion}{Proceso}
Nueva relación jurídica (tesis de Bülow) que surge del conflicto llevado ante la jurisdicción. Se desarrolla sobre la base del Código Procesal y concluye con la sentencia.
\end{definicion}

\subsection{El proceso como sistema}

El proceso puede mirarse no solo como una relación jurídica entre tres sujetos (actor, juez y demandado), sino como un \textbf{sistema} complejo, siguiendo la Teoría General de Sistemas de von Bertalanffy. Un sistema es un conjunto de partes interrelacionadas que persiguen un objetivo común; son ejemplos análogos el cuerpo humano, un automóvil o una computadora.

\begin{table}[H]
\centering
\renewcommand{\arraystretch}{1.25}
\begin{tabularx}{\textwidth}{>{\RaggedRight\arraybackslash}p{0.22\textwidth}>{\RaggedRight\arraybackslash}X}
\toprule
\textbf{Elemento del sistema} & \textbf{En el proceso} \\
\midrule
Insumos & Actores y demandados, juez, empleados y funcionarios judiciales, tecnología, edificio, recursos económicos e insumos instrumentales (papel, sistemas informáticos). \\
\addlinespace
Procesador & Las etapas: \textbf{introductoria} (demanda y contestación), \textbf{probatoria} (ofrecimiento, admisión y producción de prueba), \textbf{conclusional} (alegatos de bien probado) y \textbf{sentencia}. \\
\addlinespace
Resultado & La sentencia, norma individual que resuelve el conflicto. \\
\addlinespace
Retroalimentación & Los \textbf{recursos procesales} (apelación, etc.), que controlan el resultado y permiten correcciones. \\
\bottomrule
\end{tabularx}
\end{table}

\begin{ejemplo}{Retroalimentación}
El juez dicta sentencia y la cámara la revoca. El juez advierte que aplicó mal el derecho o valoró incorrectamente los hechos, y en un caso futuro similar aplicará el criterio rectificado. También la parte que pierde, si tiene otro caso análogo, deberá replantear su estrategia.
\end{ejemplo}

\section{Los tres poderes del Estado y la función jurisdiccional}

El poder del Estado es uno solo; el soberano es el pueblo. Ese poder se organiza en tres departamentos con funciones diferenciadas: el \textbf{Legislativo} hace la ley, el \textbf{Ejecutivo} administra y el \textbf{Judicial} dirime los conflictos haciendo actuar la voluntad de la ley.

En la práctica las funciones se mezclan: el Ejecutivo legisla mediante DNU; el Legislativo cumple funciones administrativas (por ejemplo, el alquiler de oficinas) e incluso jurisdiccionales (enjuiciamiento de sus miembros); y el Judicial legisla cuando la CSJN dicta acordadas de aplicación obligatoria (por ejemplo, la acordada sobre notificación electrónica y formato de cédulas). Pero la función básica de cada poder es la descripta, y cuando se habla de jurisdicción en sentido estricto se alude a esa función esencial del Poder Judicial.

\begin{nota}
Las fuentes registran que estaba en debate una reforma judicial. La Asociación Argentina de Derecho Procesal emitió una declaración señalando que la reforma propuesta no era una reforma judicial propiamente dicha, sino una reasignación de competencias en el fuero federal. Una verdadera reforma requeriría cambiar el sistema procesal, que seguiría siendo el mismo aunque se modifique la distribución de causas.
\end{nota}

\begin{cornell}{jurisdicción, acción, pretensión y proceso}
\cqa{¿Qué es la jurisdicción y en qué se distingue de la competencia?}{Potestad del Poder Judicial de «decir el derecho» mediante una norma individual; es una función constitucional. La competencia es la organización del trabajo de los tribunales, fijada por el legislador (territorio, materia, turno, grado).}
\cqa{¿Cuáles son los cinco momentos de la jurisdicción?}{\textit{Notio} (conocer), \textit{vocatio} (convocar), \textit{coertio} (coerción), \textit{iudicium} (sentencia) y \textit{executio} (ejecución).}
\cqa{¿Por qué jurisdicción e imperio no son lo mismo?}{El imperio (uso de la fuerza pública) también lo tienen el Ejecutivo y el Legislativo; decir el derecho es privativo del Judicial. Los árbitros tienen jurisdicción pero no imperio.}
\cqa{¿Qué exige la CSJN de la jurisdicción administrativa?}{Que todo procedimiento administrativo habilite el control del Poder Judicial (\textit{Fernández Arias c/ Poggio}).}
\cqa{¿Qué aportaron Windscheid--Muther, Bülow y Chiovenda?}{Separaron la acción del derecho sustancial (1856), concibieron el proceso como nueva relación jurídica (1868) y consolidaron la autonomía científica de la materia (1913).}
\cqa{¿Qué es la pretensión y qué relación genera?}{La acción ejercida en concreto ante la jurisdicción; genera el triángulo procesal (actor, juez imparcial, demandado).}
\cqa{¿Cómo se describe el proceso como sistema?}{Insumos, procesador (etapas introductoria, probatoria, conclusional y sentencia), resultado (sentencia) y retroalimentación (recursos).}
\cornellsummary{Jurisdicción, acción/pretensión y proceso son el pedestal común del derecho procesal. La jurisdicción es la potestad de decir el derecho; la pretensión, la acción concretada ante el juez; el proceso, la nueva relación jurídica que concluye con la sentencia y puede analizarse como sistema con retroalimentación mediante recursos.}
\end{cornell}
\part{Organización judicial y competencia}

\chapter{La competencia judicial}
\label{cap:competencia}

Si la jurisdicción es el poder de decir el derecho, queda por responder qué juez resuelve cada conflicto. Este capítulo ubica la competencia dentro de la teoría general, expone sus caracteres esenciales, describe la doble estructura del sistema judicial argentino, enumera los criterios de distribución, explica la competencia de la Corte Suprema y presenta los mecanismos procesales para discutir la competencia.

\section{Ubicación de la competencia en la teoría general}

La competencia debe vincularse con los otros grandes temas de la materia: proceso, jurisdicción y acción.

\begin{table}[H]
\centering
\renewcommand{\arraystretch}{1.3}
\begin{tabularx}{\textwidth}{>{\RaggedRight\arraybackslash}p{0.19\textwidth}>{\RaggedRight\arraybackslash}X>{\RaggedRight\arraybackslash}p{0.22\textwidth}}
\toprule
\textbf{Concepto} & \textbf{Definición} & \textbf{Rol en el sistema} \\
\midrule
\textbf{Jurisdicción} & Mecanismo por el cual el Estado quita a los ciudadanos la posibilidad de resolver conflictos usando la fuerza (teoría de la sustitución). & El poder estatal \\
\addlinespace
\textbf{Acción} & La posibilidad de acceder al sistema jurisdiccional para resolver el conflicto. & El acceso al sistema \\
\addlinespace
\textbf{Proceso} & La regla de juego que fija los límites de intervención de la jurisdicción. & La regla de juego \\
\addlinespace
\textbf{Competencia} & La distribución material de los distintos conflictos entre los distintos jueces; la medida de la jurisdicción. & Quién resuelve qué conflicto \\
\bottomrule
\end{tabularx}
\end{table}

\begin{definicion}{Competencia}
Distribución material de los distintos conflictos entre los distintos jueces; es la \textbf{medida de la jurisdicción}.
\end{definicion}

En el lenguaje coloquial se usa equívocamente «competencia» en lugar de «jurisdicción». Todos los jueces están investidos con el mismo poder jurisdiccional: no hay mayor jurisdicción en un juez de la Corte que en uno de primera instancia, sino una distribución distinta de la tarea. A mayor concentración demográfica, mayor especialización de los jueces.

\begin{ejemplo}{La competencia como especialidad médica}
La jurisdicción es el sistema de salud pública en su conjunto; la competencia es la especialidad de cada médico (no cualquier doctor puede operar del corazón); el proceso son los protocolos y pasos de cada intervención.
\end{ejemplo}

\begin{ejemplo}{Especialización según la densidad poblacional}
En la Ciudad de Buenos Aires hay jueces sumamente especializados (de familia, civiles patrimoniales, comerciales). En el interior de la Provincia de Buenos Aires, los juzgados civiles y comerciales atienden esos temas conjuntamente. En el sur del país los tribunales tienen competencia universal: el juez conoce en todos los conflictos del ámbito territorial, porque hay poca población y menos conflictos.
\end{ejemplo}

\section{Caracteres de la competencia}

\subsection{Orden público}

La primera regla es que la competencia es una cuestión de \textbf{orden público}: surgido el conflicto, las personas no pueden llevarlo a cualquier juez, porque una norma estatal previa determina quién es el juez competente. Ello responde a la garantía constitucional del \textbf{juez natural}.

\subsection{Improrrogabilidad y su excepción}

\begin{regla}{artículo 1.º del CPCCN}
«La competencia atribuida a los tribunales nacionales es improrrogable.» No puede ser pactada por las partes en conflicto.
\end{regla}

\begin{importante}
\textbf{Excepción: prórroga en materia patrimonial y territorial.} Solo se puede prorrogar la competencia cuando se cumplen \emph{simultáneamente} dos condiciones: (1) que se trate de materia \textbf{patrimonial} (no de familia, penal, laboral ni federal) y (2) que la prórroga sea respecto de la competencia \textbf{territorial}. La prórroga puede ser \textbf{expresa} (pactada, por ejemplo, en un contrato) o \textbf{tácita} (el demandado contesta la demanda sin cuestionar la competencia territorial).
\end{importante}

\begin{ejemplo}{Contrato de alquiler con cláusula de prórroga}
Las partes pueden pactar que cualquier conflicto del contrato se lleve a un juzgado de la Ciudad de Buenos Aires, de la Provincia de Buenos Aires o de la Provincia de Santa Fe. Es válido porque es materia patrimonial y se pacta sobre competencia territorial.
\end{ejemplo}

En la prórroga tácita, el actor lleva el caso al juez que elige en función del territorio, sin pacto previo; si el demandado nada dice al contestar, prorroga tácitamente esa atribución, aunque por las reglas generales (por ejemplo, el lugar donde sucedió el hecho o el domicilio del demandado, fijadas en el CPCCN) podría llevarse ante otro magistrado.

\subsection{Indelegabilidad}

El juez al que se le asigna una competencia \textbf{no puede excusarse de entender} en el caso: no puede derivarlo a un colega de igual grado ni a la segunda instancia. Intervenir en todos los asuntos asignados a su competencia es un deber. Ello no significa que acepte sin más cualquier atribución: puede y debe controlar su competencia.

La mayor cantidad de fallos de la Corte Nacional se refieren precisamente a discusiones de competencia, lo que refleja que las normas no siempre son claras y los mismos conflictos pueden interpretarse de uno u otro modo. Además, los litigantes eligen ciertos fueros porque saben que son más proclives a acceder al reclamo, de modo que también allí se discute quién es el juez competente.

\section{Doble estructura del sistema judicial argentino}

\begin{table}[H]
\centering
\renewcommand{\arraystretch}{1.3}
\begin{tabularx}{\textwidth}{>{\RaggedRight\arraybackslash}p{0.16\textwidth}>{\RaggedRight\arraybackslash}X>{\RaggedRight\arraybackslash}X}
\toprule
 & \textbf{Jurisdicción federal} & \textbf{Jurisdicción ordinaria / local} \\
\midrule
\textbf{Alcance} & Presente en todo el territorio nacional & Organizada por cada provincia (poder no delegado) \\
\addlinespace
\textbf{Órgano máximo} & Corte Suprema de Justicia de la Nación & Superiores Tribunales de cada provincia \\
\addlinespace
\textbf{Código aplicable} & CPCCN y leyes federales & Código procesal propio de cada provincia \\
\bottomrule
\end{tabularx}
\end{table}

El origen histórico es el siguiente: cuando en 1853-1860 se conforma el Estado argentino, lo hace a partir de estados previos ---las provincias--- que delegaron ciertas facultades en el Estado federal y retuvieron las restantes. La materia procesal, a diferencia del derecho sustancial, fue entendida (no sin discusión en su momento) como función \textbf{no delegada}; por eso cada provincia tiene su propio código procesal y organiza su justicia local.

\subsection{Los tribunales nacionales de la Ciudad de Buenos Aires}

Los \textbf{tribunales nacionales} dependen funcionalmente de la esfera federal (Corte Suprema, Consejo de la Magistratura de la Nación), pero no conocen en materia federal sino en asuntos locales ordinarios (accidentes de tránsito, relaciones laborales, cuestiones penales no federales, etc.). Un juez nacional es designado tras concursar en el Consejo de la Magistratura, con acuerdo del Senado y por decreto presidencial.

\begin{nota}
\textbf{Contexto histórico.} La Ciudad de Buenos Aires (antes Capital Federal) era territorio de la Provincia de Buenos Aires, que al ingresar a la Federación por el Pacto de San José de Flores lo cedió al gobierno federal para asentar las autoridades nacionales. En 1853-1860 existían además numerosos territorios nacionales aún no provincializados (Patagonia, Tierra del Fuego, el norte). Como la Ciudad, densamente poblada, necesitaba tribunales y era territorio cedido, el gobierno federal organizó los actuales tribunales nacionales. Como dato ilustrativo, el Banco de la Provincia de Buenos Aires, preexistente a la Federación, es territorio provincial dentro de la capital federal.
\end{nota}

En 1994 la reforma constitucional creó la \textbf{Ciudad Autónoma de Buenos Aires (CABA)}, con autonomía prácticamente igual a la de una provincia, y en 1996 se sancionó su estatuto (su constitución). Lo lógico hubiera sido que los tribunales nacionales pasaran a la esfera local; sin embargo, desde entonces subsiste una disputa y solo se han transferido ciertas competencias penales. El núcleo de los juzgados nacionales (penal, laboral, civil y comercial) permanece en la esfera federal.

\section{Criterios de distribución de la competencia}

El legislador utiliza distintos criterios, que se \textbf{superponen}: en cada conflicto puede interactuar más de uno para determinar el juez competente.

\begin{enumerate}
\item \textbf{Fuero federal u ordinario (local).} Ciertos conflictos interesan al Estado federal y otros quedan en la órbita local. Una discusión sobre un registro de marca (cuestión civil sobre titularidad) se resuelve en un Juzgado Civil y Comercial Federal de Primera Instancia, mientras que un accidente de tránsito (también materia civil) se resuelve en la justicia ordinaria, en la Ciudad de Buenos Aires ante un Juzgado Civil.
\item \textbf{Materia.} El tipo de conflicto determina la especialidad del tribunal: civil y comercial (conflictos patrimoniales, contratos), laboral (relaciones de trabajo), familia (divorcio), penal no federal (delitos del Código Penal sin carácter federal) y contencioso-administrativo (conflictos con el Estado local).
\item \textbf{Territorio.} Se prefiere el juez más próximo al lugar del conflicto y a las partes. En familia, la competencia territorial se asigna al último lugar donde convivió la familia (la sede del hogar); ese tribunal entiende en los conflictos posteriores a la separación, como el régimen de comunicación con los hijos y los alimentos.
\item \textbf{Valor o monto.} Las causas de menor cuantía pueden asignarse a tribunales específicos (en muchas provincias, juzgados de pequeñas causas). En el CPCCN opera a través del monto de apelabilidad: según el art.~242, las causas de monto inferior al que determina la Corte Nacional no son apelables. El monto fue modificado desde el comienzo de 2025; al momento de la clase se indicó que era de \$300.000.
\item \textbf{Grado.} Se vincula con las instancias de revisión: primera instancia, segunda instancia (cámaras de apelaciones), superiores tribunales provinciales (en algunos estados) y Corte Suprema (por recurso extraordinario federal, no como instancia ordinaria).
\item \textbf{Turno.} Distribuye los asuntos que surgen en un momento dado, sobre todo en materia penal y de familia. Hay tribunales abiertos las 24 horas para recibir las denuncias de ese período; todas esas causas se asignan al tribunal de turno, que previene y queda con el conocimiento definitivo.
\end{enumerate}

\section{Competencia de la Corte Suprema}

\subsection{Competencia originaria}

\begin{regla}{artículo 117 de la Constitución Nacional}
La Corte interviene como tribunal de \textbf{primera y única instancia} en los conflictos que involucren diplomáticos, cónsules extranjeros o algún Estado extranjero en conflicto con un Estado local o ciudadano local, y en los conflictos entre dos provincias o entre un ciudadano de una provincia y otra provincia.
\end{regla}

El fundamento es evitar conflictos internos entre provincias o situaciones que comprometan las relaciones exteriores.

\subsection{Política restrictiva: el caso Barreto}

Desde hace más de diez años (desde 2007 en adelante) la Corte restringe su competencia, tanto originaria como apelada. Se contraponen dos modelos: uno de alta cantidad de causas (históricamente el argentino) y otro de menor cantidad y mayor relevancia institucional (el de la Corte Suprema de los Estados Unidos).

A partir del caso \textbf{Barreto} se definió la «causa civil»: se interviene en instancia originaria solo cuando la cuestión de derecho sustancial es \textbf{central}, no meramente tangencial. Cuando la discusión requiere analizar normas locales o administrativas provinciales, la causa debe iniciarse ante un \textbf{tribunal federal de primera instancia} con asiento en ese estado. Es una línea difusa de trazar.

\subsection{Competencia apelada}

\begin{table}[H]
\centering
\renewcommand{\arraystretch}{1.3}
\begin{tabularx}{\textwidth}{>{\RaggedRight\arraybackslash}p{0.24\textwidth}>{\RaggedRight\arraybackslash}X}
\toprule
\textbf{Vía} & \textbf{Descripción} \\
\midrule
\textbf{Competencia ordinaria apelada} & Procedía en asuntos donde fuera parte el Estado federal y se superara un determinado monto. Fue declarada inconstitucional en el caso \textbf{Anadón}, porque el Estado no debía tener el privilegio de llegar a una tercera (o cuarta) instancia de revisión; quedó prácticamente clausurada. \\
\addlinespace
\textbf{REF (art.~14 Ley 48)} & Recurso Extraordinario Federal; procede cuando se suscita una cuestión federal. \\
\addlinespace
\textbf{Arbitrariedad} & Creación jurisprudencial: procede cuando la sentencia inferior carece de fundamento o es irrazonablemente fundada. \\
\addlinespace
\textbf{Gravedad institucional} & Creación jurisprudencial: casos de singular relevancia para las instituciones del país. \\
\addlinespace
\textbf{\textit{Per saltum}} & Originariamente jurisprudencial, luego recibió sanción legislativa. Permite saltear instancias intermedias cuando la índole de la cuestión requiere intervención rápida; son casos muy excepcionales. \\
\bottomrule
\end{tabularx}
\end{table}

\begin{regla}{artículo 280 del CPCCN: «certiorari negativo»}
En línea con su política restrictiva, la Corte aplica el art.~280 para desestimar recursos extraordinarios sin más cuando considera que los casos no alcanzan el umbral de trascendencia requerido, lo que le permite seleccionar los casos que revisa.
\end{regla}

\section{Mecanismos procesales para discutir la competencia}

\subsection{Control de oficio}

La primera función del juez al recibir un caso es analizar si es competente; si no lo es, debe declarar su incompetencia y remitirlo al tribunal que considere competente. Hay una excepción: si la materia es patrimonial y la incompetencia es territorial (prorrogable), el juez no puede declararse incompetente de entrada y debe esperar a que el demandado cuestione la competencia; si no lo hace, habrá prórroga tácita. En los demás casos (competencia improrrogable), si el juez se considera competente corre traslado de la demanda; si no, lo declara por resolución fundada, apelable por las partes.

\subsection{Declinatoria e inhibitoria}

\begin{definicion}{Declinatoria (art.~8 CPCCN)}
Excepción procesal por la cual el demandado se presenta ante el mismo juez que entiende en la causa y le plantea su incompetencia. Se enmarca en las excepciones previstas desde los arts.~346-347 del CPCCN. Se discute con la otra parte y el juez resuelve aceptarla (enviando la causa al juez competente) o rechazarla.
\end{definicion}

\begin{definicion}{Inhibitoria}
Vía alternativa: el demandado acude al \textbf{juez que considera competente} y le pide que se declare competente y requiera al otro juez que se inhiba de seguir conociendo.
\end{definicion}

La inhibitoria es mucho menos frecuente en la práctica. Fue relevante, por ejemplo, en el conflicto de competencias entre tribunales nacionales y de CABA: los jueces nacionales rechazaban la asignación de competencia a los tribunales locales, y los abogados del gobierno de la Ciudad optaron por acudir directamente a un juez de la Ciudad, con la causa ya iniciada ante un juez nacional, para que se declarara competente y pidiera la inhibición de aquel.

\begin{importante}
La declinatoria y la inhibitoria son vías \textbf{alternativas y excluyentes}: si se usa una, no puede usarse la otra. En ambos casos el demandado debe \textbf{también contestar la demanda} en el mismo acto, porque en el CPCCN las excepciones se oponen junto con la contestación; aunque se use la inhibitoria, la presentación debe ser completa, aunque se haga ante el juez que no tiene la causa.
\end{importante}

\subsection{Contiendas positivas y negativas}

\begin{table}[H]
\centering
\renewcommand{\arraystretch}{1.3}
\begin{tabularx}{\textwidth}{>{\RaggedRight\arraybackslash}p{0.27\textwidth}>{\RaggedRight\arraybackslash}X}
\toprule
\textbf{Situación} & \textbf{Consecuencia} \\
\midrule
Acuerdo entre los jueces & Queda asignada la competencia; una parte puede apelar, pero el caso sigue. \\
\addlinespace
Contienda positiva & Ambos jueces se consideran competentes y quieren retener la causa; interviene el tribunal superior común. \\
\addlinespace
Contienda negativa & Ambos se consideran incompetentes y ninguno quiere asumir la causa; interviene el tribunal superior común. \\
\bottomrule
\end{tabularx}
\end{table}

Resuelve el \textbf{tribunal superior común} de los tribunales en contienda. Para evitar que todas las contiendas entre tribunales nacionales lleguen a la Corte, una acordada establece que las resuelve, sean positivas o negativas, la \textbf{Cámara del tribunal que previno}: si el caso se inició en un juzgado comercial, la Cámara Comercial; si en uno civil, la Cámara Civil.

Interviene inevitablemente la Corte cuando la contienda se da entre un tribunal de la CABA y un tribunal nacional: aunque ambos tienen asiento en la ciudad, uno pertenece a la justicia local y el otro a la federal, de modo que no tienen superior común salvo la Corte. Por eso la jurisprudencia de la Corte contiene muchos más fallos sobre competencia que sobre derecho sustancial.

\begin{importante}
Lo central en un curso de grado es conocer las normas fundamentales y las reglas que siempre estarán presentes: los caracteres de la competencia, los de su asignación y la organización de los tribunales. En casos atípicos habrá que estudiar a qué tribunal corresponde la causa, y hay cuestiones opinables, con vaivenes jurisprudenciales.
\end{importante}

\begin{cornell}{competencia judicial}
\cqa{¿Qué es la competencia y cómo se relaciona con la jurisdicción?}{Distribución material de los conflictos entre los jueces; es la medida de la jurisdicción. Todos los jueces tienen el mismo poder jurisdiccional; solo cambia la distribución de la tarea.}
\cqa{¿Cuáles son sus caracteres?}{Orden público (juez natural), improrrogabilidad (art.~1 CPCCN) e indelegabilidad (el juez no puede excusarse de intervenir).}
\cqa{¿Cuándo se puede prorrogar?}{Solo si la materia es patrimonial y la competencia territorial; la prórroga puede ser expresa (pacto) o tácita (el demandado contesta sin cuestionar).}
\cqa{¿Cómo es la doble estructura judicial y qué son los tribunales nacionales?}{Federal (Corte Suprema, CPCCN) y local (superiores tribunales provinciales, códigos propios). Los nacionales de la Ciudad dependen de la esfera federal pero conocen en asuntos locales, por razones históricas del territorio cedido.}
\cqa{¿Cuáles son los criterios de distribución?}{Fuero federal u ordinario, materia, territorio, valor, grado y turno; se superponen.}
\cqa{¿Cuál es la competencia de la Corte y qué política sigue?}{Originaria (art.~117 CN, delimitada por \textit{Barreto}) y apelada (ordinaria, cerrada por \textit{Anadón}; extraordinaria: REF, arbitrariedad, gravedad institucional, \textit{per saltum}). Restringe su intervención, incluso con el art.~280 CPCCN.}
\cqa{¿Cómo se discute la competencia?}{De oficio por el juez, por declinatoria (ante el juez de la causa) o por inhibitoria (ante el juez que se considera competente), vías excluyentes; en ambos casos debe contestarse la demanda. Las contiendas las resuelve el superior común; entre CABA y nacionales, la Corte.}
\cornellsummary{La competencia reparte los conflictos entre jueces que comparten un mismo poder jurisdiccional. Es de orden público, improrrogable (salvo materia patrimonial y territorial) e indelegable. Se distribuye por fuero, materia, territorio, valor, grado y turno dentro de una doble estructura federal y local. La Corte Suprema restringe su competencia y las discusiones se canalizan por control de oficio, declinatoria e inhibitoria.}
\end{cornell}
\part{Garantías, sujetos y vías de resolución de conflictos}

\chapter{El debido proceso legal}
\label{cap:debido}

El debido proceso legal es la garantía madre del ordenamiento jurídico y el criterio con el que se evalúan todos los institutos que siguen en este libro. El capítulo recorre su origen histórico, su recepción en la Argentina, su estructura en vertientes, principios y sistemas, y su proyección sobre la sentencia.

\section{Importancia y origen histórico}

\subsection{Importancia para el abogado}

La noción de debido proceso legal es un conocimiento que el abogado debe dominar con rigor: quien no sepa qué es no puede ejercer la profesión. Su aplicación excede la litigación: sirve también para preparar y redactar contratos, brindar asesoramiento preventivo para evitar litigios y trabajar con legisladores en la elaboración de leyes. Se observa, además, que una característica de la idiosincrasia argentina es el escaso respeto por la letra de la ley, conducta que debilita las instituciones.

Aun en los casos que se pierden, el abogado puede aspirar a que el cliente pierda a través de un juicio justo o de una decisión razonablemente aceptable.

\begin{nota}
Se recomienda la obra \textit{La razonabilidad de las leyes}, de Juan Francisco Linares, antiguo profesor de la Facultad, perteneciente a la escuela ecológica del derecho. % TODO: verificar consistencia entre fuentes originales: denominación de la escuela del derecho a la que pertenece el autor.
Fue premiada con el Premio Facultad y su subtítulo, según las notas, es «El debido proceso legal como garantía denominada de la Constitución Nacional». % TODO: verificar referencia presente en las fuentes originales: «denominada» / «nominada».
Su tesis es que el debido proceso es una garantía constitucional que no aparece nominada específicamente en ningún artículo de la Constitución Nacional. Escrita aproximadamente en la década del 40 del siglo pasado, sigue siendo la base para entender el control de constitucionalidad y, a partir de él, el actual control de convencionalidad.
\end{nota}

\subsection{De la Carta Magna al Estado constitucional de derecho}

\begin{table}[H]
\centering
\renewcommand{\arraystretch}{1.25}
\begin{tabularx}{\textwidth}{>{\RaggedRight\arraybackslash}p{0.2\textwidth}>{\RaggedRight\arraybackslash}X}
\toprule
\textbf{Momento} & \textbf{Aporte al debido proceso} \\
\midrule
Carta Magna (1215) & El rey Juan sin Tierra extraía más recursos mediante impuestos de nobles y trabajadores; los nobles se aliaron con el pueblo para frenar su poder y surgió la Carta Magna. Establecía que nadie puede ser privado de su libertad, su vida o su propiedad sino por un juicio de sus pares. Entonces no se llamaba \textit{due process of law} sino «ley de la tierra» (\textit{law of the land}). \\
\addlinespace
\textit{Habeas corpus} (siglo XIII) & Corporizó la noción en sus primeros tiempos: resguardar la libertad y, a través de ella, la vida de la persona. \\
\addlinespace
Fines del siglo XVIII & Por las colonias inglesas pasa a América. La Constitución americana incorpora el \textit{due process} (quinta enmienda, aprox.\ 1791) y la federaliza (decimocuarta enmienda); la Argentina la tomó como modelo, hasta decirse que la nuestra está vaciada en su molde. La Revolución francesa derriba la monarquía, instala la Primera República y cristaliza las ideas liberales. \\
\addlinespace
Mediados del siglo XX & Tras la Segunda Guerra Mundial surgen las Naciones Unidas y los primeros tratados de derechos humanos; flaquean conceptos como jurisdicción y soberanía y el Estado se orienta a proteger a los más desvalidos. El debido proceso se amplifica porque \textbf{los nuevos derechos necesitan sus garantías}. \\
\bottomrule
\end{tabularx}
\end{table}

\section{El debido proceso en la Argentina}

\subsection{Los derechos por generaciones}

\begin{table}[H]
\centering
\renewcommand{\arraystretch}{1.2}
\begin{tabularx}{\textwidth}{>{\RaggedRight\arraybackslash}p{3.6cm}>{\RaggedRight\arraybackslash}p{1.9cm}>{\RaggedRight\arraybackslash}X>{\RaggedRight\arraybackslash}p{2.2cm}}
\toprule
\textbf{Reforma} & \textbf{Año} & \textbf{Derechos} & \textbf{Tipo} \\
\midrule
Constitución original & 1853/1860 & Políticos y civiles & 1.ª generación \\
Art.~14 bis & 1957 & Laborales y sociales & 2.ª generación \\
Reforma constitucional & 1994 & Medio ambiente, consumidor, usuario, patrimonio cultural & 3.ª generación \\
\bottomrule
\end{tabularx}
\end{table}

El derecho del consumidor y el ambiental no se enseñaban antiguamente porque no existían: son concepciones modernas surgidas del desarrollo de los derechos de las convenciones internacionales.

\subsection{La reforma de 1994, el bloque de constitucionalidad y las Reglas de Brasilia}

La reforma de 1994 modificó la pirámide jerárquica del art.~31 CN: la nueva pirámide resulta de la lectura integrada de los arts.~31 y 75 incs.~22, 23 y 24, el llamado \textbf{bloque de constitucionalidad}.

\begin{importante}
\textbf{Art.~75 inc.~22 CN.} Los tratados internacionales de derechos humanos suscriptos por la Argentina adquieren jerarquía constitucional, y el Congreso puede incorporar nuevos tratados con el voto de los dos tercios de la totalidad de los miembros de cada Cámara. Así, el debido proceso \textbf{siempre puede ampliarse}.
\end{importante}

Las 100 Reglas de Brasilia no son un tratado suscripto por la Argentina, pero la CSJN les dio tanta importancia que, por la \textbf{Acordada n.º 5 de 2009}, dispuso observar sus recaudos en los procesos con personas vulnerables. La principal implicancia es que el juez debe dar primacía a la \textbf{oralidad} para asegurar el acceso a la justicia de quienes se encuentran en esa situación.

\subsection{El retorno a la democracia en 1983}

Con el advenimiento de la democracia (1983), que dejó atrás la década del 70, se restableció la vigencia de la Constitución y la CSJN comenzó a dictar líneas doctrinarias que marcaron un cambio fundamental.

\begin{importante}
\textbf{Del Estado de derecho positivista al constitucional.} Antes de 1983 primaba el positivismo clásico: «lo que dice la ley es lo que vale». Desde entonces, la ley vale a la luz de los principios de la Constitución y de los tratados de derechos humanos.
\end{importante}

\section{Estructura del debido proceso: vertientes, principios y sistemas}

El debido proceso puede describirse como un \textbf{sistema a través del cual se respetan los principios que reflejan las cartas fundamentales, los tratados internacionales y los derechos humanos}. Es tan abstracto que no puede definirse con precisión: es una noción \textbf{dinámica} que se amplía a medida que aparecen nuevos derechos que requieren nuevas garantías.

\subsection{Las dos vertientes}

\begin{table}[H]
\centering
\renewcommand{\arraystretch}{1.2}
\begin{tabularx}{\textwidth}{>{\RaggedRight\arraybackslash}X>{\RaggedRight\arraybackslash}X}
\toprule
\textbf{Vertiente adjetiva (formal)} & \textbf{Vertiente sustantiva} \\
\midrule
Se refiere al proceso formal y a los principios constitucionales que deben respetarse durante su desarrollo. & Se refiere a la razonabilidad de la decisión final: el proceso puede haberse desarrollado formalmente bien y la sentencia ser, sin embargo, arbitraria. \\
\addlinespace
\textbf{Herramienta:} el sistema creado por el legislador. & \textbf{Herramienta:} control de razonabilidad (art.~28 CN). \\
\addlinespace
\textit{Ejemplos:} bilateralidad, escritura/oralidad, preclusión, juez natural. & \textit{Ejemplo:} sentencia arbitraria por falta de fundamento. \\
\bottomrule
\end{tabularx}
\end{table}

\subsection{Los principios constitucionales}

\begin{definicion}{Principio}
Presupuesto político con contenido jurídico, fundante de un ordenamiento procesal.
\end{definicion}

Los principios de la vertiente adjetiva surgen de la Constitución y tienen \textbf{una sola cara}: su contrario sería inconstitucional.

\begin{table}[H]
\centering
\renewcommand{\arraystretch}{1.2}
\small
\begin{tabularx}{\textwidth}{>{\RaggedRight\arraybackslash}p{2.5cm}>{\RaggedRight\arraybackslash}p{2.0cm}>{\RaggedRight\arraybackslash}X>{\RaggedRight\arraybackslash}p{3.2cm}}
\toprule
\textbf{Principio} & \textbf{Fuente} & \textbf{Contenido} & \textbf{Su contrario} \\
\midrule
Igualdad & Art.~16 CN & Todas las personas son iguales ante la ley. & Desigualdad \\
Legalidad & Art.~17 CN & Nadie puede ser privado de su propiedad sin sentencia fundada en ley. & Ilegalidad \\
Defensa en juicio & Art.~18 CN & La defensa en juicio de la persona y sus derechos es inviolable. & Indefensión \\
Juez natural & Art.~18 CN & Nadie puede ser sometido a comisiones especiales. & Juez \textit{ad hoc} \\
Razonabilidad & Art.~28 CN & Los principios, derechos y garantías no pueden ser alterados por las leyes que reglamentan su ejercicio. & Irrazonabilidad \\
Tutela efectiva inmediata & CADH art.~25 & Tutela judicial efectiva, acceso rápido a la justicia. & Denegación de justicia \\
\bottomrule
\end{tabularx}
\end{table}

\begin{table}[H]
\centering
\renewcommand{\arraystretch}{1.2}
\begin{tabularx}{\textwidth}{>{\RaggedRight\arraybackslash}p{3.2cm}>{\RaggedRight\arraybackslash}X}
\toprule
\textbf{Norma} & \textbf{Texto según las notas} \\
\midrule
Art.~16 CN & Todos sus habitantes son iguales ante la ley, y admisibles en los empleos sin otra condición que la idoneidad. La igualdad es la base del impuesto y de las cargas públicas. \\
\addlinespace
Art.~17 CN & La propiedad es inviolable, y ningún habitante de la Nación puede ser privado de ella, sino en virtud de sentencia fundada en ley. \\
\addlinespace
Art.~18 CN & Ningún habitante de la Nación puede ser penado sin juicio previo fundado en ley anterior al hecho del proceso, ni juzgado por comisiones especiales, o sacado de los jueces designados por la ley antes del hecho de la causa. Nadie puede ser obligado a declarar contra sí mismo\ldots{} La defensa en juicio de la persona y de los derechos es inviolable. \\
\addlinespace
Art.~28 CN & Los principios, garantías y derechos reconocidos en los anteriores artículos no podrán ser alterados por las leyes que reglamenten su ejercicio. \\
\bottomrule
\end{tabularx}
\end{table}

\subsection{Los sistemas procesales}

\begin{definicion}{Sistema}
Forma metódica a través de la cual los principios cobran vida dentro de un ordenamiento procesal.
\end{definicion}

A diferencia de los principios, los sistemas tienen \textbf{dos facetas} entre las que el legislador elige según el tiempo y el lugar.

\begin{table}[H]
\centering
\renewcommand{\arraystretch}{1.2}
\small
\begin{tabularx}{\textwidth}{>{\RaggedRight\arraybackslash}p{3cm}>{\RaggedRight\arraybackslash}X>{\RaggedRight\arraybackslash}X}
\toprule
\textbf{Principio que protege} & \textbf{Opción A} & \textbf{Opción B} \\
\midrule
Igualdad / defensa & Escritura (CPCCN actual) & Oralidad \\
Celeridad / seguridad jurídica & Preclusión (etapas que se cierran) & Unidad de vista (sin plazos perentorios) \\
Juez natural & Tribunal unipersonal (juzgado de primera instancia) & Tribunal colegiado (Cámara de Apelaciones) \\
Juez natural & Única instancia (ej.: justicia laboral provincial) & Doble instancia \\
\bottomrule
\end{tabularx}
\end{table}

\paragraph{Bilateralidad.} Una parte accede a la jurisdicción e interpone un planteo y el juez se lo hace saber a la otra para que pueda contestarlo, de modo que ambas están en igualdad de condiciones. El sistema reporta al principio y por eso es constitucionalmente válido. Si un juez condenara a alguien sin escucharlo, el sistema sería inconstitucional por violar el principio de igualdad.

\paragraph{Preclusión.} Es la \textbf{pérdida, extinción o consumación de una facultad procesal}: se cierra la etapa y el proceso continúa hacia adelante. El sistema alternativo es el de \textbf{unidad de vista}, sin plazos, que acompaña a la parte interesada y puede retrotraerse. El CPCCN adopta el sistema preclusivo.

\paragraph{Juez natural y competencia.} La Constitución no establece cuál es genéricamente el juez natural: dispone que nadie sea sometido a comisiones especiales, por lo que debe juzgarlo el juez de la causa, que resulta del sistema de \textbf{competencia} creado por el legislador (capítulo~\ref{cap:competencia}). La jurisdicción es la potestad de decir el derecho; la competencia, la de decirlo en un determinado ámbito.

\paragraph{Doble instancia.} La Corte aún no reconoció jerarquía constitucional a la doble instancia en materia civil, comercial, laboral y otras no penales, cuestión que podría ser objeto de un futuro planteo. En materia penal, el art.~8 inc.~2 ap.~h CADH reconoce el derecho de recurrir el fallo ante juez o tribunal superior, y la CSJN, en \textit{Giroldi} (década del 90), le reconoció jerarquía constitucional.

\subsection{Principios y sistemas: diferencia conceptual}

\begin{table}[H]
\centering
\renewcommand{\arraystretch}{1.2}
\begin{tabularx}{\textwidth}{>{\RaggedRight\arraybackslash}p{2.4cm}>{\RaggedRight\arraybackslash}X>{\RaggedRight\arraybackslash}X}
\toprule
 & \textbf{Principios} & \textbf{Sistemas} \\
\midrule
Naturaleza & Presupuestos políticos con contenido jurídico, fundantes del ordenamiento. & Formas metódicas por las que los principios cobran vida. \\
Origen & La CN y los tratados internacionales. & El legislador, en el código procesal. \\
Caras & Una sola. & Dos: el legislador opta por una. \\
Validez & Su contrario es inconstitucional. & Es constitucional si respeta los principios; si viola alguno, es inconstitucional. \\
\bottomrule
\end{tabularx}
\end{table}

Como una receta, el proceso necesita ingredientes (los principios) y técnica (los sistemas) para producir el resultado justo: la sentencia.

\subsection{Clasificación de los principios}

El principio, en visión aristotélica, es el punto de partida de algo que se va a desarrollar.

\begin{table}[H]
\centering
\renewcommand{\arraystretch}{1.2}
\small
\begin{tabularx}{\textwidth}{>{\RaggedRight\arraybackslash}X>{\RaggedRight\arraybackslash}X>{\RaggedRight\arraybackslash}X}
\toprule
\textbf{Fundantes} & \textbf{Generales} & \textbf{Específicos} \\
\midrule
Surgen de la CN y los tratados; fundan el ordenamiento; una sola cara. & Hermenéuticos, de valor universal; pueden no estar positivizados; sirven como «muleta» interpretativa. \textit{Ej.:} \textit{pro actione}, \textit{non bis in idem}, buen padre de familia, \textit{favor debilis}. & Los pone el legislador sustancial en la ley; son operativos y el legislador procesal no puede apartarse de ellos. \textit{Ej.:} inmediación y oralidad en el proceso de familia (CCCN). \\
\bottomrule
\end{tabularx}
\end{table}

El art.~41 CN consagra la protección del medio ambiente (tercera generación) y el legislador lo bajó mediante la Ley General del Ambiente n.º 25.675; sin embargo, aún no se creó el sistema procesal específico para hacerlo valer, lo que ilustra la brecha entre el reconocimiento de un derecho y las herramientas procesales para protegerlo.

\begin{nota}
Según el criterio expuesto, inmediación y oralidad son en rigor sistemas, no principios, porque tienen contracara (la delegación y la escritura). Cuando el legislador sustancial los eleva a principio, el legislador procesal queda obligado a respetarlos al diseñar el sistema.
\end{nota}

\section{La sentencia, la razonabilidad y la arbitrariedad}

\begin{definicion}{Sentencia}
Norma individual que crea el juez al resolver un caso: una especie de ley particular para las partes. Respetando la razonabilidad, no puede alterar los principios, derechos y garantías de la primera parte de la Constitución.
\end{definicion}

\subsection{La subsunción jurídica}

\begin{table}[H]
\centering
\renewcommand{\arraystretch}{1.2}
\begin{tabularx}{\textwidth}{c>{\RaggedRight\arraybackslash}p{4cm}>{\RaggedRight\arraybackslash}X}
\toprule
\textbf{Paso} & \textbf{Operación} & \textbf{Explicación} \\
\midrule
1 & Fijación de hechos & El juez determina los hechos que tiene por probados con la prueba producida. \\
2 & Inserción en el plexo jurídico & Los hechos probados se «sumergen» en la norma abstracta. \\
3 & Control de constitucionalidad y convencionalidad & Se verifica que no se violen los principios de la CN ni de los tratados (arts.~1 y 2 CCCN). \\
4 & Dictado de la sentencia & La norma individual debe estar «razonablemente fundada» (art.~3 CCCN) y respetar la congruencia con lo pedido. \\
\bottomrule
\end{tabularx}
\end{table}

Por ejemplo, la ley no dice que Juan, que embistió con su auto a Pedro, debe indemnizarlo: consagra el principio \textit{alterum non laedere}. Los arts.~1 y 2 CCCN indican que la ley no es la mera letra: al aplicarla, el juez debe contemplar la Constitución y los tratados. El art.~3 exige una decisión «razonablemente fundada»; no dice que deba ser «justa».

\subsection{La sentencia arbitraria}

Puede haberse desarrollado todo el proceso con un sistema que respetó los principios y que la decisión final sea aberrante: se violó entonces el principio de razonabilidad del art.~28 CN. La Corte llamó a estos pronunciamientos \textbf{sentencias arbitrarias}; no hay definición única, sino causales identificadas por la doctrina. Hay arbitrariedad cuando el juez:

\begin{itemize}
\item deja de lado prueba conducente;
\item no aplica la ley sustantiva o se excede en su aplicación;
\item viola la congruencia (resuelve sobre lo no pedido);
\item funda la sentencia en su sola voluntad.
\end{itemize}

\begin{importante}
\textbf{Trípode de Carrió.} Causales de arbitrariedad: (1) violación al derecho de propiedad; (2) violación a la noción de debido proceso legal; (3) violación a un adecuado servicio de administración de justicia.
\end{importante}

\subsection{Sana crítica, vicios y recursos}

La sana crítica significa, en una de sus acepciones, libre de error o de vicio. Según el art.~386 CPCCN, los jueces forman su convicción de conformidad con las reglas de la sana crítica y solo deben expresar en la sentencia la valoración de las pruebas esenciales y decisivas.

El \textbf{fundamento último de los recursos es evitar el error humano}. Se distingue el vicio \textit{in iudicando} de juzgamiento (error en la fijación de los hechos) y el vicio \textit{in iudicando} de aplicación del derecho (error en la aplicación del derecho al caso). % TODO: verificar consistencia entre fuentes originales: ambos vicios figuran con la misma denominación.

\subsection{Conclusión: la analogía del gato}

Un animal con bigotes, pelaje, cola, ojos y orejas de gato lleva a concluir que es un gato. Del mismo modo, si el sistema creado por el legislador respeta los principios constitucionales y convencionales, no hay violación a ninguno dentro de él y la sentencia está suficientemente fundada, hay un debido proceso legal respetado. Esa es la regla; la excepción es la sentencia arbitraria. Cuando un sistema viola un principio, el abogado debe advertirlo y plantear la cuestión constitucional; el sistema permite corregir desde adentro (recursos) y, si no es posible, resulta una sentencia arbitraria.

\begin{cornell}{debido proceso legal}
\cqa{¿Cuál es su origen y cómo llegó a la Argentina?}{Carta Magna de 1215 («ley de la tierra»), \textit{habeas corpus}, quinta y decimocuarta enmiendas americanas (modelo de la Constitución argentina) y tratados de derechos humanos tras la Segunda Guerra Mundial.}
\cqa{¿Qué cambió con la reforma de 1994 y con 1983?}{Los tratados de derechos humanos integran el bloque de constitucionalidad (arts.~31 y 75 incs.~22-24); desde 1983 se pasa del positivismo clásico al Estado de derecho constitucional.}
\cqa{¿Cuáles son las vertientes del debido proceso?}{Adjetiva (proceso formal y principios; herramienta: el sistema) y sustantiva (razonabilidad de la decisión final; art.~28 CN).}
\cqa{¿En qué se diferencian principios y sistemas?}{Los principios surgen de la CN y tratados y tienen una sola cara; los sistemas los diseña el legislador, tienen dos facetas y son válidos si respetan los principios.}
\cqa{¿Qué es la preclusión?}{La pérdida, extinción o consumación de una facultad procesal; el CPCCN adopta este sistema frente a la unidad de vista.}
\cqa{¿Cómo opera la subsunción y cuándo hay sentencia arbitraria?}{Fijación de hechos, inserción en el plexo jurídico, control de constitucionalidad y convencionalidad y sentencia razonablemente fundada. Hay arbitrariedad si, por ejemplo, se deja de lado prueba conducente o se viola la congruencia.}
\cqa{¿Cuándo hay debido proceso según la analogía del gato?}{Cuando el sistema respeta los principios constitucionales y convencionales y la sentencia tiene fundamento suficiente.}
\cornellsummary{El debido proceso nació en el medioevo inglés y se expandió con el constitucionalismo y los derechos humanos. Se estructura en una vertiente adjetiva, que exige sistemas respetuosos de principios de una sola cara, y una sustantiva, que exige una sentencia razonablemente fundada. La tarea del abogado es detectar cuándo un sistema viola un principio o una sentencia es arbitraria.}
\end{cornell}
\chapter{Los sujetos del proceso}
\label{cap:sujetos}

Fijados el pedestal del derecho procesal, la competencia y el debido proceso, corresponde preguntarse quiénes intervienen en el proceso. Este capítulo presenta el proceso judicial como sistema, clasifica a sus intervinientes y desarrolla a los órganos de la jurisdicción, a las partes y sus colaboradores, a los sistemas y cargas que rigen su actuación, al litisconsorcio, al proceso colectivo y a los terceros.

\section{El proceso judicial como sistema de resolución de conflictos}

El proceso judicial es una de las vías para resolver un conflicto; existen otros métodos alternativos ---arbitraje, mediación y negociación--- que también son materia de estudio del derecho procesal (capítulo~\ref{cap:adr}). A lo largo de la materia se advierte que el proceso está formado por \textbf{subsistemas} que se relacionan entre sí.

\begin{definicion}{Proceso judicial}
Sistema: conjunto de actos y elementos interrelacionados que regulan la actividad de la jurisdicción y de las partes con un objetivo final, la \textbf{resolución de un conflicto}, que dicta el juez en ejercicio de su jurisdicción.
\end{definicion}

Todos los jueces ejercen la misma función jurisdiccional, pero dentro de las reglas de competencia creadas por el legislador (materia, grado y territorio). Para que la sentencia sea válida el proceso debe haber transitado el debido proceso legal en sus dos vertientes: la formal (cumplir las formas del legislador sin violar derechos constitucionales ni convencionales) y la sustancial o de razonabilidad (que la sentencia no sea arbitraria).

\subsection{Etapas del proceso ordinario}

El proceso ordinario (de conocimiento) se estructura así:

\begin{enumerate}
\item \textbf{Mediación previa obligatoria}, salvo las excepciones de la Ley de Mediación.
\item \textbf{Etapa introductoria:} se promueve la demanda y se corre traslado al demandado para que la conteste; se invocan hechos y se ofrece prueba.
\item \textbf{Audiencia preliminar} (art.~360 CPCCN), en la que el juez intenta que las partes arriben a un acuerdo conciliatorio.
\item \textbf{Etapa probatoria:} se produce la prueba necesaria para acreditar los hechos controvertidos.
\item \textbf{Etapa conclusional:} el juez dicta la sentencia (\textit{iudicium}).
\item \textbf{Etapa recursiva:} las partes pueden impugnar la sentencia, principalmente por apelación ante la Cámara.
\item \textbf{Etapa de ejecución (eventual):} si el condenado no cumple voluntariamente la obligación de dar o de hacer, se ejecuta forzadamente.
\end{enumerate}

\subsection{Clasificación de los intervinientes}

Los intervinientes son de tres tipos: (1) los \textbf{órganos de la jurisdicción} (el juez y sus auxiliares); (2) las \textbf{partes} (actor y demandado) y sus colaboradores (abogados, consultores técnicos); y (3) los \textbf{terceros} intervinientes.

\section{Los órganos de la jurisdicción}

\subsection{El juez natural}

El juez es, por excelencia, el sujeto procesal: el magistrado que, investido del poder jurisdiccional, resuelve el conflicto conforme a las reglas de competencia, organizadas respetando el principio del \textbf{juez natural}.

\begin{definicion}{Juez natural}
Aquel juez que es \textbf{imparcial} (no tiene interés personal en el resultado), \textbf{impartial} (no es parte del conflicto) e \textbf{independiente} (no está subordinado a ninguna de las partes).
\end{definicion}

\subsection{Organización judicial en la Ciudad de Buenos Aires}

Dentro de la justicia ordinaria de la Ciudad hay jueces nacionales de primera instancia en lo Civil, en lo Comercial y del Trabajo, y jueces nacionales en lo Civil y Comercial Federal (justicia federal). En segunda instancia intervienen las Cámaras (Civil y Comercial de la Nación), integradas por \textbf{salas de tres jueces} (tribunales colegiados). Por encima está la Corte Suprema, que entiende en competencia originaria (art.~117 CN) o cuando hay derecho federal en juego; el recurso extraordinario federal por sentencia arbitraria permite además el acceso a la Corte aun en supuestos de derecho común. Son normas de referencia el art.~75 inc.~12 CN (derecho común y derecho federal) y el art.~117 CN.

El juez de cada proceso se asigna \textbf{por sorteo}, tanto en primera instancia como en las salas de Cámara. La primera vez que el expediente llega a la Cámara se sortea la sala y, por el \textbf{principio de prevención}, esa misma sala interviene en las apelaciones ulteriores del expediente.

\subsection{Recusación}

La \textbf{recusación} es la facultad de las partes de excluir al juez que interviene y solicitar su reemplazo, para preservar el juez natural.

\begin{table}[H]
\centering
\renewcommand{\arraystretch}{1.25}
\begin{tabularx}{\textwidth}{>{\RaggedRight\arraybackslash}p{0.2\textwidth}>{\RaggedRight\arraybackslash}X>{\RaggedRight\arraybackslash}X}
\toprule
 & \textbf{Sin causa (art.~14 CPCCN)} & \textbf{Con causa (arts.~17 y 30 CPCCN)} \\
\midrule
Fundamento & Se excluye al juez sin expresar fundamento. & Debe invocarse y demostrarse una causal taxativa del art.~17 (parentesco, enemistad u odio, trato familiar o frecuente con una parte). \\
\addlinespace
Oportunidad & El actor en la demanda; el demandado en su primera presentación; para jueces de Cámara, el día siguiente de notificada la primera resolución de la sala. & Las mismas, y además si la causa es \textbf{sobreviniente}: dentro de los cinco días desde que se tomó conocimiento. \\
\addlinespace
Límites & Una sola vez en primera instancia y una en Cámara; no procede en juicios ejecutivos, desalojos ni sumarísimos. & El art.~30 impone al juez el deber de excusarse de oficio si advierte una causal. \\
\addlinespace
Quién resuelve & --- & La Cámara si el recusado es de primera instancia; el mismo tribunal, sin el recusado, si es miembro de Cámara o de la Corte. \\
\bottomrule
\end{tabularx}
\end{table}

% TODO: verificar consistencia entre fuentes originales: el cuadro de síntesis de la clase describía la recusación sin causa como aplicable «solo en proceso ordinario», mientras que el texto enumera procesos en que no procede.

\subsection{Deberes y facultades del juez (arts.~34 a 37 CPCCN)}

Como funcionario público, el juez tiene derechos (estabilidad, intangibilidad de las remuneraciones), pero también deberes y facultades:

\begin{itemize}
\item \textbf{Art.~34, deberes generales:} dictar resoluciones dentro de los plazos legales (providencias simples y sentencia definitiva) % TODO: verificar consistencia entre fuentes originales: no se especifican los plazos concretos.
y asistir personalmente a la audiencia preliminar del art.~360, acto indelegable por el principio de inmediación.
\item \textbf{Art.~35, potestades disciplinarias:} excluir de una audiencia a quien la perturbe y aplicar las sanciones del Código.
\item \textbf{Art.~36, facultades ordenatorias e instructorias:} convocar a las partes a una audiencia en cualquier etapa para intentar un acuerdo y dictar \textbf{medidas de mejor proveer}, destinadas a esclarecer los hechos controvertidos sin afectar el derecho de defensa ni suplir la negligencia de las partes, incluida la convocatoria de testigos y peritos para pedirles aclaraciones.
\item \textbf{Art.~37, sanciones conminatorias:} sanciones pecuniarias progresivas (\textit{astreintes}) a partes o terceros para compelerlos a cumplir sus decisiones.
\end{itemize}

\subsection{Nombramiento y remoción}

\begin{table}[H]
\centering
\renewcommand{\arraystretch}{1.25}
\begin{tabularx}{\textwidth}{>{\RaggedRight\arraybackslash}p{0.2\textwidth}>{\RaggedRight\arraybackslash}X>{\RaggedRight\arraybackslash}X}
\toprule
 & \textbf{Jueces de la Corte} & \textbf{Resto de los jueces} \\
\midrule
Nombramiento & El Presidente con acuerdo de dos tercios del Senado presentes en sesión pública (art.~99 inc.~4 CN). & También el Presidente con acuerdo de dos tercios del Senado, pero antes el Consejo de la Magistratura eleva una terna vinculante surgida de concurso público de oposición y antecedentes. \\
\addlinespace
Remoción & Juicio político: acusación ante el Senado por la Cámara de Diputados. & Jurado de Enjuiciamiento (miembros del Legislativo, del Judicial y abogados de la matrícula). \\
\bottomrule
\end{tabularx}
\end{table}

El Consejo de la Magistratura, órgano dependiente del Poder Judicial, se renueva cada cuatro años y se integra con miembros de los tres poderes, abogados de la matrícula y académicos de reconocido prestigio. Las causales de remoción son las del art.~53 CN: comisión de delito o mal desempeño.

\subsection{Auxiliares dentro del juzgado y organismos de apoyo}

El \textbf{Secretario} y el \textbf{Prosecretario} (arts.~38 y 39 CPCCN) no están investidos de la función jurisdiccional, pero dictan las resoluciones que el Código les asigna. Es relevante para los recursos: los mecanismos impugnativos difieren según quién firmó la resolución.

Otros organismos asisten a la jurisdicción: el \textbf{Cuerpo Médico Forense} (dictámenes médicos a solicitud de los magistrados), los \textbf{Cuerpos Interdisciplinarios} (asistentes sociales, psicopedagogos y psicólogos, que trabajan sobre todo con los jueces de familia) y la \textbf{Oficina de Mandamientos y Notificaciones}, cuyos oficiales notificadores llevan las cédulas a los domicilios y cuyos oficiales de justicia cumplen mandas (desocupación de inmuebles en desalojos, traba de embargos sobre maquinaria empresarial).

\subsection{El Ministerio Público}

Es un auxiliar de la jurisdicción con tres ramas:

\begin{table}[H]
\centering
\renewcommand{\arraystretch}{1.25}
\begin{tabularx}{\textwidth}{>{\RaggedRight\arraybackslash}p{0.27\textwidth}>{\RaggedRight\arraybackslash}X}
\toprule
\textbf{Rama} & \textbf{Función} \\
\midrule
Ministerio Público Fiscal & Vela por la \textbf{legalidad} del proceso. Interviene en cuestiones de orden público (competencia del tribunal, inconstitucionalidad de una norma) y emite un \textbf{dictamen}. \\
\addlinespace
Defensor de Menores e Incapaces & Interviene en todo proceso con menores o incapaces, en forma \textbf{paralela} a la representación legal de padres o curador; se le corre vista en los distintos estadios. \\
\addlinespace
Defensor de Pobres y Ausentes & Interviene cuando hay ausentes o cuando las personas carecen de recursos para una defensa particular, para garantizar su defensa. \\
\bottomrule
\end{tabularx}
\end{table}

\section{Las partes y sus colaboradores}

\subsection{Actor y demandado}

La relación es \textbf{bipolar}: las partes son siempre dos, con independencia de cuántas personas integren cada polo. El \textbf{actor} (sujeto activo) ejerce la pretensión y pone en actividad la jurisdicción; el \textbf{demandado} (sujeto pasivo) es aquel a quien se dirige la pretensión. Pueden ser personas humanas o jurídicas (de derecho público, como el Estado Nacional, o privado, como las sociedades), que actúan por sus representantes legales (presidente del directorio en una SA; socio gerente en una SRL).

Al contestar la demanda, el demandado puede \textbf{reconvenir} al actor si hay identidad de objeto (mismo hecho generador) y de sujetos. Entonces el actor es actor-reconvenido y el demandado, demandado-reconviniente: las partes siguen siendo dos, con doble calidad.

\subsection{Capacidad, legitimación y representación}

\begin{definicion}{Capacidad civil}
Aptitud para ser titular de derechos y obligaciones (CCyC) y para ejercer, durante el proceso, los deberes, cargas y facultades que impone el Código Procesal.
\end{definicion}

La regla es la capacidad; son excepciones (art.~24 CCyC) las personas por nacer, los menores de edad y los declarados incapaces por sentencia, que actúan por representantes legales cuya representación se acredita documentalmente.

\begin{definicion}{Legitimación en la causa}
Ser titular de la relación jurídica sustancial preexistente al proceso en virtud de la cual se inicia el juicio. Es \textbf{activa} (el actor es titular del derecho reclamado) o \textbf{pasiva} (el demandado es titular de la obligación reclamada).
\end{definicion}

\begin{importante}
La falta de legitimación \textbf{no invalida la calidad de parte}: las partes lo son durante todo el proceso. Pero influye en el resultado: el juez reconocerá o rechazará la pretensión en la sentencia según exista o no legitimación.
\end{importante}

\begin{ejemplo}{Falta de legitimación activa}
Quien pidió prestado un vehículo y sufre un accidente por culpa de un tercero que cruzó con semáforo en rojo promueve demanda por los daños materiales al vehículo. Es parte actora, pero quien está legitimado para reclamar es el propietario, titular de la relación sustancial.
\end{ejemplo}

\begin{ejemplo}{Falta de legitimación pasiva}
Un trabajador despedido demanda a varias empresas, pero una de ellas nunca fue su empleadora. Será parte, pero carece de legitimación pasiva porque no era titular del contrato de trabajo.
\end{ejemplo}

La \textbf{legitimación en el proceso} (representación) corresponde a quienes no pueden estar por sí en el proceso. Puede ser \textbf{legal} (el menor, por sus padres, con partida de nacimiento; el incapaz, por su curador, con testimonio de la sentencia que restringió la capacidad y lo designó) o \textbf{convencional} (poder a un abogado). Normas de referencia: art.~24 CCyC; arts.~152 y 153 CCyC.

\subsection{Poder, mandato y representación}

\begin{table}[H]
\centering
\renewcommand{\arraystretch}{1.25}
\begin{tabularx}{\textwidth}{>{\bfseries\RaggedRight\arraybackslash}p{0.2\textwidth}>{\RaggedRight\arraybackslash}X}
\toprule
Concepto & Características \\
\midrule
Poder & \textbf{Instrumento} (generalmente escritura pública) por el cual el poderdante faculta al abogado a representarlo. Solo intervienen el poderdante y el escribano. \\
\addlinespace
Mandato & \textbf{Contrato} (art.~1.319 y ss.\ CCyC) que rige la relación entre cliente y abogado; el apoderado sí está presente. Puede haber mandato sin representación. \\
\addlinespace
Representación & Puede existir \textbf{con mandato} o \textbf{sin mandato} (por ejemplo, la de un hijo menor por sus padres, sin base contractual). \\
\bottomrule
\end{tabularx}
\end{table}

El art.~56 CPCCN establece que en todas las actuaciones es \textbf{obligatoria} la asistencia letrada: los escritos llevan la firma de la parte y del abogado patrocinante, o bien la parte actúa por un abogado apoderado.

\subsection{Excepciones previas de personería y legitimación (art.~347 CPCCN)}

En la etapa introductoria pueden oponerse como excepciones previas la \textbf{falta de personería} (inc.~2: la parte carece de capacidad para estar por sí o la representación invocada no está bien acreditada) y la \textbf{falta de legitimación} (inc.~3: el actor o el demandado no es titular de la relación sustancial). La segunda puede resolverse en la etapa introductoria o \textbf{diferirse a la sentencia} si requiere prueba; las partes mantienen su calidad y la excepción solo influye en el pronunciamiento.

\subsection{Domicilio real y domicilio procesal (art.~40 CPCCN)}

En su primera presentación actor y demandado deben denunciar su \textbf{domicilio real} (residencia habitual de las personas físicas; sede social registrada de las jurídicas; arts.~73, 152 y 153 CCyC), donde se practican las notificaciones que el Código dispone (traslado de la demanda, citación a absolver posiciones), y constituir un \textbf{domicilio procesal} (legal o \textit{ad litem}) dentro del radio territorial del juzgado: la Ciudad de Buenos Aires si tramita en un juzgado civil ordinario de la CABA, el Departamento Judicial de Lomas de Zamora si tramita allí, etc.

\begin{importante}
Si no se constituye domicilio procesal, las notificaciones que debían cursarse allí quedan notificadas \textbf{por ministerio de la ley} los días martes o viernes siguientes a la resolución (art.~133 CPCCN), sin cédula.
\end{importante}

Por normas complementarias, las cédulas a los domicilios constituidos se practican de modo electrónico, en el \textbf{domicilio electrónico} vinculado a la matrícula del abogado en el Colegio; las dirigidas al domicilio real siguen siendo en papel.

\section{Sistemas procesales, cargas y deberes}

Los principios (igualdad, legalidad) tienen una sola cara; los sistemas, creados por el legislador en el Código, tienen doble cara y permiten optar entre alternativas opuestas sin violar principio alguno.

\begin{itemize}
\item \textbf{Preclusión y unicidad.} En el sistema de preclusión, adoptado por el CPCCN, el proceso se divide en etapas y \textbf{siempre avanza, nunca retrocede} (salvo planteo de nulidad); lo no ejercido en una etapa no puede ejercerse en la siguiente. En la unicidad hay un único período para ejercer todas las facultades y cargas.
\item \textbf{Dispositivo y oficiosidad.} En el proceso civil y comercial, de sistema dispositivo, el impulso está a cargo de las partes, y la inactividad puede conducir a la \textbf{caducidad de la instancia}. En el proceso laboral, de oficiosidad, el impulso lo lleva la jurisdicción, por lo que no existe la caducidad de instancia.
\end{itemize}

\begin{nota}
La caducidad de instancia no interrumpe la prescripción: si el proceso caduca, hay que reiniciarlo si la acción aún no prescribió.
\end{nota}

\begin{table}[H]
\centering
\renewcommand{\arraystretch}{1.25}
\begin{tabularx}{\textwidth}{>{\bfseries\RaggedRight\arraybackslash}p{0.17\textwidth}>{\RaggedRight\arraybackslash}X}
\toprule
Concepto & Características \\
\midrule
Carga procesal & Facultad de realizar ciertos actos en ciertos momentos; «imperativo del propio interés». Su no ejercicio no trae sanción; a lo sumo, un beneficio para el otro litigante. \\
\addlinespace
Deber procesal & Imperativo legal cuyo incumplimiento \textbf{sí trae sanción}. \\
\addlinespace
Obligación & Término del derecho privado: vínculo jurídico en que el deudor debe una prestación a su acreedor, bajo apercibimiento de ejecución. \\
\bottomrule
\end{tabularx}
\end{table}

\begin{ejemplo}{Carga}
El demandado puede no contestar la demanda en los 15 días del proceso ordinario. No hay sanción, pero el no ejercicio de esa carga beneficia al actor: el juez tendrá cierta presunción de verdad sobre los hechos invocados, que el demandado no negó ni contó a su manera.
\end{ejemplo}

\begin{ejemplo}{Deber}
El testigo notificado debe comparecer, declarar y decir la verdad; si no comparece puede ser multado y si miente incurre en falso testimonio. Del mismo modo, el juez debe dictar sentencia en el plazo legal, bajo posibles sanciones.
\end{ejemplo}

\section{Litisconsorcio y acumulación de procesos}

El litisconsorcio se presenta cuando hay \textbf{multiplicidad de actores o de demandados}.

\begin{itemize}
\item \textbf{Necesario (obligatorio):} es indispensable integrar a \textbf{todos} los litisconsortes para que la sentencia sea útil y eficaz; si la litis no se integra, el juez puede detener el proceso hasta completarla, de oficio o a pedido de parte. \textit{Ejemplo:} cuatro condóminos (25\,\% cada uno) y uno promueve la división solo contra dos, omitiendo al cuarto; para que la sentencia sea útil y oponible a todos hay que integrar a todos.
\item \textbf{Facultativo (voluntario):} varias partes pueden reunirse por \textbf{conexidad en la pretensión}, un hecho o título único; no es obligatorio y pueden promoverse procesos separados. \textit{Ejemplo:} un conductor embiste a tres vehículos, con daños y lesiones a sus ocupantes; los damnificados pueden demandar juntos o por separado.
\end{itemize}

Si optaron por procesos separados procede la \textbf{acumulación}: todos tramitan ante el mismo juez, el que \textbf{previno} (aquel en cuyo proceso se notificó primero la demanda), y se dicta una única sentencia, evitando el «escándalo jurídico» de pronunciamientos contradictorios. Si los procesos están en etapas muy avanzadas y distintas, puede denegarse para evitar demoras injustificadas.

\section{Proceso colectivo}

El proceso colectivo \textbf{no debe confundirse con el litisconsorcio}: en este cada parte tiene un derecho subjetivo individual y todas son individualizables; en aquel hay un \textbf{hecho generador común} que puede afectar a un número indeterminado de personas. No está regulado en el CPCCN ni en una ley específica: debe buscarse en la Constitución (art.~41, medio ambiente sano; art.~42, usuarios y consumidores; art.~43, 2.º párrafo, que habilita el amparo contra la discriminación y en lo relativo a ambiente, competencia, usuario, consumidor y derechos de incidencia colectiva en general, al afectado, al defensor del pueblo y a las asociaciones) y en los precedentes de la Corte.

\begin{table}[H]
\centering
\renewcommand{\arraystretch}{1.25}
\begin{tabularx}{\textwidth}{>{\bfseries\RaggedRight\arraybackslash}p{0.27\textwidth}>{\RaggedRight\arraybackslash}X}
\toprule
Tipo de derecho & Características \\
\midrule
Subjetivos individuales & Propios del proceso ordinario y del litisconsorcio; titulares identificables y determinables. \\
\addlinespace
Difusos o de incidencia colectiva & Un hecho afecta a un número \textbf{indefinido e indeterminable} de personas por la lesión de un \textbf{bien colectivo indivisible} (por ejemplo, el medio ambiente). \\
\addlinespace
Individuales homogéneos & El mismo hecho afecta a un número numeroso pero \textbf{determinable} de personas; cada una puede ser resarcida individualmente. Son «homogéneos» por la causa común. \\
\bottomrule
\end{tabularx}
\end{table}

\begin{ejemplo}{Derechos individuales homogéneos}
Proceso iniciado por los usuarios afectados por un aumento arbitrario de la tarifa de un servicio público domiciliario: cada uno tiene un perjuicio cuantificable, pero comparten el hecho generador. El art.~43 CN lo habilita.
\end{ejemplo}

\section{Terceros intervinientes}

Los \textbf{terceros propiamente dichos} intervienen por un \textbf{interés propio}; no lo son, en sentido estricto, quienes participan sin interés propio (peritos, testigos). La intervención puede ser \textbf{obligada} (arts.~94 y 96 CPCCN) o \textbf{voluntaria} (arts.~90 y 91), esta última litisconsorcial o adhesiva.

\begin{itemize}
\item \textbf{Intervención obligada.} Las partes pueden citar a un tercero en la etapa introductoria demostrando que la controversia le es \textbf{común} y que la sentencia \textbf{puede alcanzarlo}; se sustancia con traslado a la otra parte y resuelve el juez. \textit{Ejemplos:} el fiador demandado que cita al deudor principal (podría luego repetir contra él); el demandado por daños que cita en garantía a su aseguradora conforme a la Ley de Seguros, pudiendo ejecutarse la sentencia contra ella como si fuera parte, salvo defensas que excedan el marco del debate (art.~96).
\item \textbf{Voluntaria litisconsorcial.} El tercero demuestra que, desde el derecho sustancial, estaba \textbf{legitimado para ser actor o demandado}; admitido, tiene las mismas facultades que la parte original. \textit{Ejemplos:} otro socio que se suma a quien impugna la validez de una asamblea; otro copropietario que se suma al que demanda al consorcio por la reparación de una parte común.
\item \textbf{Voluntaria adhesiva.} El tercero \textbf{no tiene rol autónomo}: acredita sumariamente un \textbf{interés propio conexo}, adhiere a una parte y queda limitado por su actividad; la sentencia no lo alcanza directamente. \textit{Ejemplo:} el acreedor que se presenta en los procesos de su deudor, porque si este pierde y su patrimonio se afecta, podría perjudicarse su crédito.
\end{itemize}

\begin{cornell}{sujetos del proceso}
\cqa{¿Quiénes intervienen en el proceso y cuáles son sus etapas?}{Órganos de la jurisdicción, partes y colaboradores, y terceros. Etapas: mediación previa, introductoria, audiencia preliminar (art.~360), probatoria, conclusional, recursiva y ejecución eventual.}
\cqa{¿Qué es el juez natural?}{El juez imparcial, impartial e independiente; la recusación (sin causa, art.~14; con causa, art.~17) y la excusación (art.~30) lo preservan.}
\cqa{¿Qué deberes y facultades tiene el juez?}{Arts.~34-37 CPCCN: dictar resoluciones en plazo, asistir a la audiencia preliminar, disciplinar, conciliar, dictar medidas de mejor proveer e imponer \textit{astreintes}.}
\cqa{¿En qué se distinguen capacidad, legitimación y representación?}{Capacidad: aptitud para derechos y para actuar en el proceso. Legitimación: titularidad de la relación sustancial (no afecta la calidad de parte, sí el resultado). Representación: actuar por quien no puede estar por sí (legal o convencional).}
\cqa{¿Qué diferencian poder, mandato y representación?}{Poder: instrumento. Mandato: contrato entre cliente y abogado. Representación: puede haber con o sin mandato.}
\cqa{¿Qué consecuencia tiene no constituir domicilio procesal?}{Las notificaciones quedan hechas por ministerio de la ley los martes o viernes (art.~133 CPCCN).}
\cqa{¿Carga, deber y obligación?}{Carga: imperativo del propio interés, sin sanción. Deber: su incumplimiento tiene sanción. Obligación: vínculo de derecho privado ejecutable.}
\cqa{¿Qué distingue litisconsorcio necesario de facultativo, y proceso colectivo de litisconsorcio?}{En el necesario, la sentencia exige integrar a todos; en el facultativo hay conexidad y es opcional. El proceso colectivo parte de un hecho común que afecta a un número indeterminado o numeroso de personas.}
\cqa{¿Qué clases de intervención de terceros existen?}{Obligada (arts.~94 y 96), voluntaria litisconsorcial y voluntaria adhesiva (arts.~90 y 91).}
\cornellsummary{El proceso es un sistema de actos en que intervienen el juez y sus auxiliares, las partes con sus representantes y colaboradores, y los terceros con interés propio. El Código regula recusación, deberes judiciales, legitimación, domicilios, cargas, litisconsorcio y terceros, bajo los sistemas de preclusión y dispositivo.}
\end{cornell}
\chapter{El conflicto y los sistemas de resolución de conflictos}
\label{cap:adr}

El proceso judicial es una de las vías para resolver un conflicto, pero no la única. Este capítulo estudia primero la teoría del conflicto, que permite al operador jurídico identificar y mapear el conflicto; luego clasifica los sistemas de resolución y desarrolla la negociación, la mediación, la conciliación y el arbitraje.

\section{Teoría del conflicto}

\subsection{El abogado como operador de conflictos}

La teoría del conflicto tiene una fase \textbf{estática} (identificar cuándo se está ante un conflicto y de qué tipo) y una fase \textbf{dinámica} (las escaladas y las herramientas para frenarlas). El análisis estático no es de la esencia del conflicto, que es dinámico, pero sirve para \textbf{mapearlo} y gestionarlo luego con las herramientas adecuadas.

Todo abogado tiene como materia prima el conflicto, y todo funcionario o empleado del Poder Judicial, y claramente el juez, tiene como tarea resolverlo. Sin embargo, la teoría del conflicto no es obligatoria en las universidades y en algunos planes es optativa, de modo que se forman operadores de conflictos sin todas las herramientas. Con frecuencia se trata del mismo modo un conflicto patrimonial, de familia, con un sindicato, con una empresa o con un trabajador, aunque todos son específicos y requieren su propio mapeo. Ante el cambio que supone el proceso por audiencias, la inmediación y el juego de las partes, es necesaria una formación que hoy los estudiantes no reciben.

\subsection{Resolución y solución del conflicto}

La finalidad del derecho es la paz social, es decir, poner fin a los conflictos. La sentencia establece un ganador y un perdedor; en algunos casos pone fin al conflicto y en otros no. En un conflicto de familia, por ejemplo, lo que determine el juez puede ser solo un eslabón, porque intervienen emociones, sentimientos y la cuestión familiar misma.

\begin{importante}
Debe distinguirse la \textbf{resolución} del conflicto ---que está formalmente en la sentencia--- de la \textbf{solución} del conflicto, que puede tener que ver con otra cosa. De allí la importancia de que las partes, asistidas por operadores de conflictos, puedan darse su propia solución, que siempre será mejor que la que disponga el juez.
\end{importante}

\subsection{Concepto de conflicto}

\begin{definicion}{Conflicto según Cossio (Escuela Egológica del Derecho)}
«Una interferencia intersubjetiva de intereses»: cuando en la conducta humana existe interferencia entre dos personas en una relación intersubjetiva cuyos intereses son incompatibles, nace el conflicto.
\end{definicion}

\begin{definicion}{Conflicto según Entelman}
Hay conflicto cuando dos o más personas tienen objetivos incompatibles entre sí en una relación de interdependencia. Sus elementos son tres: dos o más personas, objetivos incompatibles y relación de interdependencia.
\end{definicion}

Remo Entelman es el creador de la Teoría del Conflicto (estática y dinámica). Se recomienda también \textit{Mapeo de Conflictos}, de Raúl Calvo Soler, con proyecciones de casos concretos.

Los abogados suelen creer que hay conflicto solo cuando está en juego una norma que prohíbe una acción; pero dos personas pueden tener objetivos incompatibles aunque ambas conductas estén permitidas. Y aunque en el común de la gente el conflicto se asocia con violencia, puede no haberla.

\subsection{Elementos del conflicto: análisis estático}

\begin{description}[leftmargin=0pt,style=nextline]
\item[Comunicación e información.] Deben estudiarse los elementos de la comunicación (emisor, receptor, mensaje, canal y código): la elección del canal (un WhatsApp, un correo electrónico, una carta documento) puede evitar o no las escaladas. El operador debe además recabar y organizar los datos, pues la información correctamente analizada genera la inteligencia que permite decidir.
\item[Los actores.] Son los protagonistas del conflicto, quienes tienen una intención y una preocupación por el resultado y pueden cambiar su dinámica; no son necesariamente el actor y el demandado del proceso. Pueden ser individuales o colectivos, con liderazgo formal o informal y con organización o sin ella (un grupo que ocupa un inmueble no es lo mismo que un sindicato). Hay que preguntarse si están en la mesa todos los que tienen que estar: en un divorcio, la nueva pareja de una de las partes puede incidir en el resultado (por ejemplo, en la división de bienes), y aunque no siempre deba sentarse a la mesa, sí puede ser necesario incorporarla a la reunión del abogado con su cliente.
\item[Conciencia y percepción.] La \textbf{conciencia} es el acto intelectual por el cual el actor admite encontrarse, respecto de otro, en una situación en que cree, o ambos creen, tener objetivos incompatibles. Hay quienes no la tienen porque ninguna norma les prohíbe algo: si los trabajadores reclaman mejoras y amenazan con parar la planta y el abogado le dice al presidente de la compañía que no está obligado porque no se firmaron paritarias, lo lleva a creer que no hay conflicto; por eso trabajar su conciencia es fundamental. Sin ella, un actor puede no legitimar al otro en la mesa («no tiene derecho a reclamarme»), y así no se puede avanzar. La \textbf{percepción} es el contenido con que llegan al intelecto las opiniones, amenazas o riesgos, y está vinculada con el conocimiento e información que se tiene del otro.
\item[Los objetivos.] Lo que quieren los protagonistas y perciben como incompatible. Pueden ser \textbf{tangibles} (verbalizables, como «mi pretensión son 100 mil pesos») u \textbf{ocultos}, por razones estratégicas, de prestigio o vergüenza, o porque, si el actor los verbaliza, la otra parte automáticamente los negará.
\item[Las relaciones de poder.] El poder son los \textbf{recursos} de cada actor; conviene enlistar los propios y los del contrario: amenaza de sanción o promesa de premio, tiempo de espera (no es lo mismo no necesitar cobrar ahora que necesitarlo urgentemente), paciencia, alternativa al acuerdo, posibilidad de hacer la primera oferta, información sobre el caso y conocimiento del otro actor (muy importante en familia).
\item[Suma cero y suma variable.] Quien cree que todo lo que gana el otro lo pierde él tiene mentalidad de \textbf{suma cero} (conflictos \textbf{puros}, con un ganador y un perdedor). Los conflictos \textbf{impuros} son de \textbf{suma variable}: hay una distribución de ganancias, y de forma cooperativa y creativa pueden alcanzarse soluciones que beneficien a todos. En la formación del operador está la clave.
\item[Terceros y coaliciones.] Hay que evaluar qué terceros pueden incidir en el resultado y si cabe una coalición o alianza con alguno.
\end{description}

\subsection{El mapeo del conflicto}

Mapear es estudiar, en cada conflicto particular: quiénes son los actores (individual o colectivo, cohesión, liderazgo, posibilidad de fragmentarse); qué objetivos tienen (tangibles u ocultos); qué grado de conciencia tienen; qué recursos de poder tienen (alternativas al acuerdo, paciencia, tiempo, conocimiento del otro, información); qué terceros pueden intervenir; si es posible una coalición; y si alguna parte puede amenazar con sanciones o prometer premios. Estos elementos inciden en la toma de decisiones en la dinámica del conflicto (escaladas y desescaladas mediante técnicas de negociación).

\section{Clasificación de los sistemas de resolución}

\begin{table}[H]
\centering
\renewcommand{\arraystretch}{1.25}
\begin{tabularx}{\textwidth}{>{\RaggedRight\arraybackslash}p{0.2\textwidth}>{\RaggedRight\arraybackslash}X>{\RaggedRight\arraybackslash}X}
\toprule
 & \textbf{Sistemas no adversariales} & \textbf{Sistemas adversariales} \\
\midrule
Matriz y formas & La \textbf{negociación}: directa (autocomposición entre partes) o asistida por un tercero (mediación, conciliación, facilitadores, tercero neutral del sistema anglosajón). & Un \textbf{tercero decide}: proceso judicial y arbitraje (de derecho, de amigables componedores y pericia arbitral). \\
\addlinespace
Intervención de tercero & No siempre (solo en la negociación asistida). & Siempre (juez o árbitro que decide). \\
\addlinespace
Mecanismo procedimental & Rige la informalidad. & Fijado por ley (proceso judicial) o por las partes (arbitraje). \\
\addlinespace
Prueba & No es necesaria. & Necesaria: el tercero evalúa la prueba. \\
\addlinespace
Quién resuelve & Las propias partes. & El juez o árbitro impone la decisión. \\
\addlinespace
Resultado & Acuerdo (con valor de sentencia si es homologado). & Sentencia o laudo. \\
\bottomrule
\end{tabularx}
\end{table}

En una línea imaginaria, la negociación directa (nula intervención de terceros) está en un extremo y el proceso judicial (máxima intervención) en el otro. El calificativo «alternativos» debe ponerse entre comillas: implicaría que el proceso judicial fue previo, y no es así, pues la mediación, la conciliación y el arbitraje existen prácticamente desde que existe la humanidad (tribus, clanes, derecho romano).

\section{Negociación: el Método de Harvard}

Fisher y Ury, de la Escuela de Derecho de la Universidad de Harvard, fundaron las técnicas de negociación en cuatro puntos:

\begin{enumerate}
\item \textbf{Separar a las personas del problema:} ser duro con el problema y suave con las personas. La persona no es el problema; muchas veces cambiar los interlocutores a través de los profesionales mejora la comunicación.
\item \textbf{Concentrarse en los intereses, no en las posiciones:} las posiciones parecen inamovibles, pero el \textbf{interés} es lo que está detrás y puede ser compatible con el de la otra parte.
\item \textbf{Generar opciones (torbellino de ideas):} ensayar variedad de acuerdos posibles, saliendo de la hipótesis de ganador y perdedor, sin que todas las ideas deban ser factibles.
\item \textbf{Buscar criterios objetivos} que fundamenten el acuerdo y convenzan de que es viable y razonable: casos similares de la jurisprudencia, uno o varios tasadores para el valor de un inmueble o campo.
\end{enumerate}

\begin{ejemplo}{El conflicto de las naranjas}
Dos personas quieren la misma naranja (posición). Al indagar el interés, uno la quiere para comerla (el jugo o la pulpa) y el otro para hacer dulce (la cáscara). Con esa diferencia se puede alcanzar un acuerdo.
\end{ejemplo}

Si las partes no se dan su propia solución, el camino es el sistema adversarial.

\section{Mediación}

La mediación está regulada por la Ley 26.589 y su decreto reglamentario 1.467/2011, que forman parte de las leyes complementarias del Código Procesal. En el ámbito nacional es una instancia \textbf{previa y obligatoria} en los asuntos civiles o comerciales, en general de contenido patrimonial, con las excepciones que fija la ley. Agotada la instancia se extiende un acta que deja constancia de las partes, del objeto y de que no se alcanzó acuerdo; ese acta de cierre es uno de los requisitos al promover la demanda y habilita la instancia judicial.

\begin{itemize}
\item \textbf{El mediador:} abogado matriculado con dos años de antigüedad y otros requisitos, habilitado por el Ministerio de Justicia (Poder Ejecutivo). Su función es acercar a las partes y mejorar su comunicación para que ellas encuentren su propia solución; no les propone directamente fórmulas conciliatorias.
\item \textbf{Confidencialidad:} lo dicho, manifestado o exhibido en la mediación no puede utilizarse en el proceso judicial, ni puede convocarse al mediador como testigo. Ello permite negociar con libertad.
\item \textbf{Mediación oficial y privada:} en la oficial se pide a la cámara del fuero el sorteo de un mediador. En la privada se elige un mediador de confianza, que convoca a la otra parte por carta documento; esta puede elegir entre cinco mediadores ofrecidos (incluido el del requirente) y, si no elige, queda el del requirente. La conciliación laboral previa no tiene esta posibilidad.
\item \textbf{Valor del acuerdo:} el acta con el acuerdo vale para las partes como sentencia y, ante incumplimiento, se ejecuta por la vía de ejecución de sentencias (art.~499 y ss.; para acuerdos de mediación, art.~500 CPCC).
\end{itemize}

\section{Conciliación}

\subsection{Conciliación intraprocesal}

El art.~360 CPCCN (modificado por la Ley 26.589) permite al juez derivar el caso a mediación o pedir su reapertura si advierte que el conflicto podría resolverse por métodos alternativos, dentro de un plazo de 30 días, y prevé como primera etapa de la audiencia preliminar que el juez invite a las partes a conciliar (capítulo~\ref{cap:audiencia}). Según el art.~36 inc.~2, el juez puede además convocar a las partes a conciliar en cualquier momento antes de la sentencia (etapa introductoria, probatoria o conclusional).

\begin{table}[H]
\centering
\renewcommand{\arraystretch}{1.25}
\begin{tabularx}{\textwidth}{>{\RaggedRight\arraybackslash}p{0.25\textwidth}>{\RaggedRight\arraybackslash}X>{\RaggedRight\arraybackslash}X}
\toprule
\textbf{Aspecto} & \textbf{Juez conciliador} & \textbf{Mediador} \\
\midrule
Proponer fórmulas & Sí, con criterios objetivos, la causa y la prueba, sin incurrir en prejuzgamiento. & No. \\
Confidencialidad & No rige. & Rige. \\
Informalidad & No rige; se cumplen las formas del proceso. & Rige. \\
Dependencia & Poder Judicial. & Ministerio de Justicia (Poder Ejecutivo). \\
\bottomrule
\end{tabularx}
\end{table}

\begin{regla}{art.~309 CPCCN}
«Los acuerdos conciliatorios celebrados por las partes ante el juez y homologados por éste tendrán autoridad de cosa juzgada.» La conciliación es un modo anormal de terminación del proceso (el normal es la sentencia); si no se cumple, se ejecuta como sentencia.
\end{regla}

\subsection{Conciliación laboral obligatoria: SECLO}

Los casos de raíz laboral deben pasar por la conciliación laboral obligatoria del SECLO (Servicio de Conciliación Laboral Obligatoria), dependiente del Ministerio de Trabajo. La Ley 24.635 lo crea y modifica la ley de procedimiento laboral 18.345; el acta de cierre es requisito para promover el juicio laboral. El conciliador laboral, a diferencia del mediador, puede proponer fórmulas y debe controlar el acuerdo para su homologación por el Ministerio de Trabajo, porque están comprometidos el orden público y la protección de los derechos laborales.

\begin{importante}
Va a mediación la materia civil y comercial (y algunas cuestiones federales); va a conciliación laboral la materia laboral. Ambas son instancias previas y obligatorias a la promoción del juicio.
\end{importante}

\section{Arbitraje}

El arbitraje es un sistema adversarial: un tercero, el árbitro, decide. Está regulado en el CPCCN desde el art.~736. Solo puede someterse a arbitraje la \textbf{materia disponible}, transable y no afectada por el orden público; no los alimentos ni el divorcio. Su origen es un contrato que incorpora una \textbf{cláusula compromisoria}, por la cual las partes se comprometen a someter sus conflictos futuros a arbitraje.

\begin{table}[H]
\centering
\renewcommand{\arraystretch}{1.25}
\begin{tabularx}{\textwidth}{>{\RaggedRight\arraybackslash}p{0.27\textwidth}>{\RaggedRight\arraybackslash}X}
\toprule
\textbf{Clase} & \textbf{Características} \\
\midrule
Arbitraje \textit{ad hoc} & Creado para un conflicto concreto; las partes designan árbitros (por ejemplo, dos que eligen a un tercero, formando un tribunal de tres). \\
\addlinespace
Arbitraje institucional & Se somete a un tribunal existente con reglamento y árbitros permanentes designados por concurso (Bolsa de Comercio de Buenos Aires, Bolsa de Cereales, Colegio de Abogados de San Isidro). \\
\addlinespace
Arbitraje de derecho & El árbitro resuelve como un juez, ante una jurisdicción privada, aplicando la norma. Suele ser más rápido y económico y es el método más utilizado a nivel internacional. La decisión es el \textbf{laudo}, equiparable a una sentencia; los árbitros no pueden ejecutarlo por sí, y se ejecuta judicialmente (art.~499 y ss.). \\
\addlinespace
Amigables componedores & Se resuelve por \textbf{equidad}, atendiendo a la norma; Barrios de Angelis habla de «dulcificación de la letra de la ley». \\
\addlinespace
Pericia arbitral (art.~773) & Remite al art.~516. Procede cuando hay dificultad en las operaciones por liquidaciones muy complejas: un experto emite un dictamen para acercar o dirimir las diferencias. El Dr.~Rojas se pregunta si puede ser otro modo de conclusión del proceso. \\
\bottomrule
\end{tabularx}
\end{table}

\begin{cornell}{conflicto y sistemas de resolución}
\cqa{¿Qué es un conflicto?}{Interferencia intersubjetiva de intereses (Cossio); objetivos incompatibles de dos o más personas en una relación de interdependencia (Entelman). Puede no haber norma prohibitiva ni violencia.}
\cqa{¿Diferencia entre resolución y solución?}{La resolución está en la sentencia; la solución puede ser otra cosa, y la que se dan las propias partes siempre es mejor.}
\cqa{¿Qué elementos se analizan al mapear un conflicto?}{Actores, objetivos (tangibles u ocultos), conciencia y percepción, recursos de poder, suma cero o variable, terceros y coaliciones.}
\cqa{¿Cómo se clasifican los sistemas?}{No adversariales (negociación directa o asistida: mediación, conciliación) y adversariales (proceso judicial y arbitraje).}
\cqa{¿Cuáles son los cuatro puntos del Método de Harvard?}{Separar personas del problema, intereses antes que posiciones, generar opciones y buscar criterios objetivos.}
\cqa{¿Cómo funciona la mediación obligatoria?}{Ley 26.589: previa y obligatoria en lo civil y comercial; mediador abogado con dos años; confidencialidad; el acta de cierre es requisito de la demanda; el acuerdo se ejecuta como sentencia.}
\cqa{¿En qué se distingue el juez conciliador del mediador?}{El juez puede proponer fórmulas, sin confidencialidad ni informalidad; el mediador no propone y rige la confidencialidad. Los acuerdos homologados ante el juez tienen autoridad de cosa juzgada (art.~309).}
\cqa{¿Qué es el SECLO?}{La conciliación laboral obligatoria previa (Ley 24.635), con conciliador que propone fórmulas y acuerdo homologado por el Ministerio de Trabajo.}
\cqa{¿Cómo se diferencian las modalidades de arbitraje?}{De derecho (aplica la norma, laudo), amigables componedores (equidad) y pericia arbitral (art.~773, liquidaciones complejas). Requieren materia disponible y cláusula compromisoria.}
\cornellsummary{Los conflictos pueden resolverse por vías no adversariales, en que las partes se dan su propia solución (negociación, mediación, conciliación), o adversariales, en que un tercero decide (juicio y arbitraje). El operador jurídico debe mapear el conflicto y conocer las instancias previas obligatorias —mediación civil y comercial y SECLO— antes de promover la demanda.}
\end{cornell}
\part{Los actos procesales}

\chapter{Concepto y clasificación de los actos procesales}
\label{cap:actos}

El proceso es una cadena de actos que se suceden desde su inicio hasta la sentencia. Este capítulo define el acto procesal, identifica quiénes lo realizan, lo distingue del hecho procesal y expone la clasificación de Palacio; los capítulos siguientes estudian su forma, lugar, tiempo (\ref{cap:forma}), los plazos (\ref{cap:plazos}), las notificaciones (\ref{cap:notificaciones}) y los escritos judiciales (\ref{cap:escritos}).

\section{Concepto y sujetos}

El concepto de acto procesal guarda con el proceso la misma relación que el de acto jurídico con el derecho civil.

\begin{definicion}{Acto procesal}
Hecho voluntario que tiene por efecto inmediato lograr la constitución, el desenvolvimiento o la extinción de un proceso judicial.
\end{definicion}

El proceso es una cadena sucesiva de actos que va del inicio a la \textbf{conclusión normal}, la sentencia, a la que el proceso generalmente tiende. Puede terminar antes por los \textbf{modos anormales}: desistimiento, conciliación, transacción, allanamiento y caducidad de instancia.

Realizan actos procesales: las \textbf{partes} (actor, demandado y, en caso de litisconsorcio, sus litisconsortes: actos de petición, impulso, inicio, desenvolvimiento o extinción); los \textbf{auxiliares de las partes} (letrados patrocinantes y procuradores apoderados, en representación de la parte o autónomamente, pero a su servicio); el \textbf{órgano judicial} (el juez, mediante resoluciones, providencias y decisiones, y sus auxiliares: secretario, prosecretario administrativo, oficiales y agentes); y los \textbf{terceros vinculados} (peritos, oficiales de justicia, oficiales notificadores, martilleros y testigos).

\subsection{Hechos y actos procesales}

\begin{definicion}{Hecho procesal}
Circunstancia que impacta directamente en el curso del proceso pero que no responde a la voluntad de las partes, sino a acontecimientos naturales o ajenos a la voluntad humana.
\end{definicion}

\begin{ejemplo}{Hechos procesales}
La muerte de una de las partes durante el proceso, o la restricción de su capacidad o una enfermedad mental que la lleva a perder la plena capacidad y a ser representada por otro.
\end{ejemplo}

Los actos procesales, en cambio, responden a hechos voluntarios actuados por la persona humana (partes, letrados, juez, secretario y demás sujetos).

\section{Clasificación de los actos procesales (Palacio)}

Se sigue la clasificación del Dr.~Lino Palacio por su claridad; otros autores proponen otras.

\subsection{Actos de iniciación}

Dan comienzo al proceso y pertenecen a la \textbf{etapa constitutiva}: las \textbf{diligencias preliminares} (previas al proceso, orientadas a obtener la información necesaria para armar la demanda), la \textbf{demanda} (plantea la pretensión) y la \textbf{contestación de demanda} (postura del demandado). Véanse los capítulos~\ref{cap:demanda} y~\ref{cap:contestacion}.

\subsection{Actos de desarrollo}

\paragraph{Actos de instrucción.} Se asocian con los \textbf{hechos}, que en el proceso civil aportan las partes. Los \textbf{actos de alegación} incorporan los hechos mediante el relato (en la demanda, al plantear un hecho nuevo o uno sobreviniente); los \textbf{actos de prueba} reconstruyen los hechos alegados aportando al juez elementos para verificarlos.

\paragraph{Actos de dirección.}

\begin{table}[H]
\centering
\renewcommand{\arraystretch}{1.25}
\begin{tabularx}{\textwidth}{>{\RaggedRight\arraybackslash}p{0.24\textwidth}>{\RaggedRight\arraybackslash}X}
\toprule
\textbf{Tipo} & \textbf{Descripción} \\
\midrule
Ordenación & Organizan el proceso: de \textbf{impulso}, de \textbf{resolución} y de \textbf{impugnación}. \\
\addlinespace
Transmisión & Ponen en conocimiento de las partes las decisiones o los pedidos formulados; por ejemplo, los traslados. \\
\addlinespace
Documentación & Incorporan al proceso elementos documentales de prueba en cualquier soporte (papel, filmación, informes de organismos). \\
\addlinespace
Cautelares & Conservan situaciones de hecho o patrimoniales para asegurar que la sentencia pueda ejecutarse. \\
\bottomrule
\end{tabularx}
\end{table}

\begin{itemize}
\item \textbf{Actos de impulso.} Hacen avanzar el proceso hacia la sentencia, pasando la cadena de un eslabón al otro, y son los que \textbf{interrumpen el plazo de caducidad de instancia}. % TODO: verificar consistencia entre fuentes originales: el desarrollo dice que los actos de impulso «interrumpen» la caducidad; el cuadro de repaso original los presentaba como «únicos».
No lo son, por ejemplo, el cambio de abogado de un litigante (relevante, pero no hace avanzar el proceso) ni la justificación de inasistencia a una audiencia por un viaje.
\item \textbf{Actos de resolución o decisión.} El tribunal responde a las peticiones de las partes, decide si hace lugar o no y da sus razones.
\item \textbf{Actos de impugnación.} Buscan modificar o reemplazar una actuación judicial o actos anteriores, por recursos o por incidente de nulidad. Los recursos son la revocatoria o reposición (ante el mismo órgano), la apelación (ante el tribunal de segunda instancia) y la aclaratoria, que no se categoriza exactamente como recurso. El \textbf{incidente de nulidad} procura que se decrete la nulidad de un acto cumplido con algún vicio, anulando sus efectos y los de los actos ulteriores que lo tomaron como base.
\item \textbf{Actos cautelares.} Evitan que, al dictarse la sentencia, no haya material para hacer realizable la condena. Una medida de no innovar mantiene una situación de hecho; si el bien que se pretende recuperar desaparece del patrimonio del demandado (lo enajenó, donó o se deshizo de él), la sentencia no podrá cumplirse.
\end{itemize}

\subsection{Actos de conclusión}

El acto prototípico es la \textbf{sentencia}. Los modos anormales son:

\begin{table}[H]
\centering
\renewcommand{\arraystretch}{1.25}
\begin{tabularx}{\textwidth}{>{\RaggedRight\arraybackslash}p{0.19\textwidth}>{\RaggedRight\arraybackslash}X>{\RaggedRight\arraybackslash}p{0.2\textwidth}}
\toprule
\textbf{Modo} & \textbf{Descripción} & \textbf{Voluntad} \\
\midrule
Desistimiento & El actor decide no seguir con el juicio y abandona voluntariamente su pretensión. & Unilateral (actor) \\
\addlinespace
Allanamiento & El demandado se somete unilateralmente a la pretensión. & Unilateral (demandado) \\
\addlinespace
Conciliación & Ambas partes llegan a un avenimiento dentro del proceso, generalmente en audiencia, con concesiones recíprocas (véase el capítulo~\ref{cap:adr}). & Bilateral, en el proceso \\
\addlinespace
Transacción & Contrato extrajudicial de concesiones recíprocas para poner fin a derechos litigiosos; se presenta al juicio. & Bilateral, fuera del proceso \\
\addlinespace
Caducidad de instancia & Pasado el plazo legal sin actos de impulso, se presume el abandono del proceso; es a la vez una sanción procesal. & Presunta (inactividad) \\
\bottomrule
\end{tabularx}
\end{table}

\begin{regla}{caducidad de instancia (art.~310 CPCCN)}
En el juicio ordinario el plazo es de \textbf{seis meses}; en los procesos de ejecución e incidentes, de \textbf{tres meses}. Sin actos de impulso en ese lapso se presume el abandono y puede decretarse la caducidad, de oficio o a petición de la contraparte.
\end{regla}

\begin{cornell}{concepto y clasificación de los actos procesales}
\cqa{¿Qué es un acto procesal y en qué se distingue del hecho procesal?}{Hecho voluntario que tiene por efecto constituir, desenvolver o extinguir el proceso. El hecho procesal es ajeno a la voluntad (por ejemplo, la muerte de una parte).}
\cqa{¿Quiénes realizan actos procesales?}{Las partes, sus auxiliares, el órgano judicial con sus auxiliares y terceros vinculados (peritos, oficiales, martilleros, testigos).}
\cqa{¿Cómo clasifica Palacio los actos?}{De iniciación, de desarrollo (instrucción: alegación y prueba; dirección: ordenación, transmisión, documentación y cautelares) y de conclusión.}
\cqa{¿Por qué son importantes los actos de impulso?}{Hacen avanzar el proceso e interrumpen el plazo de caducidad de instancia.}
\cqa{¿Cuáles son los modos anormales de terminación?}{Desistimiento, allanamiento, conciliación, transacción y caducidad de instancia.}
\cqa{¿Qué plazos de caducidad fija el art.~310?}{Seis meses en el juicio ordinario; tres en ejecución e incidentes.}
\cornellsummary{El proceso es una cadena de actos voluntarios que se clasifican en iniciación, desarrollo y conclusión. Los actos de impulso evitan la caducidad de instancia, y la sentencia, modo normal de terminar el proceso, convive con modos anormales.}
\end{cornell}
\chapter{Forma, lugar y tiempo de los actos procesales}
\label{cap:forma}

\section{La forma de los actos procesales}

La forma es la manera en que el acto se exterioriza y se vuelca en el proceso. Tiene dos dimensiones: el \textbf{modo de expresión} y el \textbf{modo de recepción e incorporación}.

\subsection{Idioma y características de los escritos}

El idioma debe ser el \textbf{idioma nacional}; los actos realizados en el extranjero o en otro idioma se incorporan con traducción de traductor público nacional. Si una persona sordomuda no puede expresarse en idioma nacional, se designa un intérprete de lengua de señas, registrado y habilitado, que pase al idioma nacional escrito lo que expresa.

Acordadas de la Corte regulan cómo deben presentarse los escritos: en \textbf{tinta negra} (sin colores ni subrayados en color) y con una \textbf{suma}, título que resume el contenido («Impugna pericia médica», «Apela sentencia», «Presenta memorial», «Opone excepciones»). Su encabezamiento debe contener el nombre de quien peticiona, el del letrado patrocinante o apoderado (con tomo y folio ante el Colegio Público de Abogados), la \textbf{carátula} del expediente (actor, demandado y objeto; por ejemplo, «Carlos O. Funes contra Empresa de Transporte La Nueva Ideal S.A. sobre Daños y Perjuicios») y el \textbf{domicilio electrónico}, que hoy reemplaza al domicilio procesal físico y donde son válidas las notificaciones fehacientes. La estructura completa del escrito se desarrolla en el capítulo~\ref{cap:escritos}.

\subsection{Notas manuscritas y copias}

\begin{regla}{art.~117 CPCCN}
Se admite la nota manuscrita en el expediente para peticiones de mero trámite; se trata como un escrito y se le coloca cargo como a una presentación formal.
\end{regla}

\begin{regla}{art.~120 CPCCN}
De todo escrito del que deba darse traslado se acompañan tantas copias como partes contrarias haya. Hoy se cumple de manera virtual: dentro de las \textbf{24 horas} de presentado el escrito en mesa de entradas, el abogado debe subir el archivo al sistema de gestión judicial (NextGen).
\end{regla}

\begin{importante}
La omisión puede hacer tener por no presentado el escrito, con gravísimo perjuicio para el cliente: se puede perder un juicio por no haber presentado las copias de un escrito central, ya que el expediente se manda a desglosar y se lo tiene por no presentado.
\end{importante}

\subsection{Modo de recepción: el cargo}

Los escritos se reciben con el \textbf{cargo}, constancia que indica el tribunal, la fecha de presentación y si vino firmada por letrado o sin su firma. Lo suscribe típicamente el prosecretario administrativo (aunque pueden hacerlo el secretario o el juez) y da \textbf{fecha cierta} al escrito, transformándolo en un \textbf{acto público} incorporado al expediente.

\subsection{Las audiencias: actas, videofilmación y publicidad}

Las audiencias se incorporan mediante \textbf{actas escritas} (resumen de lo acontecido con las firmas de los comparecientes) y \textbf{videofilmación} (archivo de video incorporado al sistema de gestión judicial). Hoy se usan en forma complementaria: el acta documenta comparecencia y firmas; el video contiene el desarrollo completo, y puede verse desde el estudio.

En principio las audiencias son \textbf{públicas}; son reservadas las de familia (por la intimidad familiar y los intereses de menores) y aquellas en que las partes lo pidan o el juez lo disponga por razones fundadas. La audiencia preliminar (art.~360) debe ser presidida por el juez bajo pena de nulidad (en ella suele tomarse también la confesional); las de prueba (testimoniales, etc.) pueden ser presenciadas por el juez o delegadas al secretario.

\begin{regla}{art.~125 inc.~3 CPCCN}
Las audiencias se fijan con antelación no menor a \textbf{tres días}, salvo urgencia que justifique un plazo menor. La parte que no concurre pierde los derechos que solo puede ejercer con su presencia.
\end{regla}

\section{El lugar de los actos procesales}

Lo habitual es que los actos se cumplan en la \textbf{sede del tribunal}. Son excepciones:

\begin{enumerate}[label=\alph*)]
\item audiencias en el domicilio del declarante, si una parte o testigo no puede trasladarse por salud, edad u otras razones justificadas;
\item reconocimiento judicial de lugares, cuando el tribunal debe trasladarse;
\item actos en extraña jurisdicción mediante oficio de la Ley 22.172;
\item notificaciones en el domicilio real por el oficial notificador;
\item actuaciones del oficial de justicia (constataciones, lanzamientos, embargos de muebles en domicilios).
\end{enumerate}

\begin{ejemplo}{Street View como herramienta judicial}
En un accidente ocurrido en 2011 en Avenida 9 de Julio y Avenida Santa Fe, concurrir hoy al lugar sería inútil porque la dinámica del tráfico cambió (en 2011 no existía el Metrobús). Street View permite acceder a imágenes del lugar de entonces y brinda una reconstrucción más fiel que la inspección ocular actual; hoy tiene amplia aceptación en los tribunales.
\end{ejemplo}

\begin{ejemplo}{Oficio Ley 22.172}
Un juez de Capital Federal necesita constatar el estado de plantaciones de arándanos en un campo de Entre Ríos. No tiene jurisdicción allí, de modo que libra un oficio al juez de Entre Ríos, quien ordena a la policía y al oficial de justicia del lugar realizar la constatación y devolver el resultado.
\end{ejemplo}

\begin{definicion}{Ley 22.172}
Convenio interjurisdiccional de asistencia recíproca entre jueces de distintas jurisdicciones, que regula la comunicación de decisiones judiciales entre jueces de distintas provincias o entre el fuero federal y el provincial.
\end{definicion}

\section{El tiempo de los actos procesales: horarios hábiles}

\begin{itemize}
\item \textbf{En la sede del tribunal:} de 7:30 a 13:30 horas (acordada de la Corte). Si una audiencia en curso se prolonga más allá, se habilita el horario sin resolución expresa mientras no demande muchas más horas; de lo contrario, el juez puede suspenderla y continuarla al día siguiente.
\item \textbf{Fuera de la sede} (notificaciones, oficial de justicia): hasta las 21:00 horas. Con habilitación horaria dispuesta por el juez, el oficial puede cumplir la diligencia fuera del horario hábil.
\end{itemize}

\begin{cornell}{forma, lugar y tiempo}
\cqa{¿Qué exigencias formales tienen los escritos?}{Idioma nacional (con traducción si corresponde), tinta negra, suma, encabezamiento con datos del peticionante, letrado, carátula y domicilio electrónico.}
\cqa{¿Qué dispone el art.~120 y qué ocurre si se omite?}{Copias para los traslados, hoy subidas al sistema dentro de 24 horas; la omisión puede hacer tener por no presentado el escrito.}
\cqa{¿Qué es el cargo?}{Constancia de recepción con tribunal, fecha y firma del letrado; da fecha cierta y convierte el escrito en acto público.}
\cqa{¿Cómo se documentan y quién preside las audiencias?}{Actas y videofilmación. La preliminar la preside el juez bajo pena de nulidad; las de prueba pueden delegarse al secretario. Se fijan con tres días de antelación (art.~125).}
\cqa{¿Cuáles son los horarios hábiles?}{7:30 a 13:30 en la sede del tribunal; hasta las 21:00 fuera de ella.}
\cornellsummary{La forma de los actos incluye idioma, características del escrito, copias, cargo y registro de audiencias. Los actos se cumplen normalmente en la sede del tribunal, con excepciones por oficio, reconocimiento o diligencia externa, y en horarios hábiles diferenciados.}
\end{cornell}

\chapter{Plazos procesales}
\label{cap:plazos}

Los plazos son el tiempo dentro del cual es preciso cumplir cada acto procesal.

\section{Clasificación}

\begin{table}[H]
\centering
\renewcommand{\arraystretch}{1.25}
\begin{tabularx}{\textwidth}{>{\RaggedRight\arraybackslash}p{0.24\textwidth}>{\RaggedRight\arraybackslash}X}
\toprule
\textbf{Criterio y tipo} & \textbf{Descripción} \\
\midrule
Legales & Los fija la ley (procesal o de fondo): 5 días para apelar; 15 días hábiles para contestar la demanda. \\
\addlinespace
Judiciales & Los fija el juez ante la ausencia de plazo legal: 15 o 20 días al perito para presentar su dictamen. \\
\addlinespace
Convencionales & Los acuerdan las partes, que pueden pactar plazos distintos de los legales. \\
\addlinespace
Perentorios, fatales o preclusivos & El Código adopta el sistema de preclusión: vencido el plazo, caduca la potestad de ejercer el acto, sin que nadie deba pedirlo. Son sinónimos en la práctica forense. \\
\addlinespace
Prorrogables e improrrogables & Algunos pueden prorrogarse a petición de parte antes de su vencimiento (un plazo judicial de intimación de 10 días, revisable por el juez); otros son improrrogables (los 5 días para apelar). \\
\addlinespace
Individuales y comunes & La generalidad son individuales y corren para cada parte desde su propia notificación; los comunes corren para todas desde la última notificación (plazo para alegar y, antes de la reforma, plazo de prueba). % TODO: verificar consistencia entre fuentes originales: no se especifica a qué reforma se refiere la mención.
\\
\bottomrule
\end{tabularx}
\end{table}

\begin{ejemplo}{Preclusión}
Si la parte tiene 5 días para apelar desde la notificación y pasan 6 sin hacerlo, precluyó la posibilidad. Lo mismo ocurre con los 15 días hábiles para contestar la demanda.
\end{ejemplo}

\section{Ampliación por distancia (art.~158 CPCCN)}

\begin{regla}{art.~158 CPCCN}
Si el notificado vive lejos de la sede del tribunal, los plazos legales se amplían \textbf{un día por cada 200 km}. Si la fracción restante supera los 100 km se agrega un día; si es inferior, no.
\end{regla}

\begin{ejemplo}{Cálculo}
A 470 km: 400 km equivalen a 2 días y la fracción de 70 km no suma otro; total, 2 días. A 520 km: 2 días más 1 por la fracción de 120 km; total, 3 días.
\end{ejemplo}

\begin{cornell}{plazos}
\cqa{¿Cómo se clasifican los plazos por su origen?}{Legales, judiciales y convencionales.}
\cqa{¿Qué significa que un plazo sea perentorio?}{Que, vencido, precluye la facultad, sin necesidad de petición de la contraria.}
\cqa{¿Cómo se calcula la ampliación por distancia?}{Un día por cada 200 km, más un día si la fracción supera 100 km (art.~158).}
\cqa{¿Qué distingue los plazos individuales de los comunes?}{Los individuales corren desde la notificación de cada parte; los comunes, desde la última notificación.}
\cornellsummary{Los plazos pueden ser legales, judiciales o convencionales, y en el sistema preclusivo vencido el plazo se pierde la facultad. Se amplían por distancia y son individuales salvo los comunes, como el de alegar.}
\end{cornell}

\chapter{Comunicaciones y notificaciones procesales}
\label{cap:notificaciones}

\section{Traslados, vistas y oficios}

\begin{definicion}{Traslado}
Cuando una parte formula una petición, el juez la pone en conocimiento de las otras para que ejerzan su derecho de defensa. El traslado de la demanda es el traslado por antonomasia.
\end{definicion}

\begin{definicion}{Vista}
Nombre que hoy se da a los traslados dados a funcionarios del Ministerio Público (fiscal, defensor de menores e incapaces, representante del fisco). Entre partes, todo son traslados.
\end{definicion}

Los \textbf{oficios} son comunicaciones libradas a otros jueces, funcionarios y reparticiones públicas o privadas.

\begin{table}[H]
\centering
\renewcommand{\arraystretch}{1.25}
\begin{tabularx}{\textwidth}{>{\RaggedRight\arraybackslash}p{0.34\textwidth}>{\RaggedRight\arraybackslash}X>{\RaggedRight\arraybackslash}p{0.17\textwidth}}
\toprule
\textbf{Destinatario} & \textbf{Quién firma} & \textbf{Norma} \\
\midrule
Presidente, ministros, secretarios de Estado; otros jueces de la misma circunscripción & El juez & Art.~38 CPCCN \\
Reparticiones públicas o privadas (bancos, empresas, organismos) & El secretario, transcribiendo la providencia del juez & Art.~38 CPCCN \\
Jueces de otras provincias & Secretario o juez, con formato especial & Ley 22.172 \\
Jueces de otros países & Vía diplomática (Relaciones Exteriores) & Exhorto diplomático \\
\bottomrule
\end{tabularx}
\end{table}

% TODO: verificar consistencia entre fuentes originales: la unidad sobre sujetos procesales cita los arts. 38 y 39 CPCCN como deberes del secretario y prosecretario, mientras que esta unidad cita el art. 38 para los oficios.

\begin{ejemplo}{Exhorto diplomático}
Para obtener los movimientos de una cuenta bancaria suiza de un investigado, o para que declare un testigo que vive en Francia o Perú, se libra un exhorto con colaboración del Ministerio de Relaciones Exteriores, que lo hace llegar al juez extranjero; este cita al testigo y devuelve el resultado.
\end{ejemplo}

\section{Notificación por ministerio de la ley}

\begin{regla}{art.~133 CPCCN}
Todas las providencias se notifican por ministerio de la ley los días \textbf{martes y viernes}. El solo transcurso del martes o viernes siguiente al dictado de la resolución notifica a todas las partes y sus auxiliares.
\end{regla}

Se basa en la presunción de que esos días los expedientes están a disposición de las partes. Si el día es feriado, la notificación se cumple el siguiente martes o viernes. Si el expediente no estaba disponible (estaba «a despacho», en poder de un empleado para proveer), el abogado deja una \textbf{nota de asistencia} (antes en un libro físico, hoy virtual) y ese día no corre la notificación.

\section{Notificación personal o por cédula (art.~135 CPCCN)}

Los actos más trascendentes se notifican personalmente (la parte concurre al expediente y deja nota) o por cédula (en papel o electrónica, según haya o no domicilio electrónico constituido), o por acta notarial, telegrama colacionado o carta documento (estos últimos para la mayoría de los actos, salvo el traslado de demanda y la notificación de sentencia).

\begin{regla}{art.~135 CPCCN}
Se notifican por cédula o personalmente: el traslado de la demanda y la reconvención con sus documentos; el traslado de las excepciones previas; la citación a absolver posiciones; la declaración de cuestión de puro derecho; la citación a audiencia preliminar (art.~360); las intimaciones o apercibimientos no establecidos directamente por la ley; la devolución del expediente de cámara; el traslado de liquidaciones; la primera providencia tras volver el expediente del archivo o de estar más de tres meses fuera de secretaría; la citación a personas extrañas al proceso (testigos); las sentencias definitivas e interlocutorias con fuerza de tales; y las que ponen fin al proceso (por ejemplo, la caducidad de instancia).
\end{regla}

La cédula sigue siendo en \textbf{papel} para el traslado de la demanda (el demandado aún no constituyó domicilio electrónico), la declaración de rebeldía, la sentencia definitiva al rebelde o a quien no compareció y las notificaciones a ajenos al proceso sin domicilio electrónico.

\subsection{Procedimiento del oficial notificador en el traslado de demanda}

El oficial va al domicilio real denunciado por el actor y pregunta si la persona vive allí. Si le dicen que sí, deja aviso de que volverá al día siguiente en cierto horario; al volver, si la persona está la notifica y, si no, deja la cédula por haberle informado que vive allí. Si no vive allí, no puede dejar la cédula: informa lo manifestado y la devuelve al juzgado, que espera instrucciones del actor. Existe también la notificación \textbf{bajo responsabilidad del actor}: el oficial notifica aunque le digan que la persona no vive allí, modalidad frecuentemente atacable por nulidad (véase el capítulo~\ref{cap:notificacion}).

\section{Notificación por edictos}

Se utiliza para \textbf{personas inciertas} (por ejemplo, en las sucesiones, posibles herederos o acreedores desconocidos del causante) y para \textbf{personas de paradero desconocido}, que no pueden localizarse en ningún domicilio registrado (RENAPER, padrón electoral, AFIP, etc.). Se publican en periódicos de amplia tirada y en el Boletín Oficial; no siempre aseguran que la persona se entere, pero cumplen formalmente con la notificación.

\begin{cornell}{comunicaciones y notificaciones}
\cqa{¿Qué son traslado y vista?}{Traslado: poner en conocimiento de la contraria una petición. Vista: el traslado dado al Ministerio Público.}
\cqa{¿Quién firma los oficios?}{El juez (autoridades y jueces de la circunscripción); el secretario (reparticiones); Ley 22.172 para otras provincias; exhorto diplomático para otros países.}
\cqa{¿Cómo opera la notificación por ministerio de la ley?}{Los martes y viernes siguientes a la resolución (art.~133); no corre si el expediente no estaba disponible y se dejó nota.}
\cqa{¿Qué actos requieren cédula o notificación personal?}{Los enumerados en el art.~135: traslado de la demanda, excepciones, citación a posiciones, audiencia preliminar, sentencias, entre otros.}
\cqa{¿Cuándo procede la notificación por edictos?}{Para personas inciertas o de paradero desconocido.}
\cornellsummary{La comunicación procesal se vale de traslados, vistas, oficios y notificaciones. La regla es la notificación por ministerio de la ley los martes y viernes; los actos trascendentes se notifican por cédula o personalmente, y por edictos a personas inciertas o de paradero desconocido.}
\end{cornell}
\chapter{El escrito judicial}
\label{cap:escritos}

Con orientación práctica, este capítulo estudia el escrito judicial, herramienta cotidiana del abogado litigante. Se distingue de la demanda, que es objeto de desarrollo específico (capítulo~\ref{cap:demanda}). La estructura gráfica del escrito se encuentra en un recurso disponible en la página de la cátedra.

\begin{definicion}{Escrito judicial}
Instrumento formal que el abogado utiliza habitualmente para comunicarse con el tribunal. Mediante él se \textbf{peticiona} algo para que el tribunal lo \textbf{despache}, cumpliendo las formalidades requeridas.
\end{definicion}

\section{Los actos procesales según su finalidad}

Los actos procesales pueden clasificarse también en actos de \textbf{postulación}, de \textbf{comunicación} y de \textbf{obtención} o petición. Las clasificaciones no son verdades en sí mismas, sino herramientas más o menos útiles según la finalidad didáctica.

\begin{table}[H]
\centering
\renewcommand{\arraystretch}{1.25}
\begin{tabularx}{\textwidth}{>{\RaggedRight\arraybackslash}p{0.25\textwidth}>{\RaggedRight\arraybackslash}X}
\toprule
\textbf{Categoría} & \textbf{Contenido} \\
\midrule
Postulación & Ejercicio del \textit{ius postulandi}; son los \textbf{actos principales}, que propenden a la \textbf{sentencia de mérito}: demanda (inicia el proceso), contestación de demanda, alegato (escrito final que evalúa la prueba producida) y recurso (impugnación ante instancia superior). \\
\addlinespace
Comunicación & Se realiza por escrito, presentado físicamente y digitalizado para que la otra parte lo conozca. Ante la petición, el juez da traslado a la otra parte, manteniendo el equilibrio entre ambas. Incluye notificaciones, cédulas, oficios y exhortos. \\
\addlinespace
Obtención o petición & Pedidos del abogado a la jurisdicción para desarrollar el proceso, que el juez puede o no admitir; por ejemplo, el libramiento de un oficio, que es un pedido de informes a una entidad. \\
\bottomrule
\end{tabularx}
\end{table}

\section{Partes del escrito}

El escrito se compone de: suma, destinatario, encabezamiento (datos del presentante, domicilio, identificación de los autos y fórmula de cierre), objeto, petitorio, cierre con firmas y cargo.

\subsection{Suma o sumario}

Es el \textbf{título del escrito}, que sintetiza la petición. Se coloca en el ángulo superior; la posición (izquierda, centro, derecha) importa menos que ser \textbf{consistente}, lo que permite luego ubicar con rapidez los propios escritos en el expediente. Debe expresar lo que se pide: «líbrese oficio», «agréguese partida», «solicítese aclaración», «denúnciese hecho nuevo».

\subsection{Destinatario}

A quién se dirige: en la mayoría de los casos el juez de primera instancia, pero también la Cámara o la Corte Suprema.

\subsection{Encabezamiento}

Contiene los datos de quien se presenta y varía según la actuación:

\begin{table}[H]
\centering
\renewcommand{\arraystretch}{1.25}
\begin{tabularx}{\textwidth}{>{\RaggedRight\arraybackslash}X>{\RaggedRight\arraybackslash}X}
\toprule
\textbf{Patrocinio letrado} & \textbf{Representación con poder} \\
\midrule
El cliente actúa por su propio derecho y el abogado lo patrocina. Firman ambos: el cliente a la derecha y el letrado a la izquierda. & El abogado actúa en nombre y representación del cliente, invocando el poder otorgado. Firma solo el abogado. \\
\bottomrule
\end{tabularx}
\end{table}

El art.~56 CPCCN establece el patrocinio letrado obligatorio: el cliente puede actuar por su propio derecho y el abogado dirige estratégicamente el proceso como letrado patrocinante sin representarlo formalmente.

\begin{ejemplo}{Encabezamiento en primera persona con patrocinio letrado}
\textit{Juan Alberto Pérez, por su propio derecho, con domicilio constituido junto a mi letrado patrocinante, Dr. Jorge Armando Rojas, C.P.A.C.F. Tomo \blanco[0.8cm] Folio \blanco[0.8cm], Responsable Inscripto en IVA N.º \blanco[1.5cm], en la calle Montevideo 514, 2.º piso, oficina B, C1058 de la Ciudad de Buenos Aires, y domicilio electrónico N.º (CUIT) \blanco[2.5cm],}
\end{ejemplo}

El escrito puede redactarse en primera o tercera persona, pero el estilo elegido debe mantenerse. En cada escrito el abogado repite todos sus datos de registración y matriculación (C.P.A.C.F., con tomo y folio). % TODO: verificar consistencia entre fuentes originales: el Cornell original menciona la categoría fiscal «monotributista o responsable inscripto» y el «poder notarial»; el desarrollo solo ejemplifica «Responsable Inscripto en IVA» y alude al «poder otorgado».

\subsection{Domicilio procesal y electrónico}

El \textbf{domicilio procesal} (\textit{ad litem}, o legal) es el constituido para el proceso, en general el del abogado, donde se reciben las notificaciones; hay una superposición en que el cliente constituye domicilio en el estudio de su abogado, y \textbf{no puede haber dos domicilios constituidos separados}. Debe estar dentro del \textbf{radio territorial del tribunal}: si el pleito tramita en la Capital Federal no puede constituirse domicilio en el estudio de un colega de San Isidro, aunque el abogado viva en Martínez o La Lucila; si tramita en San Isidro, Morón, San Martín o Quilmes, debe constituirse allí.

Con la digitalización se incorporó el \textbf{domicilio electrónico}, identificado por el CUIT o CUIL del abogado; tras la Ley 26.685, que crea el expediente electrónico, la Corte desarrolló las notificaciones electrónicas. Se consigna a continuación del domicilio constituido. Los distintos tipos de domicilio (real, legal, especial, social y procesal) se retoman al estudiar la demanda.

\subsection{Identificación de los autos}

«Autos» designa al \textbf{proceso} en los usos forenses (la resolución que comienza «autos y vistos» significa «visto el expediente, resuelvo»). Se caratulan con el nombre del actor y del demandado y el objeto, y la barra \textbf{c/} reemplaza a «contra». Por ejemplo: «Pérez, Juan Alberto c/ García, Pedro s/ Ordinario --- Daños y Perjuicios», o «García, Ana María c/ Juan Alberto Pérez s/ Alimentos» (iniciado por la madre contra el padre que no los paga). «Fojas» (fs.) equivale a «hoja» y se mantiene por los usos forenses.

\subsection{Fórmula de cierre del encabezamiento, objeto y petitorio}

El encabezamiento cierra con una fórmula: «\ldots en los autos caratulados \ldots, a Vuestra Señoría digo» (en tercera persona, «dice»). El \textbf{objeto} es el cuerpo principal, donde se formula el pedido; por ejemplo, solicitar que se libre un oficio retenido a fojas tantas o no respondido por el ente requerido. Si la petición es compleja (como denunciar un hecho nuevo con documentación adjunta) puede haber un \textbf{petitorio}: primero se señala el objeto y luego se enumeran los pedidos (Primero: que se agregue la documentación; Segundo: que se tenga presente el hecho nuevo y se corra traslado).

\subsection{Cierre, firmas y cargo}

Son fórmulas habituales de cierre: «Proveer de conformidad será justicia», «Solicito se provea, así la justicia» y «Será justicia que se provea». Luego van las firmas: la parte a la derecha y el letrado patrocinante a la izquierda. El \textbf{cargo} es la constancia del tribunal con fecha, hora, si se presentó con o sin copias y la firma del letrado.

\section{Tramitación del escrito}

Presentado el escrito, mesa de entrada lo incorpora y un empleado proyecta la resolución que firmará el juez: el \textbf{auto de despacho}. Por ejemplo: «Buenos Aires, 10 de marzo de 2009. Atendiendo al objeto de la petición, líbrese el oficio solicitado.» Lo firma posiblemente el secretario, por ser un auto de mero trámite (el secretario si el oficio es a entidades privadas; el juez si es a otros tribunales). Así, a través de actos de petición, se va \textbf{impulsando el proceso}.

Al iniciar el proceso se sortea la demanda en la Cámara de Apelaciones por vía de computación; el sistema entrega una hoja que encabeza la documentación y se convierte en la \textbf{carátula del proceso} (por ejemplo, «Pérez, Juan Alberto c/ García, Pedro s/ Ordinario --- Daños y Perjuicios --- Juzgado N.º 25»). Con ella se lleva la demanda al juzgado sorteado, donde se arma el expediente físico, una carpeta plástica con la hoja inicial en la tapa.

\begin{importante}
En el expediente la \textbf{documentación va adelante} y el \textbf{escrito que se deja va atrás}, porque el tribunal lo retiene para trabajarlo y despacharlo.
\end{importante}

\begin{ejemplo}{Escrito integrador}
Caso \textit{Pérez, Juan Alberto c/ García, Pedro s/ Ordinario --- Daños y Perjuicios}. \textbf{Suma:} LÍBRESE OFICIO. \textbf{Destinatario:} Señor Juez. \textbf{Encabezamiento:} Juan Alberto Pérez, por su propio derecho, con domicilio procesal constituido junto a su letrado patrocinante, Dr.~Jorge A.~Rojas, CPACF Tomo \blanco[0.6cm] Folio \blanco[0.6cm], \ldots{} y domicilio electrónico N.º (CUIT) \blanco[1.5cm], en los autos caratulados «Pérez, Juan Alberto c/ García, Pedro s/ Ordinario», a Vuestra Señoría digo. \textbf{Objeto:} vengo a solicitar se libre el oficio dirigido a XX Sociedad Anónima, toda vez que ha vencido el plazo para su contestación y no fue respondido. \textbf{Cierre:} proveer de conformidad será justicia. \textbf{Firmas:} Juan Alberto Pérez (derecha) y Dr.~Jorge A.~Rojas (izquierda). \textbf{Cargo:} fecha, hora, copias y firma del letrado. \textbf{Auto de despacho:} «Buenos Aires, 10 de marzo de 2009. Atendiendo al objeto de la petición, líbrese el oficio solicitado. Fernández, Secretario.»
\end{ejemplo}

\begin{cornell}{el escrito judicial}
\cqa{¿Qué es un escrito judicial y cómo se clasifican los actos según su finalidad?}{Instrumento formal de comunicación con el tribunal para peticionar. Actos de postulación (demanda, contestación, alegato, recurso), de comunicación y de obtención o petición.}
\cqa{¿Cuáles son las partes del escrito?}{Suma, destinatario, encabezamiento (datos, domicilio, autos, fórmula de cierre), objeto, petitorio, cierre y firmas, cargo.}
\cqa{¿Qué diferencia hay entre patrocinio y poder?}{Con patrocinio firman cliente (derecha) y abogado (izquierda); con poder firma solo el abogado en representación.}
\cqa{¿Cómo debe ser el domicilio procesal?}{Único, dentro del radio territorial del tribunal, con domicilio electrónico (CUIT o CUIL).}
\cqa{¿Qué son los autos y la carátula?}{Autos es el proceso; la carátula nombra actor, demandado (c/) y objeto.}
\cqa{¿Qué sucede tras presentarse el escrito?}{Se coloca el cargo, se incorpora al expediente y se dicta el auto de despacho.}
\cornellsummary{El escrito judicial es el instrumento cotidiano de petición, con partes formales constantes que permiten identificarlo y despacharlo. Su encabezamiento varía según haya patrocinio o poder, y su tramitación culmina en un auto de despacho que impulsa el proceso.}
\end{cornell}
\part{Etapa introductoria del proceso de conocimiento}

\chapter{Diligencias preliminares}
\label{cap:diligencias}

Antes de la demanda, el Código permite ciertas actuaciones que preparan el proceso o resguardan la prueba. Están reguladas en el Libro Segundo (procesos de conocimiento), Título Primero, Capítulo Segundo, \textbf{arts.~323 a 329} CPCCN, y comprenden dos institutos de distinta naturaleza: las \textbf{medidas preparatorias} y la \textbf{prueba anticipada}.

\section{Medidas preparatorias}

\begin{definicion}{Medidas preparatorias}
Diligencias que permiten a quien va a ser parte (actor o demandado) requerir judicialmente, si no puede lograrlo extrajudicialmente, que se cumplan determinados actos para preparar adecuadamente la demanda o dotar de mayor eficacia al proceso.
\end{definicion}

Asisten a quien va a demandar, o a quien con fundamento prevé que será demandado, para obtener datos o elementos sin los cuales no podría iniciarse válidamente la demanda; el Código prevé consecuencias si el requerido no cumple la orden judicial.

\begin{importante}
Las medidas preparatorias \textbf{no constituyen un proceso}: son meras diligencias que se agotan con su dictado.
\end{importante}

Están previstas para los procesos de conocimiento (ordinario y sumarísimo, derogado el sumario) y, según la doctrina, también para los especiales que tramitan por esas formas. En el sucesorio, el art.~323 permite a quien se crea heredero o legatario pedir la exhibición de un testamento.

\subsection{Caracteres}

\begin{enumerate}[label=\alph*)]
\item \textbf{Excepcionales y taxativas:} solo las de los once incisos del art.~323, aunque la jurisprudencia morigeró la taxatividad y admitió otras cuando, sin ellas, se perjudicaría la eficacia de la demanda o se imposibilitaría interponerla.
\item \textbf{No abren la instancia:} la instancia se abre con la demanda.
\item \textbf{Interrumpen la prescripción:} cualquier petición judicial lo hace (art.~2546 CCyC).
\item \textbf{Juez competente:} el que corresponda al proceso que habrá de promoverse.
\item \textbf{Oportunidad:} antes de la demanda (lo habitual) o después, pero antes de la notificación del traslado, cuando queda trabada la litis.
\end{enumerate}

\subsection{Contenido (art.~323)}

\begin{table}[H]
\centering
\renewcommand{\arraystretch}{1.25}
\begin{tabularx}{\textwidth}{>{\RaggedRight\arraybackslash}p{0.14\textwidth}>{\RaggedRight\arraybackslash}X>{\RaggedRight\arraybackslash}X}
\toprule
\textbf{Inciso} & \textbf{Medida} & \textbf{Finalidad} \\
\midrule
1.º & Declaración jurada sobre datos de personalidad del futuro demandado & Determinar capacidad, domicilio y legitimación procesal \\
2.º & Exhibición de cosas y documentos & Conocer anticipadamente el estado de cosas para encauzar la demanda \\
3.º & Designación de tutor o curador & Cuando el futuro demandado es menor o incapaz sin representante \\
9, 10 y 11 & Reconocimiento de la obligación de rendir cuentas, reconocimiento de mercaderías y mensura & No caducan a los 30 días \\
\bottomrule
\end{tabularx}
\end{table}

% TODO: verificar consistencia entre fuentes originales: la tabla del apunte indica «Arts. 9, 10 y 11»; el texto se refiere a los incisos 9, 10 y 11 del art. 323.

Los datos pedidos en el inciso 1.º \textbf{no deben apuntar al fundamento de la pretensión}: no es el momento de acreditar hechos que se discutirán luego. Por ejemplo, para promover una reivindicación de cosa mueble cuyo estado se ignora (si existe, si está dañada o destruida), la medida permite conocerlo antes de demandar.

\subsection{Trámite, sanciones y caducidad}

Se pide por escrito con las formalidades de los arts.~115 y ss.\ CPCCN y 45 a 47 del Reglamento de la Justicia Nacional, explicando su vinculación con el futuro proceso. Se dicta \textbf{sin sustanciación}.

\begin{definicion}{Sustanciar}
Transitar una cuestión por los carriles adecuados para que quede en estado de ser resuelta, lo que puede o no dar lugar a controversia. «Sin sustanciación» no significa sin trámite, sino \textbf{sin controvertir con la otra parte}.
\end{definicion}

Si se deniega, la resolución es apelable; si se concede, aunque con alcance distinto del pedido, no.

El art.~329 prevé, según la medida, \textbf{multa}, \textbf{presunción en contra} del omiso, \textbf{secuestro} con eventual allanamiento (si no se exhiben cosas o documentos) y \textbf{astreintes}. Si no se contesta afirmativamente la declaración jurada sobre personalidad, los hechos se tienen por ciertos, sin perjuicio de la prueba en contrario en el proceso. Si una persona se ausentará del país, se le requiere constituir domicilio especial en cinco días (arts.~40 y 41); de no hacerlo, las notificaciones se cursan por ministerio de la ley, los martes o viernes siguientes, y desde allí corren los plazos de impugnación.

Cumplidas, las medidas \textbf{caducan a los 30 días} de la resolución que las tiene por cumplidas si no se demanda, salvo las de los incisos 9, 10 y 11.

\section{Prueba anticipada}

Está en el \textbf{art.~326}, en cuatro incisos (el último incorporado por la Ley 25.488) y una oración final sobre la prueba confesional. Probar es demostrar en el proceso la ocurrencia de los hechos. La anticipación se refiere al \textbf{momento procesal}: en los escritos constitutivos las partes \emph{ofrecen} la prueba, pero se \emph{produce} en la etapa probatoria, abierta la causa a prueba en la audiencia preliminar. Se puede anticipar cuando hay riesgo de que la prueba se pierda o se torne ineficaz por el tiempo del proceso.

Supuestos: (1) declaración de testigo de avanzada edad o próximo a ausentarse del país; (2) reconocimiento judicial de lugares o cosas que pueden modificarse; (3) dictamen pericial sobre un estado de cosas que pueda alterarse; (4) obtención, resguardo o exhibición de documentos, incluidas historias clínicas e información complementaria. Otro ejemplo: una carta documento cuya autenticidad se necesita acreditar antes de perder los elementos para verificarla. % TODO: verificar consistencia entre fuentes originales: la redacción del apunte sobre el correo y la carta documento no es clara.

Tiene \textbf{naturaleza cautelar}: «cautela» alude a obrar con previsión frente a un perjuicio, aquí la pérdida de la prueba por el paso del tiempo; comparte con las cautelares la actuación ante el peligro en la demora.

\begin{importante}
A diferencia de las medidas preparatorias, en la prueba anticipada \textbf{debe citarse a la contraparte} para que la controle; de lo contrario la prueba es \textbf{nula} (bilateralidad y audiencia). Si no conviene citarla porque se desnaturalizaría la medida, debe citarse al \textbf{defensor oficial}. Por ejemplo, en un reconocimiento judicial por ruidos molestos (el vecino que toca la batería hasta las 4 de la mañana) no se avisa al vecino, y la prueba se produce con el defensor oficial.
\end{importante}

Denegada, es apelable, excepción al art.~379 (que declara inapelable lo relativo a la denegación, sustanciación o producción de prueba); y \textbf{no caduca} si no se demanda.

\begin{cornell}{diligencias preliminares}
\cqa{¿Qué instrumentos comprenden los arts.~323 a 329?}{Las medidas preparatorias y la prueba anticipada, de distinta naturaleza.}
\cqa{¿Cuáles son los caracteres de las preparatorias?}{Excepcionales y taxativas (art.~323), no abren la instancia, interrumpen la prescripción, juez del futuro proceso, antes de trabada la litis; se dictan sin sustanciación.}
\cqa{¿Qué sanciones prevé el art.~329 y cuándo caducan?}{Multa, presunción en contra, secuestro y astreintes; caducan a los 30 días salvo los incisos 9, 10 y 11.}
\cqa{¿Qué es la prueba anticipada y qué la distingue?}{Producción de prueba antes de la etapa probatoria por riesgo de pérdida (art.~326); naturaleza cautelar; exige citar a la contraparte o al defensor oficial; es apelable si se deniega y no caduca.}
\cornellsummary{Las diligencias preliminares preparan la demanda (medidas preparatorias, que caducan y no controvierten) o resguardan la prueba en peligro (prueba anticipada, bilateral y sin caducidad).}
\end{cornell}

\chapter{La demanda}
\label{cap:demanda}

\section{Acción, pretensión y demanda}

La \textbf{acción} es la posibilidad de peticionar a las autoridades (art.~14 CN) que, en el proceso, es el poder de hacer valer una pretensión ante la jurisdicción. Suscitado el conflicto, debe llevarse a un tribunal, que tiene la potestad de decir el derecho y hacer actuar la voluntad de la ley.

\begin{table}[H]
\centering
\renewcommand{\arraystretch}{1.25}
\begin{tabularx}{\textwidth}{>{\RaggedRight\arraybackslash}p{0.17\textwidth}>{\RaggedRight\arraybackslash}X>{\RaggedRight\arraybackslash}p{0.24\textwidth}}
\toprule
\textbf{Concepto} & \textbf{Definición} & \textbf{Carácter} \\
\midrule
Acción & Derecho de peticionar que todos tienen en potencia; poder de hacer valer una pretensión ante la jurisdicción. & Potencia (abstracto) \\
Pretensión & La acción en sentido concreto: lo que se pide al juez que otorgue en la sentencia. & Acto (concreto) \\
Demanda & Acto formal de petición que contiene la pretensión y busca iniciar el proceso. & Continente de la pretensión \\
\bottomrule
\end{tabularx}
\end{table}

\begin{definicion}{Demanda}
Acto procesal por el que se formula una petición, con formalidades, para iniciar el proceso en que se ventilará un conflicto entre dos o más personas. Contiene una pretensión.
\end{definicion}

La pretensión tiene un \textbf{elemento objetivo} (objeto inmediato y mediato), la \textbf{causa} (los hechos que la fundan) y \textbf{elementos subjetivos} (quien reclama y aquel frente a quien se reclama, con el juez como director). La demanda es un \textbf{acto de postulación} de la parte con su letrado patrocinante, pues según el art.~56 no pueden ejercerse actos postulatorios sin asistencia técnica.

Provoca la \textbf{apertura de la instancia}, ámbito de las demás peticiones. Quien la inició debe impulsar el proceso: la \textbf{caducidad de la instancia} es peligrosa, pues la demanda se tiene por no interpuesta, con graves consecuencias en costas y en la responsabilidad del abogado que la permitió.

\section{Clases de demanda}

\begin{table}[H]
\centering
\renewcommand{\arraystretch}{1.2}
\begin{tabularx}{\textwidth}{>{\RaggedRight\arraybackslash}p{0.17\textwidth}>{\RaggedRight\arraybackslash}X>{\RaggedRight\arraybackslash}X}
\toprule
\textbf{Clase} & \textbf{Descripción} & \textbf{Ejemplo} \\
\midrule
Principal & Objeto no subordinado a otra pretensión. & Daños y perjuicios por accidente de tránsito. \\
Incidental & Se relaciona con otra pretensión principal. & Incidente de redargución de falsedad de instrumento público. \\
Simple & Una sola pretensión. & Cobro de suma de dinero. \\
Compleja & Varias pretensiones articuladas (art.~87). & Daños y perjuicios, daño moral y psicológico. \\
Original / reconvencional & La del actor / la del demandado contra el actor en el mismo proceso. & --- \\
Unilateral / bilateral (art.~336) & Una persona por parte / actor y demandado presentan demanda y contestación en forma conjunta. & --- \\
Litisconsorcial & Más de una parte en la posición actora o demandada. & Varios actores o demandados. \\
\bottomrule
\end{tabularx}
\end{table}

\subsection{Acumulación objetiva de pretensiones (art.~87)}

Requiere que las pretensiones no sean contradictorias, puedan sustanciarse por los mismos trámites y sean de la competencia del mismo juez; el trámite de mayor envergadura absorbe al menor.

\begin{table}[H]
\centering
\renewcommand{\arraystretch}{1.2}
\begin{tabularx}{\textwidth}{>{\RaggedRight\arraybackslash}p{0.17\textwidth}>{\RaggedRight\arraybackslash}X>{\RaggedRight\arraybackslash}X}
\toprule
\textbf{Tipo} & \textbf{Descripción} & \textbf{Ejemplo} \\
\midrule
Acumulativa & Varias pretensiones que pueden progresar indistintamente. & Daño emergente, lucro cesante y daño moral. \\
Alternativa & Una pretensión y, si es rechazada, otra (principio de eventualidad). & Simulación y, subsidiaria, acción revocatoria. \\
Subsidiaria & Progresa solo si progresa la principal. & Escrituración y, subsidiaria, daños y perjuicios. \\
\bottomrule
\end{tabularx}
\end{table}

\section{Demanda mínima, admisible y fundada}

\begin{itemize}
\item \textbf{Mínima:} petición escrita y firmada, que activa la jurisdicción (art.~50 del Reglamento de la Justicia Nacional: ningún escrito puede ser repelido de entrada sin resolución). Si solo cuenta con ello, se rechaza in limine por el art.~337.
\item \textbf{Admisible:} demanda mínima que cumple los requisitos del art.~330. Si no, el juez, en cumplimiento del deber de sanear el proceso (art.~34), puede rechazarla in limine tras intimar a subsanar, bajo apercibimiento de tenerla por desistida.
\item \textbf{Fundada:} admisible y con pretensión válida y factible de discusión. Si el objeto es ilícito (un contrato de esclavitud o para cometer un delito), será rechazada desde el inicio por improponibilidad objetiva (art.~337).
\end{itemize}

\section{Requisitos de la demanda (art.~330)}

La demanda se redacta por escrito, pues el proceso es preponderantemente escrito; los actos constitutivos (demanda, reconvención y su contestación, excepciones previas y su contestación) conforman la etapa introductoria.

\begin{regla}{art.~330 CPCCN}
Nombre y domicilio del demandante; nombre y domicilio del demandado; la cosa demandada, designada con toda exactitud; los hechos en que se funde, explicados claramente; el derecho expuesto sucintamente; y la petición en términos claros y positivos.
\end{regla}

\begin{description}[leftmargin=0pt,style=nextline]
\item[Demandante (inc.~1).] Se indica si se presenta por su propio derecho (con capacidad procesal) o por representante. Persona humana: nombre, apellido y, convenientemente, DNI (por los homónimos), con \textbf{domicilio real} (residencia habitual, art.~73 CCyC). Persona jurídica: denominación o razón social y \textbf{domicilio legal} (art.~11 Ley 19.550), inscripto en el Registro Público. Además se constituye domicilio procesal, físico y electrónico (CUIT del abogado), arts.~40 a 42.
\item[Demandado (inc.~2).] Rigen las mismas reglas; el domicilio normalmente se conoce al contestar. Para personas jurídicas puede pedirse un informe a la Inspección General de Justicia sobre su domicilio legal.
\item[Cosa demandada (inc.~3).] El objeto del litigio, vinculado al \textbf{principio de congruencia}: en el sistema dispositivo el juez se limita a las pretensiones (art.~163 inc.~6). Se señalan el \textbf{objeto inmediato} (tipo de sentencia: constitutiva, declarativa o de condena) y el \textbf{objeto mediato} (el bien de la vida: restitución de la tenencia de un inmueble, suma por daños, cese de ruidos).
\item[Hechos (inc.~4).] Son la \textbf{causa} de la pretensión, expuestos clara y circunstanciadamente (principio de sustanciación), en orden cronológico y pensando en la norma invocada, pues deben encajar en su antecedente (art.~377: quien invoca la norma prueba el antecedente de hecho). Es el \textbf{eje central de la demanda}: sobre ellos girará la prueba. Los hechos son privativos de las partes; el saber privado del juez no cuenta.
\end{description}

\begin{definicion}{Hecho}
Acontecimiento del mundo exterior que se percibe por la modificación que genera en la fuente de prueba, las cosas o las personas.
\end{definicion}

El hecho pasado se reconstruye comenzando por \textbf{afirmarlo}, con oraciones narrativas, afirmativas y claras. Puede estructurarse en una presentación (la situación inicial, a modo de fotografía), un nudo (cómo se llegó a la relación jurídica, a modo de película) y una síntesis que concluya que los hechos encuadran en la norma rectora y que, por tanto, quien embistió debe resarcir y ser condenado.

\begin{description}[leftmargin=0pt,style=nextline]
\item[Derecho (inc.~5).] Se expone \textbf{sucintamente}: rige \textit{iura novit curia}; el juez es dueño del derecho aunque no se haya invocado bien o se haya omitido, y las partes lo son de los hechos. Cuando la cuestión es eminentemente jurídica (qué norma rige, cómo interpretarla), la carga de exponer las razones jurídicas es de la parte.
\item[Petición (inc.~6).] En términos claros y positivos, con el \textbf{monto reclamado}, salvo que no sea posible determinarlo o que ello perjudique la prescripción; entonces se deja sujeto a la prueba, evitando limitar al juez a una suma que los daños o gastos posteriores minimizarían.
\end{description}

\subsection{Prueba que se acompaña (art.~333)}

Con la demanda se acompaña la \textbf{prueba documental} que está en poder del actor y se ofrecen los demás medios (testimonial, confesional, pericial, reconocimiento judicial, informativa). Si el documento está en poder de la contraparte o de un tercero, se lo individualiza. Los documentos nuevos pueden incorporarse hasta el llamamiento de autos para sentencia. Si no se ofrece la prueba junto con los escritos constitutivos, \textbf{precluye la facultad}.

\section{Efectos de la demanda}

\begin{table}[H]
\centering
\renewcommand{\arraystretch}{1.2}
\begin{tabularx}{\textwidth}{>{\RaggedRight\arraybackslash}p{0.28\textwidth}>{\RaggedRight\arraybackslash}X}
\toprule
\textbf{Efecto} & \textbf{Descripción} \\
\midrule
\multicolumn{2}{l}{\textit{Procesales}} \\
Apertura de la instancia & Habilita los actos ante el juez y hace correr el plazo de caducidad. \\
Fijación de la competencia & Si es prorrogable y el actor eligió un juez, queda fijada para él; el demandado puede oponer incompetencia. \\
Pérdida de la recusación sin causa & Solo procede en la primera oportunidad; después, solo con causa (art.~17). \\
Confesión de los hechos afirmados & El actor no puede luego pretender que no son ciertos. \\
Fijación del objeto litigioso & El \textit{tema decidendum} queda fijado para el actor, sin perjuicio de hechos nuevos. \\
Prueba anticipada confesional & Permite ofrecer y producir como prueba anticipada la propia confesión. \\
\addlinespace
\multicolumn{2}{l}{\textit{Sustanciales}} \\
Interrupción de la prescripción & Art.~2546 CCyC: por petición ante autoridad judicial, aun incompetente, también desde la medida preparatoria; opera incluso en el plazo de gracia del art.~124 CPCCN (primeras dos horas hábiles del día siguiente al vencimiento). \\
Impedimento de caducidad de derechos & Por ejemplo, demandar al Estado Nacional dentro de los 90 días de firme la denegatoria de la reclamación administrativa. \\
Extinción de opciones del actor & Quien demandó por resolución del contrato no puede luego demandar por cumplimiento. \\
Habilitación de ciertas cautelares & Intervención judicial, remoción del administrador, etc. \\
\bottomrule
\end{tabularx}
\end{table}

% TODO: verificar consistencia entre fuentes originales: la oración del apunte sobre lo que no podía hacerse antes de la demanda está incompleta y la fila sobre cautelares no está desarrollada con claridad.

\subsection{Requisitos fiscales y mediación previa}

La demanda es una acción altamente imponible: debe pagarse la tasa de justicia, salvo beneficio de litigar sin gastos. % TODO: verificar referencia presente en las fuentes originales: el apunte menciona la «ley 3.898» para la tasa de justicia sin mayor precisión.
Como requisito de admisibilidad debe acreditarse la mediación previa obligatoria (Ley 26.589, ex Ley 24.573) con el \textbf{acta de cierre}, donde conste la comparecencia o los motivos de la incomparecencia de la contraria; todos los sujetos del proceso deben haber participado. Si no se acredita, el juez manda nuevamente a mediación. % TODO: verificar referencia presente en las fuentes originales: el apunte cita «art. 96.589» en este punto.

\begin{cornell}{la demanda}
\cqa{¿Qué relación hay entre acción, pretensión y demanda?}{Acción: potencia abstracta de peticionar; pretensión: la acción concretada; demanda: el acto formal que la contiene.}
\cqa{¿Cuándo una demanda es mínima, admisible y fundada?}{Mínima: escrita y firmada. Admisible: además cumple el art.~330. Fundada: además tiene pretensión válida (no improponible).}
\cqa{¿Cuáles son los requisitos del art.~330?}{Datos y domicilio de las partes, cosa demandada, hechos, derecho y petición.}
\cqa{¿Por qué los hechos son el eje de la demanda y el derecho se expone sucintamente?}{Los hechos son de las partes y sobre ellos gira la prueba; el juez conoce el derecho (\textit{iura novit curia}).}
\cqa{¿Qué ocurre si no se ofrece la prueba con la demanda?}{Precluye la facultad (art.~333).}
\cqa{¿Qué efectos tiene la demanda?}{Procesales (apertura de instancia, fijación de competencia y objeto, confesión) y sustanciales (interrumpe la prescripción, art.~2546 CCyC).}
\cornellsummary{La demanda contiene la pretensión, abre la instancia y debe cumplir el art.~330 con prueba documental y ofrecimiento de los demás medios. Produce efectos procesales (fija competencia y objeto litigioso) y sustanciales (interrumpe la prescripción), y exige mediación previa y tasa de justicia.}
\end{cornell}
\chapter{Notificación de la demanda y excepciones previas}
\label{cap:notificacion}

El proceso ordinario sigue una lógica: se notifica al demandado el inicio del juicio; este decide qué actitud adoptar; y, si ataca el proceso, la pretensión o el derecho del actor, aparecen las excepciones previas. Los conceptos de acción, pretensión y demanda ya fueron desarrollados (capítulos~\ref{cap:pedestal} y~\ref{cap:demanda}); aquí se agrega que la demanda, como acto de postulación, pone en marcha la maquinaria judicial y, si cumple los requisitos formales, inicia el proceso.

\section{El despacho de la demanda y el traslado}

Presentada la demanda ante el juzgado sorteado, este debe leerla y \textbf{proveerla}, es decir, dictarle una resolución o \textbf{auto}, con lugar, fecha y firma. Procesalmente, el juez tiene por presentado al actor como parte, por constituido el domicilio denunciado y ordena correr traslado al demandado; así el actor se incorpora como contendiente en un debate racional y controlable por las partes.

Correr traslado es \textbf{llamar al demandado} para comunicarle el juicio, a fin de que comparezca y dé su versión de los hechos, para determinar si lo afirmado por el actor es verdadero o falso. Se vincula con el principio de igualdad: al ordenar el traslado, el juez notifica al demandado para que esté a derecho, y así la igualdad se materializa a través de los sistemas procesales.

\section{Tipos de notificación y traba de la litis}

\begin{itemize}
\item \textbf{Por ministerio de la ley o por nota (art.~133 CPCCN).} Las resoluciones quedan notificadas los martes y viernes, o el siguiente día hábil si alguno fuera feriado, aunque la parte no concurra; en la práctica se consulta el expediente y se deja constancia en el libro de asistencia. Hace nacer los plazos para impugnar.
\item \textbf{Tácita.} Se produce cuando la parte toma conocimiento efectivo de lo resuelto: quien concurre al juzgado, anota que se corre traslado y prepara la cédula en el estudio queda tácitamente notificado del despacho.
\item \textbf{Por cédula (art.~135 CPCCN).} Entre las resoluciones que deben notificarse personalmente o por cédula figura el traslado de la demanda. La \textbf{cédula} es un formulario estructurado por disposición de la Corte o del reglamento (no el documento de identidad).
\end{itemize}

El circuito de la cédula es: el abogado del actor copia el auto y la prepara en el estudio; se presenta en el juzgado junto con copias de la demanda y documentación; pasa a la Oficina de Mandamientos y Notificaciones, que la distribuye entre los oficiales notificadores según las zonas de la Ciudad; el oficial concurre al domicilio denunciado y deja la cédula. Si no encuentra al demandado la devuelve informando que no pudo notificarlo, y el actor presenta un escrito titulado \textbf{«Denuncia nuevo domicilio»}.

\begin{regla}{art.~145 CPCCN: edictos y radiodifusión}
Si se trata de personas inciertas o de domicilio ignorado, la notificación se hace por edictos publicados por dos días en el Boletín Oficial y en un diario de los de mayor circulación del lugar del último domicilio conocido, o de la capital de la provincia, o de la Capital Federal en su caso.
\end{regla}

Para pedirlos debe declararse \textbf{bajo juramento} que se hicieron todas las gestiones al alcance para notificar en el domicilio de la demanda. El Código prevé también la radiodifusión o televisión, de uso muy infrecuente por su costo.

\subsection{Nulidad de la notificación}

Según el art.~149 CPCCN la nulidad de la notificación se rige por los principios de convalidación (art.~172: no puede pedirse la nulidad si la propia conducta de quien la invoca convalidó lo actuado) y de trascendencia (art.~173: no puede impugnarse si no se demuestra el perjuicio causado). Se declara en tanto se vea afectado el derecho de defensa del demandado, y se plantea por \textbf{incidente de nulidad}, un pequeño proceso dentro del principal.

\begin{importante}
\textit{Pas de nullité sans grief}: no hay nulidad sin perjuicio. Siempre debe señalarse el perjuicio concreto que sufrió la parte.
\end{importante}

\subsection{La traba de la litis: citación y emplazamiento}

Un error frecuente, que se remonta a Alsina, sostiene que la traba de la litis ocurre con la contestación de la demanda; no es tan grave, pero conceptualmente es incorrecto.

\begin{definicion}{Traba de la litis}
Se produce con la \textbf{notificación del traslado de la demanda}. Allí nacen dos cargas para el demandado: la \textbf{citación} (comparecer y estar a derecho) y el \textbf{emplazamiento} (el plazo para contestar, generalmente 15 días en el proceso ordinario).
\end{definicion}

Ello explica que el Código estructure un régimen de rebeldía. En sentido sustancial, la litis también puede entenderse trabada con la contestación, que enmarca el objeto del litigio: una pretensión del actor y una defensa que la resiste.

\section{Actitud del demandado frente a la demanda}

\begin{table}[H]
\centering
\renewcommand{\arraystretch}{1.25}
\begin{tabularx}{\textwidth}{>{\RaggedRight\arraybackslash}X>{\RaggedRight\arraybackslash}X}
\toprule
\textbf{Posturas omisivas} & \textbf{Posturas positivas} \\
\midrule
Rebeldía: no comparece ni contesta. & Contestar la demanda. \\
Comparece pero no contesta (pierde el derecho). & Reconvenir (contrademandar). \\
 & Allanarse. \\
 & Oponer excepciones previas. \\
\bottomrule
\end{tabularx}
\end{table}

La \textbf{rebeldía} es la situación del demandado que no observa ninguna de las cargas; si el actor lo pide, el juez puede declararlo rebelde, con consecuencias negativas. No se distingue entre rebeldía y contumacia. Quien comparece pero no contesta no es rebelde, pero pierde el derecho a contestar por \textbf{preclusión} (el proceso siempre avanza): sigue presente en el proceso.

Al \textbf{contestar}, el demandado puede negar los hechos, agregar documentos, invocar hechos distintos o dar otra versión, o reconvenir. La \textbf{reconvención} es la contrademanda: por economía procesal, quien reconviene asume el rol de actor y el actor el de demandado, siempre que se trate de la misma relación jurídica, para resolver todo en un solo proceso. \textbf{Allanarse} es someterse a la pretensión del actor sin debatir, sin que implique reconocer todos los hechos y el derecho; el allanamiento total hace carecer de sentido al proceso, que termina.

\section{Las excepciones de previo y especial pronunciamiento}

\begin{definicion}{Excepciones de previo y especial pronunciamiento}
Mecanismos defensivos con nombre propio, reunidos por el codificador en los arts.~346, 347 y 348, que tienen un \textbf{trámite previo} y requieren un \textbf{pronunciamiento especial} para que el proceso continúe.
\end{definicion}

Según los textos citados en las notas: el art.~346 fija la oportunidad (se deducen junto con la contestación de demanda); el art.~347 enumera las excepciones admisibles (incompetencia, falta de personería, falta de legitimación manifiesta, litispendencia, defecto legal, cosa juzgada, transacción, conciliación, desistimiento del derecho, prescripción, arraigo); y el art.~348 trata la suspensión del plazo para contestar en los casos de falta de personería, defecto legal y arraigo. % TODO: verificar consistencia entre fuentes originales: en otros pasajes el arraigo se ubica en los arts. 346 y 348, y la suspensión del plazo se atribuye a la parte final del art. 346.

Palacio las distingue en \textbf{dilatorias} (dilatan el trámite hasta ser resueltas) y \textbf{perentorias} (lo culminan definitivamente). Las clasificaciones no son verdaderas ni falsas, sino más o menos útiles. Se propone esta, según el objeto de ataque:

\begin{table}[H]
\centering
\renewcommand{\arraystretch}{1.25}
\begin{tabularx}{\textwidth}{>{\RaggedRight\arraybackslash}p{0.24\textwidth}>{\RaggedRight\arraybackslash}p{0.24\textwidth}>{\RaggedRight\arraybackslash}X}
\toprule
\textbf{Categoría} & \textbf{Qué se ataca} & \textbf{Ejemplos} \\
\midrule
Impedimentos procesales & La constitución válida del proceso & Incompetencia, falta de personería, litispendencia, defecto legal, arraigo \\
Excepciones propiamente dichas & La pretensión ejercida & Prescripción (puro derecho), falta de legitimación manifiesta, defensas temporarias \\
Defensas previas & El derecho troncal del actor & Cosa juzgada, transacción, conciliación, desistimiento \\
\bottomrule
\end{tabularx}
\end{table}

\subsection{Impedimentos procesales: los presupuestos procesales}

En 1868 Von Bülow interpretó el proceso como una relación jurídica nueva, distinta de la sustancial, con la aparición de un tercero, el juez, y creó la teoría de los \textbf{presupuestos procesales}. El Código los toma en sentido negativo como excepciones previas:

\begin{table}[H]
\centering
\renewcommand{\arraystretch}{1.2}
\begin{tabularx}{\textwidth}{>{\RaggedRight\arraybackslash}X>{\RaggedRight\arraybackslash}X}
\toprule
\textbf{Presupuesto procesal} & \textbf{Excepción previa} \\
\midrule
Juez competente & Incompetencia \\
Demanda clara y precisa & Defecto legal \\
No existe otro proceso igual & Litispendencia \\
Representación suficiente & Falta de personería \\
Domicilio o bienes en el país & Arraigo \\
\bottomrule
\end{tabularx}
\end{table}

\begin{itemize}
\item \textbf{Incompetencia.} Se deduce por declinatoria (al contestar, ante el juez de la causa) o por inhibitoria (ante el juez que se considera competente); son vías excluyentes (capítulo~\ref{cap:competencia}).
\item \textbf{Defecto legal.} El demandado afirma que no sabe qué se le reclama porque la demanda no es clara.
\item \textbf{Litispendencia.} Procede si hay otro proceso anterior con el mismo objeto, la misma causa y las mismas partes; se paraliza el segundo, que se acumula al primero.
\item \textbf{Falta de personería.} A quien interviene le falta la representatividad que invoca. («Persona» proviene de las máscaras del teatro griego; de allí «personero».) Se acredita, por ejemplo, el cargo y el acta en una asociación civil; el carácter de administrador en un consorcio; el poder en una persona física; la partida de nacimiento para representar a un menor. Si un chico rompe el vidrio del vecino, su padre actúa en su representación porque el menor no tiene capacidad procesal. Sin representación debida no puede dictarse sentencia: el juez puede advertir la falta hasta la sentencia y ordenar su subsanación.
\item \textbf{Arraigo.} Se vincula con quien domicilia en el extranjero, tiene un conflicto en el país y carece de bienes o domicilio aquí. El art.~2600 CCyC no permite imponer la caución de arraigo al demandado domiciliado o residente habitual en la República, de modo que la excepción quedó restringida a casos muy puntuales (véase «Réquiem para el arraigo», del Dr.~Rojas, en su sitio web). % TODO: verificar consistencia entre fuentes originales: el arraigo se ubica en los arts. 346 y 348, en el art. 348 en el cuadro de presupuestos y en el art. 347 en la enumeración.
\end{itemize}

\subsection{Excepciones propiamente dichas}

Si el proceso está válidamente constituido no hay impedimento procesal; el demandado invoca hechos \textbf{impeditivos o extintivos}. Ya no se ataca el proceso sino la pretensión o el derecho troncal.

\begin{itemize}
\item \textbf{Prescripción liberatoria:} el derecho eventualmente existe, pero no puede ejercerse porque no se demandó en el plazo. Un trabajador despedido hace cuatro años reclama indemnización cuando la ley laboral fija dos años: la empresa opone prescripción. Tramita como de \textbf{puro derecho} si no necesita prueba y el proceso concluye; si no, el juez difiere el pronunciamiento a la sentencia definitiva, donde la trata en primer lugar.
\item \textbf{Falta de legitimación manifiesta:} el demandado señala que no integra la relación jurídica invocada, que corresponde a otra persona. Si es manifiesta, el juez la resuelve de inmediato y el proceso se extingue; si no, difiere el pronunciamiento. % TODO: verificar consistencia entre fuentes originales: el índice menciona legitimación activa y pasiva, pero el desarrollo describe únicamente la pasiva.
\item \textbf{Defensas temporarias (art.~347 inc.~8):} no atacan el proceso ni la pretensión, sino que señalan que momentáneamente no puede ejercerse hasta cumplirse ciertas condiciones: el beneficio de excusión (art.~1583 CCyC: primero se ejecuta al deudor principal que al fiador que no renunció al beneficio); las condenaciones del posesorio (art.~2272 CCyC: no puede iniciarse la acción real sin cumplir las condenaciones de la posesoria); los días de llanto y luto (art.~2289 CCyC: no puede reclamarse a los herederos lo que adeudaba el causante hasta que transcurra ese período); y las condenaciones del ejecutivo (art.~553 CPCCN: quien perdió el ejecutivo y quiere iniciar un ordinario posterior debe cumplirlas antes). El texto citado incluye la excepción de incumplimiento contractual, que no se desarrolla en la explicación. % TODO: verificar consistencia entre fuentes originales.
\end{itemize}

\subsection{Defensas previas}

Se ataca el \textbf{derecho troncal}: el actor no tiene derecho a reclamar porque su derecho dejó de existir. Son la \textbf{cosa juzgada} (ya fue juzgado por sentencia firme), la \textbf{transacción} (acuerdo extrajudicial con renuncia a reclamos futuros), la \textbf{conciliación} (acuerdo homologado ante el juez, en mediación o en lo laboral, que terminó el proceso de modo anormal) y el \textbf{desistimiento del derecho} (la parte se apartó del derecho que tenía, en el proceso o extrajudicialmente).

\subsection{Cómo se deducen y sus efectos}

Las excepciones se deducen en un solo escrito, dentro del plazo para contestar la demanda (quince días). Solo suspenden ese plazo la \textbf{falta de personería}, el \textbf{defecto legal} y el \textbf{arraigo}. Desde una perspectiva crítica, se sostiene que el defecto legal debería ser la única, pues si no se sabe qué se reclama no puede contestarse; si se rechazara y los quince días ya hubieran vencido, el demandado quedaría sin plazo.

\begin{importante}
Consejo práctico: se aprende estudiando, pero se ejerce pensando. Como el defecto legal suspende el plazo, puede plantearse el segundo o tercer día del traslado y pedir que se suspenda el proceso hasta resolverlo, ganando tiempo: si prospera, el juez manda corregir la demanda y corre nuevo traslado; si se rechaza, aún quedan días del plazo original.
\end{importante}

Efectos de la resolución: \textbf{extinción del proceso} (vicio insanable: incompetencia de distinta jurisdicción, prescripción de puro derecho, cosa juzgada); \textbf{remisión del expediente} (incompetencia dentro de la misma jurisdicción); \textbf{subsanación y continuación} (defecto legal, falta de personería, bajo apercibimiento); y \textbf{diferimiento} a la sentencia definitiva cuando no puede tramitar como de puro derecho.

\begin{cornell}{notificación de la demanda y excepciones}
\cqa{¿Cuándo se traba la litis y qué cargas nacen?}{Con la notificación del traslado: citación (comparecer) y emplazamiento (contestar en 15 días).}
\cqa{¿Qué formas de notificación existen?}{Por ministerio de la ley (art.~133), tácita, por cédula (art.~135), por edictos o radiodifusión (art.~145, previo juramento).}
\cqa{¿Cuándo es nula una notificación?}{Cuando se afecta la defensa y se demuestra perjuicio (art.~173), sin convalidación (art.~172); \textit{pas de nullité sans grief}.}
\cqa{¿Qué actitudes puede adoptar el demandado?}{Omisivas (rebeldía; comparecer sin contestar) y positivas (contestar, reconvenir, allanarse, oponer excepciones).}
\cqa{¿Cómo se clasifican las excepciones según lo que atacan?}{Impedimentos procesales (proceso), excepciones propiamente dichas (pretensión) y defensas previas (derecho troncal).}
\cqa{¿Cuáles suspenden el plazo para contestar?}{Falta de personería, defecto legal y arraigo.}
\cqa{¿Qué hace el juez si la prescripción o la falta de legitimación no son de puro derecho?}{Difiere el pronunciamiento a la sentencia definitiva.}
\cornellsummary{Notificado el traslado queda trabada la litis y el demandado debe optar entre omitir o actuar. Las excepciones previas atacan el proceso, la pretensión o el derecho del actor, se oponen en un solo escrito con la contestación y producen extinción, remisión, subsanación o diferimiento.}
\end{cornell}
\chapter{La contestación de la demanda}
\label{cap:contestacion}

El proceso se diagrama en etapas por un sistema de \textbf{preclusión} (y no de unicidad): las partes pueden ejercer ciertos actos en cada etapa, en determinada forma y tiempo, y si no lo hacen pierden la posibilidad de hacerlo después; el proceso avanza hacia la sentencia y no se retrotrae, salvo planteo de nulidad. Sus etapas son la introductoria (demanda, traslado, contestación, excepciones), la probatoria, la conclusional (alegatos), la sentencia, la recursiva (apelación ante la Cámara) y la de ejecución. Antes del proceso existen la negociación directa y la mediación previa y obligatoria (capítulo~\ref{cap:adr}); la demanda se notifica generalmente por cédula al domicilio real (también por edictos, radiodifusión, personalmente o por acta notarial), y la cédula informa el juzgado, la carátula y el número de expediente. Con la notificación nacen las cargas de comparecer (citación) y contestar (emplazamiento); a diferencia del deber, el no ejercicio de la carga no trae sanción, aunque puede beneficiar al adversario. Las actitudes posibles del demandado se describieron en el capítulo~\ref{cap:notificacion}; este capítulo desarrolla la contestación y las actitudes especiales de reconvenir, allanarse y ser rebelde.

\section{Concepto, plazo y requisitos}

\begin{definicion}{Contestación de demanda}
Acto procesal de \textbf{postulación} en el cual el demandado se defiende, pidiendo a la jurisdicción que \textbf{rechace la pretensión} del actor. Es el acto de defensa por excelencia del proceso ordinario, donde ejerce todas sus defensas de modo amplio.
\end{definicion}

\begin{table}[H]
\centering
\renewcommand{\arraystretch}{1.25}
\begin{tabularx}{\textwidth}{>{\RaggedRight\arraybackslash}p{0.27\textwidth}>{\RaggedRight\arraybackslash}X}
\toprule
\textbf{Aspecto} & \textbf{Regla} \\
\midrule
Plazo & Proceso ordinario: 15 días hábiles; sumarísimo: 5 días hábiles, desde el día siguiente a la notificación. Es \textbf{perentorio, improrrogable e individual} para cada demandado. Si el domicilio está a más de 100 km, se extiende un día cada 200 km o fracción mayor a 100 km. \\
\addlinespace
Plazo de gracia (art.~156) & El escrito no presentado dentro del horario judicial del día del vencimiento puede entregarse el día hábil siguiente durante las dos primeras horas de despacho (entre 7:30 y 9:30). % TODO: verificar consistencia entre fuentes originales: la unidad sobre la demanda atribuye el plazo de gracia al art. 124 CPCCN y esta unidad al art. 156.
\\
\addlinespace
Aspectos fiscales & No se abona tasa de justicia (ley 23.898) salvo que se reconvenga; sí el bono profesional del Colegio de Abogados en la primera presentación. % TODO: verificar consistencia entre fuentes originales: la unidad sobre la demanda cita la «ley 3.898».
\\
\addlinespace
Mediación previa & Solo es exigible al actor; el demandado que advierte que no participó puede pedir que se suspenda el proceso y se reabra la mediación. \\
\addlinespace
Forma & Escrito en idioma nacional, con firma de letrado y/o parte, denuncia del domicilio real y constitución de domicilio procesal y electrónico. \\
\bottomrule
\end{tabularx}
\end{table}

El demandado puede presentarse (a) por sí con patrocinio letrado (firman ambos), (b) por apoderado (firma solo el abogado) o (c) como \textbf{gestor} (art.~48 CPCCN), cuando no hay tiempo de obtener el poder y hay urgencia, con 40 días hábiles para acreditar la representación o que la parte ratifique lo actuado; si no, puede decretarse la nulidad y tenerse por no contestada la demanda.

Dentro de las 24 horas de la presentación en papel el letrado debe digitalizarla (escrito y documental). Si no lo hace, se intima; si tampoco cumple en el nuevo plazo (24 a 48 horas), se tiene por no presentado el escrito, con la grave consecuencia de tener por no contestada la demanda. Es una exigencia reciente, derivada del expediente electrónico y análoga a la del art.~120.

\section{Contenido (art.~356 CPCCN)}

\subsection{Reconocimiento o negación de hechos y documentos (inc.~1)}

\begin{regla}{art.~356 inc.~1 CPCCN}
El demandado debe reconocer o negar categóricamente cada hecho expuesto en la demanda y la autenticidad de los documentos que se le atribuyan y la recepción de cartas y telegramas. El silencio, las respuestas evasivas o la negativa meramente general pueden estimarse como reconocimiento de la verdad de los hechos pertinentes y lícitos; los documentos se tienen por reconocidos o recibidos, según el caso.
\end{regla}

\begin{table}[H]
\centering
\renewcommand{\arraystretch}{1.25}
\begin{tabularx}{\textwidth}{>{\RaggedRight\arraybackslash}p{0.2\textwidth}>{\RaggedRight\arraybackslash}X>{\RaggedRight\arraybackslash}X}
\toprule
\textbf{Objeto} & \textbf{Silencio o negativa genérica} & \textbf{Negación categórica} \\
\midrule
Hechos & El juez \textbf{puede} tenerlos por admitidos (facultad). & Quedan controvertidos y deberán probarse. \\
Documentos & El juez \textbf{debe} tenerlos por reconocidos y auténticos (obligación). & Queda desconocido; el actor deberá probar su autenticidad. \\
\bottomrule
\end{tabularx}
\end{table}

\begin{table}[H]
\centering
\renewcommand{\arraystretch}{1.25}
\begin{tabularx}{\textwidth}{>{\RaggedRight\arraybackslash}p{0.2\textwidth}>{\RaggedRight\arraybackslash}X>{\RaggedRight\arraybackslash}X}
\toprule
 & \textbf{Instrumento privado} & \textbf{Instrumento público} \\
\midrule
Ejemplos & Constancia médica, factura, remito. & Escritura pública ante escribano. \\
Eficacia probatoria & Menor; no hace plena fe por sí mismo. & Mayor; hace plena fe por sí mismo, pues intervino un oficial público. \\
Cómo desconocerlo & Basta negar categóricamente su autenticidad, uno por uno («niega la autenticidad en forma y contenido de la constancia médica\ldots»). & Debe impugnarse al contestar, pero no basta para restarle eficacia. \\
Cómo quitarle eficacia & La negación categórica es suficiente: el actor debe probar su autenticidad. & \textbf{Incidente de redargución de falsedad} (art.~395) dentro de los 15 días de haberlo impugnado; si no se promueve, conserva plena eficacia. \\
\bottomrule
\end{tabularx}
\end{table}

El incidente de redargución de falsedad no suspende el proceso principal salvo que el juez lo disponga. Si el demandado reconoce la firma atribuida en un instrumento privado, se tiene por auténtico todo su contenido.

\subsection{Versión de los hechos del demandado (inc.~2), derecho y petición}

El demandado debe especificar con claridad los hechos que alega como fundamento de su defensa, dando su propia versión como el actor lo hizo conforme al art.~330 (el médico en la mala praxis, para demostrar que no existió la conducta tachada; en el accidente de tránsito, cómo ocurrió realmente y por qué no es responsable). Debe además invocar con claridad el derecho (aunque rige \textit{iura novit curia} y el juez puede apartarse al sentenciar) y peticionar que se lo tenga por presentado, por parte, por contestada la demanda y por ofrecida la prueba, y que oportunamente se rechace la pretensión.

\subsection{Ofrecimiento de prueba (art.~333)}

El demandado acompaña la prueba documental que esté en su poder y ofrece los demás medios: la oportunidad son los escritos introductorios.

\begin{table}[H]
\centering
\renewcommand{\arraystretch}{1.25}
\begin{tabularx}{\textwidth}{>{\RaggedRight\arraybackslash}p{0.2\textwidth}>{\RaggedRight\arraybackslash}X}
\toprule
\textbf{Prueba} & \textbf{Requisitos al ofrecerla} \\
\midrule
Documental & Acompañarla si está en su poder. \\
Testimonial & Individualizar a los testigos (nombre, apellido, documento, profesión, domicilio) e indicar los fines y los extremos que se intentan acreditar (por ejemplo, un testigo presencial del accidente; un traumatólogo como testigo técnico). \\
Pericial & Indicar el tipo de perito (psicológica, médica, contable) y los puntos de pericia, es decir, el cuestionario. \\
\bottomrule
\end{tabularx}
\end{table}

\section{Efectos de la contestación}

Cierra la etapa introductoria; fija la competencia si no se la cuestionó; impide recusar sin causa (art.~14), excepcionar y reconvenir después, por preclusión; y fija la \textbf{cuestión litigiosa}: se identifican los hechos afirmados por cada parte, los controvertidos y los conducentes, que determinan qué se probará. Si el demandado acompañó documentos, el juez corre traslado al actor por cinco días, por cédula, para que se pronuncie conforme al art.~356 inc.~1, manteniendo la bilateralidad y la igualdad.

\section{Documentos nuevos, hechos no invocados y hechos nuevos}

Como regla, los escritos introductorios son la oportunidad de ofrecer prueba, pero el Código admite excepciones.

\begin{table}[H]
\centering
\renewcommand{\arraystretch}{1.25}
\begin{tabularx}{\textwidth}{>{\RaggedRight\arraybackslash}p{0.17\textwidth}>{\RaggedRight\arraybackslash}X>{\RaggedRight\arraybackslash}X>{\RaggedRight\arraybackslash}X}
\toprule
 & \textbf{Documentos nuevos (art.~335)} & \textbf{Hechos no invocados (art.~334)} & \textbf{Hechos nuevos (art.~365)} \\
\midrule
Qué son & Documentos de fecha posterior a la demanda o contestación, o anteriores que se jura no haber conocido. & Hechos relacionados con el relato del actor y la relación jurídica, no mencionados por este e introducidos por el demandado al contestar. & Hechos ocurridos o conocidos después de la contestación, relacionados con la cuestión. \\
Oportunidad en 1.ª instancia & Hasta los autos para sentencia. & El actor tiene 5 días desde la notificación del traslado para ampliar o modificar su prueba. & Hasta 5 días después de notificada la audiencia del art.~360. \\
En Cámara & Sí (art.~260). & No. & No. \\
\bottomrule
\end{tabularx}
\end{table}

\begin{ejemplo}{Documento de fecha posterior y hecho nuevo}
En una demanda por daños de una obra lindera, el actor reclamaba con un presupuesto; cuatro meses después paga las reparaciones y puede incorporar la factura. En una mala praxis que dañó el corazón, una nueva intervención quirúrgica posterior por la misma causa es un hecho nuevo sobreviniente.
\end{ejemplo}

Presentado un documento nuevo se corre traslado por cinco días, conforme al art.~356 inc.~1.

\section{El allanamiento}

\begin{definicion}{Allanamiento}
Acto procesal \textbf{unilateral} (sin necesidad de conformidad del actor) por el cual el demandado \textbf{se somete a la pretensión}, sin implicar que reconozca como verdaderos los hechos y el derecho invocados.
\end{definicion}

Es un modo \textbf{anormal} de terminar el proceso (art.~307 CPCCN). Puede ser \textbf{total} o \textbf{parcial}, \textbf{expreso} o \textbf{tácito} (por ejemplo, en un desalojo el demandado consigna las llaves en el expediente dentro del plazo de contestación). Puede hacerse en cualquier instancia antes de la sentencia, por consciencia del resultado probable, para evitar mayores costas o la incertidumbre de un proceso que puede durar años.

\section{La reconvención}

\begin{definicion}{Reconvención}
Demanda que el demandado dirige contra el actor dentro del mismo proceso ordinario, con una pretensión autónoma y nueva. Su fundamento es la \textbf{economía procesal}.
\end{definicion}

El actor es reconvenido y el demandado reconviniente; ambos tienen doble calidad y el juez dicta una única sentencia. Se deduce en el mismo escrito de contestación y plazo (arts.~357 y ss.), con los requisitos de la demanda; se da traslado al actor por 15 días para contestar y oponer defensas (incluidas las excepciones de los arts.~346 a 348). La etapa introductoria se cierra entonces con cuatro escritos (demanda, contestación, reconvención y contestación de la reconvención) y sigue la audiencia preliminar del art.~360.

Requisitos: los del art.~330; tasa de justicia sobre el monto reclamado (o beneficio de litigar sin gastos); identidad de sujetos; mismo tipo de proceso; proceso en trámite con traslado notificado (no en diligencias preliminares ni prueba anticipada); y relación o conexidad entre las pretensiones.

Si no se reconviene, el demandado puede iniciar un proceso nuevo (si no prescribió); ambos se \textbf{acumulan} ante el juez que previno (aquel en cuyo expediente se notificó primero la demanda) y se dicta una sola sentencia.

\begin{ejemplo}{Reconvención}
Compraventa: el comprador reclama la entrega y el vendedor reconviene por el saldo de precio. Accidente de tránsito: ambos alegan que el otro cruzó en rojo y reclaman daños. Alimentos: la madre pide aumentar una cuota de \$10.000 fijada cinco años antes y el padre reconviene pidiendo su reducción porque la madre, antes desempleada, hoy es gerente de una multinacional.
\end{ejemplo}

\section{La rebeldía}

\begin{definicion}{Rebeldía}
Situación del demandado debidamente notificado que \textbf{no comparece} y por ende \textbf{no contesta} (arts.~59 y ss. CPCCN); puede darse también cuando una parte presentada abandona el proceso.
\end{definicion}

\begin{table}[H]
\centering
\renewcommand{\arraystretch}{1.25}
\begin{tabularx}{\textwidth}{>{\RaggedRight\arraybackslash}p{0.2\textwidth}>{\RaggedRight\arraybackslash}X>{\RaggedRight\arraybackslash}X}
\toprule
 & \textbf{Comparece y no contesta} & \textbf{Rebeldía} \\
\midrule
Cargas & Cumple comparecer. & No cumple ninguna. \\
Notificaciones & Por cédula al domicilio electrónico constituido. & Por ministerio de la ley (martes y viernes), hasta decretarse la rebeldía; la declaración y la sentencia se notifican por cédula al domicilio real. \\
Actos futuros & Puede ejercer todos los pertinentes. & No puede ejercer actos de etapas anteriores a su presentación. \\
\bottomrule
\end{tabularx}
\end{table}

La rebeldía se decreta \textbf{a pedido del actor}, no de oficio, vencido el plazo de 15 días. Si no se pide, las resoluciones se notifican al demandado por ministerio de la ley. % TODO: verificar consistencia entre fuentes originales: el apunte indica que desde la declaración de rebeldía hasta la sentencia las resoluciones se notifican por ministerio de la ley y que la declaración y la sentencia se notifican por cédula al domicilio real.

\textbf{Efectos.} El proceso continúa con todas sus etapas (no se dicta sentencia automática ni se tienen los hechos por ciertos sin más); por el silencio ante el art.~356 inc.~1 el juez puede tener por admitidos los hechos y debe tener por reconocidos los documentos; hay presunción de verdad de los hechos del actor, a considerar en caso de duda; el actor puede pedir medidas cautelares (arts.~63 y 212, incluido el embargo preventivo), pues la rebeldía permite presumir verosimilitud y peligro en la demora; y las resoluciones se notifican por ministerio de la ley. Según el art.~63, el rebelde puede pedir la revocación de las cautelares si acredita que la rebeldía obedeció a un impedimento justificado que no pudo comunicar, y el juez aprecia en la sentencia los efectos de la rebeldía como elemento corroborante de las pruebas.

\textbf{Cesación.} El rebelde que se presenta solicita que cese su rebeldía y podrá ejercer los actos desde allí en adelante, sin poder contestar, excepcionar ni ofrecer prueba si ya precluyó. Antes de la etapa probatoria controla la prueba del actor; durante ella puede impugnar testimoniales y periciales. Las cautelares en principio se mantienen, pudiendo pedir su modificación, sustitución o disminución, y levantarse si demuestra que la rebeldía fue forzada por causas insuperables (por ejemplo, internación en estado de coma). Según el art.~346, solo puede oponer prescripción después del plazo de excepciones si justifica que su rebeldía obedeció a causas que no pudo superar.

\textbf{Nulidad de la notificación.} Si la cédula fue dirigida a un domicilio incorrecto, el demandado puede plantear la nulidad; si se hace lugar, el proceso se retrotrae, debe notificarse de nuevo el traslado, el demandado puede contestar y se levantan las cautelares trabadas.

\begin{cornell}{contestación, allanamiento, reconvención y rebeldía}
\cqa{¿Cuál es el plazo y qué requisitos tiene la contestación?}{15 días hábiles (5 en el sumarísimo), perentorio, improrrogable e individual, con gracia de dos horas; escrito con domicilios; digitalización en 24 horas bajo apercibimiento de tenerlo por no presentado.}
\cqa{¿Qué exige el art.~356 inc.~1 respecto de hechos y documentos?}{Reconocer o negar categóricamente. El silencio permite (hechos) u obliga (documentos) al juez a tenerlos por admitidos o reconocidos.}
\cqa{¿Cómo se desconoce un instrumento privado y uno público?}{El privado, con la negación categórica; el público, con incidente de redargución de falsedad dentro de 15 días (art.~395).}
\cqa{¿Cuándo pueden incorporarse documentos nuevos, hechos no invocados y hechos nuevos?}{Arts.~335, 334 y 365, con las oportunidades del cuadro comparativo.}
\cqa{¿Qué es el allanamiento?}{Sometimiento unilateral a la pretensión, total o parcial, expreso o tácito, modo anormal de terminar el proceso (art.~307).}
\cqa{¿Qué requisitos tiene la reconvención?}{Los de la demanda, tasa, identidad de sujetos, mismo tipo de proceso, proceso con traslado notificado y conexidad; se resuelve con una única sentencia.}
\cqa{¿Qué es la rebeldía y qué efectos tiene?}{No comparecer ni contestar; se decreta a pedido del actor; el proceso continúa, hay presunción de verdad, el actor puede pedir cautelares y las resoluciones se notifican por ministerio de la ley.}
\cqa{¿Qué ocurre si la notificación fue nula?}{El proceso se retrotrae, se notifica nuevamente, el demandado puede contestar y se levantan las cautelares.}
\cornellsummary{La contestación del art.~356 fija la cuestión litigiosa y cierra la etapa introductoria, con excepciones limitadas para documentos y hechos posteriores. Las actitudes especiales —allanamiento, reconvención y rebeldía— tienen requisitos y efectos propios, y la nulidad de la notificación retrotrae el proceso. Concluida la etapa introductoria, el proceso queda listo para la audiencia preliminar.}
\end{cornell}
\part{La audiencia preliminar y la prueba}

\chapter{La audiencia preliminar}
\label{cap:audiencia}

Concluida la etapa introductoria, el proceso ordinario llega a un acto que media entre ella y la etapa probatoria: la \textbf{audiencia preliminar}. Este capítulo explica su contexto histórico y cultural, su finalidad, el contenido del art.~360 CPCCN y los hechos que quedan exentos de prueba, con lo cual se prepara el estudio de la teoría general de la prueba (capítulo~\ref{cap:prueba}).

\begin{definicion}{Audiencia preliminar}
Acto procesal que media entre la etapa introductoria y la etapa probatoria del proceso ordinario.
\end{definicion}

No todo proceso tiene necesariamente etapa de prueba: depende de que haya hechos alegados, controvertidos y conducentes para resolver el conflicto.

\section{Antecedentes de la oralidad}

La implementación de la oralidad en el proceso civil, comercial y laboral, a diferencia del penal, tuvo una demora de muchos años y varios intentos fallidos. El sistema actual es predominantemente escrito; la oralidad se limita a la \textbf{oralidad actuada}: un empleado toma nota de lo que declara un testigo y labra un acta que se incorpora al expediente, forma de trabajo que no se ajusta a los desarrollos tecnológicos actuales, pues cualquier cámara puede registrar una declaración.

\begin{itemize}
\item \textbf{1924--1926:} el procesalista Jofre elaboró en la Universidad de Buenos Aires el primer proyecto de proceso por audiencias (de allí la expresión «100 años de demora»).
\item \textbf{1993:} primera implementación de la audiencia preliminar, que duró un año por la resistencia de los jueces y fue derogada.
\item \textbf{1995:} incorporación de una audiencia que debía tomar el juez.
\item \textbf{1996 en adelante:} se trabaja sobre la actual audiencia preliminar, incorporada al CPCCN con distintas variantes. % TODO: verificar consistencia entre fuentes originales: la audiencia preliminar figura con vigencia desde 1996, mientras que se menciona una audiencia en 1995 y otra en 1993.
\end{itemize}

\begin{table}[H]
\centering
\renewcommand{\arraystretch}{1.25}
\begin{tabularx}{\textwidth}{>{\RaggedRight\arraybackslash}X>{\RaggedRight\arraybackslash}X}
\toprule
\textbf{Código Nacional (CPCCN)} & \textbf{Código de la Provincia de Buenos Aires} \\
\midrule
Legisla la audiencia preliminar desde 1996 (más de 25 años), pero en la práctica muchos jueces no la toman o la delegan. & No incorporó la audiencia preliminar, pero mediante el Protocolo de Oralidad del Ministerio de Justicia muchos jueces y abogados la implementan con éxito. \\
\bottomrule
\end{tabularx}
\end{table}

\begin{importante}
Paradoja central: donde la norma existe hace más de 25 años la práctica es pobre; donde no existe, la práctica avanza por la voluntad de los operadores. El principal obstáculo es \textbf{cultural}, lo que obliga a repensar la formación universitaria: en universidades de Perú los estudiantes reciben además clases de teatro, oratoria y escritura, mientras que en la Facultad de Derecho argentina esas habilidades no son prioritarias.
\end{importante}

\section{Finalidad y régimen del art.~360}

El legislador busca \textbf{simplificar} un proceso cuya carga de alegaciones y de prueba supera la necesidad del caso. La causa es estructural: el actor debe ofrecer toda su prueba en la demanda sin saber qué hechos negará el demandado, y para cubrirse ofrece más de la necesaria, incluso sobre hechos que el demandado podría admitir. La audiencia depura ese exceso y concentra el proceso en los hechos realmente controvertidos. Conforme al art.~331, el actor puede modificar la demanda solo hasta que el primer demandado sea notificado; luego no puede ampliarla ni modificarla en sus elementos sustanciales.

\begin{regla}{art.~360 CPCCN (texto transcripto en las fuentes)}
El juez toma la audiencia preliminar con carácter indelegable. En ella debe: (1) invitar a las partes a una conciliación o a otra forma de solución del conflicto; (2) recibir la declaración de las partes en posiciones; (3) fijar los hechos articulados conducentes a la decisión del juicio; (4) decidir la apertura a prueba; (5) ordenar los medios de prueba conducentes; (6) declarar la cuestión de puro derecho, si correspondiere.
\end{regla}

% TODO: verificar consistencia entre fuentes originales: el texto del art. 360 transcripto no coincide con la numeración usada en el desarrollo y el índice de la clase (mediación prejudicial como inc. 2, hechos conducentes como inc. 3, confesional como inc. 4, medios de prueba como inc. 5 y puro derecho como inc. 6).

\subsection{Indelegabilidad y resistencia judicial}

La redacción de 1995 imponía la presencia del juez «bajo pena de nulidad»; en 2001 se sustituyó por el carácter indelegable y se habilitó a cualquiera de las partes a dejar constancia en el libro de asistencia del incumplimiento del juez y retirarse sin sanción. % TODO: verificar consistencia entre fuentes originales: las unidades sobre sujetos procesales y actos procesales indican que la audiencia del art. 360 debe ser presidida por el juez bajo pena de nulidad.
Según la experiencia relatada en el fuero civil nacional (con referencias al comercial), en 25 años no se declaró la nulidad de ninguna audiencia por ausencia del juez ni se sancionó a ninguno; en promedio los jueces toman a lo sumo la mitad de las audiencias. Los abogados no controlan porque no querrían como enemigo al juez que aún no ha fallado su caso; el control debería ejercerlo el \textbf{Consejo de la Magistratura}. En la Provincia de Buenos Aires funcionó mostrar a los operadores estadísticas que comparan los casos con juez presente y ausente, lo que generó una sana competencia.

\begin{importante}
El mayor beneficio de la oralidad no es la celeridad sino la \textbf{calidad del servicio}: la prueba obtenida con intervención directa del juez es de mayor valor que la producida ante un empleado.
\end{importante}

\section{Los actos de la audiencia}

\begin{description}[leftmargin=0pt,style=nextline]
\item[Conciliación.] El primer acto es invitar a conciliar, lo que exige que el juez haya leído el caso y trabaje con teoría del conflicto para ofrecer fórmulas (capítulo~\ref{cap:adr}). Si hay acuerdo, el juez lo \textbf{homologa}, vale como sentencia y es ejecutable. «Un mal acuerdo es mejor que un buen juicio».
\item[Reenvío a mediación prejudicial.] Una reforma posterior permite al juez reenviar el conflicto a mediación. Se la critica por inútil y por demorar al volver al punto de partida; el único antecedente conocido terminó con la Cámara revocando la decisión. Solo podría tener sentido en conflictos de extrema complejidad.
\item[Fijación de hechos conducentes.] Sin acuerdo, el juez fija los hechos articulados conducentes. Un \textbf{hecho controvertido} es el alegado por el actor y negado por el demandado; un \textbf{hecho conducente} es, además, aquel cuyo conocimiento es necesario para resolver. \textbf{Solo los hechos articulados, controvertidos y conducentes son objeto de prueba}; se descarta la prueba sobre los demás.
\item[Puro derecho (art.~359).] Si no hay hechos controvertidos y conducentes que acreditar (la discusión versa sobre la interpretación del derecho o los hechos no han sido desconocidos), el juez declara la causa de puro derecho y, firme la resolución, llama autos para sentencia. Puede hacerlo antes de la audiencia, en ella (art.~360) o en cualquier momento cuando resulte evidente. La secuencia lógica es intentar conciliar y, si no hay acuerdo, declarar el puro derecho, omitiendo la etapa probatoria.
\item[Prueba confesional.] La ley prevé que las partes absuelvan posiciones en la propia audiencia. Se critica su ubicación: la parte llega más nerviosa y menos concentrada en conciliar, y la prueba se produce antes de clarificar los hechos definitivamente negados; debería ubicarse al final, en la audiencia de prueba, cuando ya se produjeron los demás medios. Responde a asegurar la comparecencia personal, pues quien no asiste puede quedar confeso.
\item[Ordenamiento de los medios de prueba.] El juez ordena los medios admisibles y pertinentes, necesarios y útiles para probar un hecho controvertido y conducente, y debe concentrar la prueba testimonial en un solo acto para acortar plazos. Sin audiencia de prueba (que aún no existe en el ámbito nacional), la producción se realiza por impulso de cada parte, con considerable demora.
\end{description}

\section{Hechos eximidos de prueba}

\begin{table}[H]
\centering
\renewcommand{\arraystretch}{1.25}
\begin{tabularx}{\textwidth}{>{\RaggedRight\arraybackslash}p{0.27\textwidth}>{\RaggedRight\arraybackslash}X}
\toprule
\textbf{Categoría} & \textbf{Explicación y ejemplos} \\
\midrule
Hechos reconocidos & El hecho admitido por el demandado en su contestación queda exento: no hay controversia. \\
\addlinespace
Hechos presumidos legalmente & La negativa no basta: el demandado debe probar lo contrario o se invierte la carga. Un acta notarial tiene fe de conocimiento; para desvirtuarla hay que redargüirla de falsa y probar la falsedad. La plena capacidad a los 18 años se presume; debe probar la incapacidad quien la invoque. \\
\addlinespace
Hechos notorios & Conocidos por todos (en la clase se mencionó la pandemia y quién era el presidente). \\
\addlinespace
Hechos evidentes & Se conocen por simple percepción (el agua está mojada; a 100 grados se evapora; reglas físicas de aceptación científica universal); es difícil distinguirlos con precisión de los notorios. \\
\addlinespace
Hechos imposibles & Su acaecimiento es materialmente impensable: que un chancho pasó volando sobre la cabeza del actor y lo distrajo; que Superman causó un daño. \\
\addlinespace
Máximas de experiencia & Reglas extraídas de la observación del curso ordinario de las cosas: quien se arroja de un décimo piso probablemente muere al tocar el suelo. \\
\addlinespace
La ley & Se presume conocida por todos. Excepción: la \textbf{ley extranjera}, cuya existencia y contenido puede exigir el juez a la parte que la alega. \\
\bottomrule
\end{tabularx}
\end{table}

\begin{cornell}{audiencia preliminar y hechos eximidos de prueba}
\cqa{¿Qué es la audiencia preliminar y cuál es su finalidad?}{Acto entre la etapa introductoria y la probatoria; depura el exceso de alegaciones y prueba y concentra el proceso en los hechos controvertidos y conducentes.}
\cqa{¿Qué hace el juez según el art.~360?}{Invita a conciliar, recibe posiciones, fija hechos conducentes, abre a prueba, ordena los medios y declara el puro derecho cuando corresponde.}
\cqa{¿Qué hechos se prueban?}{Solo los articulados, controvertidos y conducentes.}
\cqa{¿Cuándo hay puro derecho (art.~359)?}{Cuando no hay hechos controvertidos; no se produce prueba y se llama autos para sentencia.}
\cqa{¿Por qué se critica la confesional en la audiencia?}{Por su nerviosismo, por distraer de la conciliación y por producirse antes de aclarar los hechos negados.}
\cqa{¿Qué hechos están eximidos de prueba?}{Reconocidos, presumidos, notorios, evidentes, imposibles, máximas de experiencia y la ley (salvo la extranjera).}
\cqa{¿Cuál es el principal obstáculo de la oralidad?}{El condicionamiento cultural, más que el jurídico; la clave es la calidad del servicio.}
\cornellsummary{La audiencia preliminar, indelegable para el juez según el art.~360, concilia, depura la prueba y fija los hechos controvertidos y conducentes que serán objeto de prueba, o declara el puro derecho. Su práctica enfrenta resistencias culturales y hay hechos que no requieren prueba.}
\end{cornell}
\chapter{Teoría general de la prueba}
\label{cap:prueba}

La teoría general de la prueba responde las preguntas básicas del sistema probatorio: qué es la prueba, para qué sirve, qué se prueba, quién prueba, cómo y cuándo se prueba y quién valora. Se apoya en lo visto sobre los hechos controvertidos y conducentes que fija la audiencia preliminar (capítulo~\ref{cap:audiencia}).

\section{Concepto de prueba}

En el uso cotidiano, probar es demostrar, comprobar o verificar. Aun restringido al proceso, el término es \textbf{multívoco}: alude al \textbf{objeto} que sirve para conocer el hecho, a la \textbf{actividad probatoria} según el medio utilizado y al \textbf{resultado} de esa actividad.

\subsection{La necesidad de probar}

Surgido el conflicto, las personas acuden al juez natural y competente mediante una \textbf{demanda}, que introduce una \textbf{pretensión} fundada en \textbf{hechos}. Por el principio dispositivo, los hechos los traen las partes, que son sus dueñas; el juez aplica el derecho a los hechos que ellas llevaron. Si el demandado los reconoce expresamente quedan exentos de prueba (relevo de prueba); si son controvertidos, deben demostrarse, pues el juez es ajeno a su conocimiento.

\begin{table}[H]
\centering
\renewcommand{\arraystretch}{1.25}
\begin{tabularx}{\textwidth}{>{\RaggedRight\arraybackslash}X>{\RaggedRight\arraybackslash}X}
\toprule
\textbf{Juez civil y comercial} & \textbf{Juez penal} \\
\midrule
Rige el principio dispositivo: depende de los hechos y pruebas aportados por las partes; no puede usar su saber privado, aun si hubiera presenciado un hecho. & Rige el principio inquisitivo: investiga los hechos más allá de lo que diga el imputado y busca la verdad jurídica objetiva. \\
Las partes ofrecen los medios; el juez es árbitro del proceso. & Puede impulsar la investigación de manera autónoma. \\
\bottomrule
\end{tabularx}
\end{table}

\begin{definicion}{Prueba (Falcón)}
Demostración en juicio de la ocurrencia de un suceso.
\end{definicion}

La definición alude al \textbf{resultado}. La demostración es una actividad compleja que varía según la prueba: técnica (perito), relato (testigo) o constancias de un documento. «Suceso» es la serie de hechos controvertidos que integran el \textit{thema decidendum} de la litis.

\section{Objeto de la prueba}

\subsection{¿Hechos o afirmaciones de hechos?}

Alsina sostenía a mediados del siglo anterior que se probaban los hechos. Pero el hecho llega al proceso según cómo impactó en una persona o un objeto y según el relato de la demanda; como el juez civil no investiga la verdad objetiva y se contenta con lo que las partes relatan en los escritos constitutivos, hay un alejamiento respecto del hecho en sí.

\begin{importante}
En el proceso civil se prueban las \textbf{afirmaciones de las partes relativas a los hechos}, no los hechos en sí. El hecho pertenece al pasado y solo puede reconstruirse a través de las huellas que dejó.
\end{importante}

Para Niño Zárate la prueba construye un \textbf{puente} entre el hecho histórico y el proceso: el juez no ve la realidad, pero sí su impronta en una cosa o en la psiquis de una persona. Consultando la \textbf{fuente de prueba} (cosas o personas) a través del \textbf{medio de prueba} más directo y adecuado, se reconstruye el hecho.

\subsection{Requisitos de los hechos: afirmados, conducentes, controvertidos y admisibles}

\begin{description}[leftmargin=0pt,style=nextline]
\item[Afirmados.] Deben afirmarse asertivamente (positivos o negativos) en los escritos constitutivos: demanda, contestación, reconvención y su contestación, excepciones y su contestación, e incidentes. Los hechos son la causa de la pretensión y cumplen una \textbf{triple función}: fundamento de la pretensión, objeto de prueba y fundamento de la sentencia. No deben afirmarse los \textbf{hechos secundarios}, que prueban los principales y se descomponen en afirmaciones menores (conviene incluirlos, pero su omisión no debería impedir preguntas o puntos de pericia sobre ellos), % TODO: verificar consistencia entre fuentes originales: la oración sobre objeciones a preguntas o puntos de pericia relativos a hechos secundarios resulta ambigua en el apunte.
ni los \textbf{hechos sobrevinientes} debidamente probados: según el art.~163 inc.~6, el juez considera al sentenciar los hechos constitutivos, modificativos o extintivos del derecho ocurridos durante el proceso y probados, aunque no hayan sido afirmados oportunamente. Por ejemplo, si se acredita que el demandado pagó al mandatario del actor con poder suficiente, ese hecho, definitorio de la pretensión, debe considerarse.
\item[Conducentes.] Deben tener peso, relevancia y trascendencia para fundar la sentencia; los \textbf{inconducentes}, aunque afirmados, son ajenos a la decisión y puede prescindirse de ellos.
\item[Controvertidos.] La controversia se produce al correr traslado de la demanda: el demandado tiene la carga del art.~356 de reconocer o negar categóricamente los hechos, y si no lo hace el juez puede tenerlos por reconocidos. Según el art.~359, si la cuestión es de puro derecho se llama autos para sentencia; si hay hechos controvertidos o no admitidos expresamente, el juez abre a prueba aunque las partes no lo pidan. La \textbf{admisión} comprende el silencio, la negativa genérica («niego todos y cada uno de los hechos que no sean objeto de especial reconocimiento»), las respuestas evasivas y la falta de contestación.
\item[Admisibles.] Deben ser posibles de probar: quien reclama por incumplimiento de un contrato de objeto ilícito no puede probar ese hecho. El Código excluye además ciertos medios: en alimentos, para acreditar el caudal económico del alimentante solo prueba informativa y documental; en desalojo por falta de pago y vencimiento del contrato, el demandado solo puede ofrecer prueba pericial, documental y confesional.
\end{description}

\subsection{La prueba del hecho negativo}

Quien afirma un hecho negativo como fundamento puede probarlo por el \textbf{afirmativo contrario o distinto}: si se reclama porque la casa no se pintó de verde como decía el contrato, se prueba que se la pintó de rojo. En cambio, alegar que la contraparte «nunca asesoró» en un contrato de asesoramiento es un hecho negativo indefinido, de \textbf{prueba diabólica}, porque se descompone en una serie casi infinita de afirmaciones; al demandado le basta demostrar que cumplió, aunque sea una vez. Cuando la carga se traslada a la parte mejor posicionada para probar, se aplica la \textbf{carga dinámica de la prueba}.

\section{Hechos exentos de prueba (profundización)}

Las categorías generales se enunciaron en el capítulo~\ref{cap:audiencia}. Se agregan estas precisiones:

\begin{itemize}
\item \textbf{Hechos notorios:} conocidos, en un lugar, tiempo y sector de la comunidad, con uniformidad sobre cómo y que sucedieron; quedan exentos no porque el juez los conozca (no usa su saber privado) sino por ser de público conocimiento donde tramita o sucedió el hecho. La cuarentena, con sus restricciones, era notoria en el lugar donde se aplicaba y podía no serlo en otra provincia que la administrara distinto. El carácter notorio no releva de afirmar el hecho ni de explicar su relación con la pretensión.
\item \textbf{Hechos evidentes:} similares a los notorios pero \textbf{generalizables y universalizables} («el todo es mayor que la parte», la ley de gravedad: quien arroja un armario desde cuatro pisos sabe que caerá sobre el auto del vecino).
\item \textbf{Hechos presumidos legalmente:} la presunción invierte la carga. Las \textbf{\textit{iuris tantum}} admiten prueba en contrario: el dueño o guardián de la cosa responde del daño que esta cause (art.~1758 CCyC) y se exime acreditando que fue usada contra su voluntad expresa o presunta. Las \textbf{\textit{iuris et de iure}} no la admiten: la plena capacidad a los 18 años, si no hubo declaración de incapacidad ni restricción judicial.
\item \textbf{Hechos reconocidos expresamente:} el reconocimiento del demandado sobre lo afirmado por el actor es una confesión plena que fija el hecho para el proceso; la excepción son los \textbf{hechos indisponibles}, que deben probarse aunque ambas partes los confiesen.
\end{itemize}

\section{Finalidad de la prueba}

Se procura la verdad material o real, o acercarse a ella, para formar el \textbf{convencimiento del juez}. Aunque el norte sea la \textbf{verdad jurídica objetiva} (según la Corte), el juez se contenta con la mayor certeza posible: una \textbf{verdad formal}, regida por las leyes jurídicas de la producción de la prueba, que con un alto grado de probabilidad le permita fijar los hechos, subsumirlos en la norma y dictar la sentencia que es la norma particular del caso.

\section{La carga de la prueba}

\begin{regla}{art.~377 CPCCN}
Incumbe la carga de la prueba a quien afirme la existencia de un hecho controvertido o de un precepto jurídico que el tribunal no tenga el deber de conocer. Cada parte debe probar el presupuesto de hecho de la norma o normas que invoque como fundamento de su pretensión, defensa o excepción.
\end{regla}

La carga es un \textbf{imperativo del propio interés}: durante el proceso no hay obligaciones en sentido estricto; quien no la ejerce no es compelido pero soporta la consecuencia procesal perjudicial. Las reglas de la carga impiden que el juez \textbf{absorba la instancia} («no puedo fallar por falta de elementos»): debe resolver; si el actor no probó lo que debía, se rechaza la demanda, y si quien no probó es el demandado, se hace lugar.

Rige el \textbf{principio de adquisición probatoria}: producida la prueba, no importa quién la ofreció; se adquiere para el proceso, por lo que una parte no podría desistir de un testigo que ya declaró, aunque lo haya hecho en su contra.

\begin{ejemplo}{Distribución de la carga según la actitud del demandado}
Una parte reclama el pago de lo trabajado conforme a un contrato. Si el demandado niega categóricamente el vínculo y la firma, el actor debe probar el antecedente de hecho (el vínculo y la obligación), y el demandado, que solo niega, no debe probar la negativa. Si reconoce el contrato (de locación de obra para hacer un banco de plaza) y dice «le pagué», el actor queda eximido y el demandado debe acreditar el hecho extintivo, el pago.
\end{ejemplo}

\textbf{Prueba del derecho.} Quien invoca un precepto que el juez no debe conocer (una ordenanza no publicada) debe probar su existencia, situación hoy muy difícil dadas las tecnologías de acceso a la legislación. El \textbf{derecho extranjero} debe acreditarlo, en principio, quien lo alega, sin perjuicio de que el juez investigue su existencia si no lo hace; se obtiene por oficio a representaciones diplomáticas o consulares o por consulta a juristas expertos. Existe un tratado con Uruguay por el cual los dos países no exigen prueba del derecho respectivo. % TODO: verificar referencia presente en las fuentes originales: la referencia «22.411» y la redacción del efecto del tratado resultan imprecisas.

\section{Producción y valoración}

\begin{table}[H]
\centering
\renewcommand{\arraystretch}{1.25}
\begin{tabularx}{\textwidth}{>{\RaggedRight\arraybackslash}p{0.2\textwidth}>{\RaggedRight\arraybackslash}X}
\toprule
\textbf{Momento} & \textbf{Descripción} \\
\midrule
Ofrecimiento & Las partes expresan con qué medios acreditarán los hechos (demanda, contestación y escritos constitutivos). \\
Producción & Se efectiviza dentro del plazo probatorio, de 40 días desde la audiencia preliminar: declara el testigo, se presenta el dictamen, se diligencian oficios. \\
Valoración & El juez pondera el peso de cada medio para formar su convicción al sentenciar, según la sana crítica (art.~386); las partes también valoran en los alegatos. \\
\bottomrule
\end{tabularx}
\end{table}

\subsection{Medios de prueba, pertinencia y admisibilidad}

Los medios de prueba son los instrumentos que ofrece el proceso para consultar las fuentes y extraer la información que reconstruya el hecho. Falcón los clasifica en \textbf{previstos} (en el Código) y \textbf{no previstos} (se usan por analogía, según lo determine el juez), y por su fuente en pruebas por documentación, por declaración y por información. Según el art.~378, se producen por los medios previstos y por los que el juez disponga a pedido de parte o de oficio, siempre que no afecten la moral, la libertad personal de los litigantes o de terceros ni estén prohibidos.

La \textbf{pertinencia} (art.~364) significa que los hechos a probar pertenecen al proceso por haber sido afirmados; si no, son \textbf{impertinentes} y la prueba debe rechazarse. La \textbf{admisibilidad} apunta al \textbf{medio de prueba}: es inadmisible si se refiere a hechos cuya prueba prohíbe la ley, por cuestiones morales (secretos de Estado o profesionales) o por no ser adecuado para el fin propuesto. El art.~364 impide producir pruebas sobre hechos no articulados y rechazar las manifiestamente improcedentes, superfluas o dilatorias. Se señala además el art.~397 como norma extrapolable para exigir que los medios consulten la fuente original e inmediata. % TODO: verificar referencia presente en las fuentes originales: el alcance del art. 397 y de su «extrapolación» no está desarrollado en el apunte.

\begin{ejemplo}{Impertinencia e inadmisibilidad}
Se ofrecen testigos para demostrar que el demandado no era apto para conducir, hecho no afirmado en la demanda. Hay un vicio de pertinencia (no se afirmó) y de admisibilidad (aun afirmado, la testimonial no es idónea para acreditar la aptitud para manejar; correspondía una pericial).
\end{ejemplo}

Rige el principio \textbf{\textit{favor probationis}}: ante la duda entre abrir a prueba o declarar el puro derecho, el juez abre a prueba, pues a lo sumo la prueba no será considerada, y eso es preferible a lesionar el derecho de defensa.

\subsection{Negligencia y caducidad de la prueba}

El plazo de prueba es de \textbf{cuarenta días} desde la audiencia preliminar y es común. Si las partes no instan la producción, pueden incurrir en negligencia o caducidad, instituciones distintas de consecuencias similares.

\begin{table}[H]
\centering
\renewcommand{\arraystretch}{1.25}
\begin{tabularx}{\textwidth}{>{\RaggedRight\arraybackslash}p{0.18\textwidth}>{\RaggedRight\arraybackslash}X>{\RaggedRight\arraybackslash}X}
\toprule
 & \textbf{Negligencia} & \textbf{Caducidad} \\
\midrule
Factor & Subjetivo: desidia o culpa de la parte, con demora injustificada. & Objetivo: vencimiento de plazos del Código, sin evaluar culpa o dolo. \\
Iniciativa & Solo a pedido de parte. & De oficio o a pedido de parte. \\
Trámite & Incidente con traslado a la acusada. & No se sustancia; se constata y decreta. \\
Recursos & Si se acredita la producción antes de contestar, se rechaza la acusación (irrecurrible); si se declara, es inapelable para el afectado. & Inapelable una vez decretada. \\
Replanteo & En cámara, si es juicio ordinario y se apela la sentencia. & Ídem. \\
\bottomrule
\end{tabularx}
\end{table}

\begin{ejemplo}{Prueba informativa}
Se ordena librar un informe al Correo Argentino y, pasados los cuarenta días, la parte no confecciona el oficio ni lo diligencia: negligencia. Caducidad (art.~402): el informante, entidad pública o privada, tiene diez días hábiles para contestar; si no lo hace y la parte, dentro de los cinco días posteriores, no pide la reiteración, la prueba caduca de pleno derecho, sin importar la culpa.
\end{ejemplo}

\subsection{Valoración final: la sana crítica}

La sana crítica supone un juez libre de errores y vicios, que se convence de cómo sucedieron los hechos según las reglas sistematizadas en el art.~386, para fijarlos, subsumir la norma y sentenciar. % TODO: verificar referencia presente en las fuentes originales: las reglas de la sana crítica no se detallan en el apunte.

\begin{importante}
Los hechos son fuente y fundamento de la pretensión, fuente de prueba y fundamento de la sentencia: entre los hechos y la prueba gira lo esencial del proceso de conocimiento, pues en virtud de los hechos el derecho nace, se modifica y se extingue.
\end{importante}

\begin{cornell}{teoría general de la prueba}
\cqa{¿Qué es la prueba y qué se prueba?}{Demostración en juicio de la ocurrencia de un suceso; en el proceso civil se prueban las afirmaciones de las partes sobre hechos, no los hechos en sí.}
\cqa{¿Qué requisitos deben tener los hechos?}{Ser afirmados, conducentes, controvertidos y admisibles.}
\cqa{¿Cómo se prueba un hecho negativo?}{Mediante el afirmativo contrario o distinto; los negativos indefinidos son prueba diabólica y llevan a la carga dinámica.}
\cqa{¿Qué presunciones existen?}{\textit{Iuris tantum} (admiten prueba en contrario) e \textit{iuris et de iure} (no la admiten).}
\cqa{¿Cuál es la finalidad de la prueba?}{Formar el convencimiento del juez mediante una verdad formal que se acerque a la verdad jurídica objetiva.}
\cqa{¿Quién debe probar (art.~377)?}{Quien afirma un hecho controvertido o un precepto que el juez no debe conocer; cada parte prueba el presupuesto de la norma que invoca. Rige la adquisición probatoria.}
\cqa{¿En qué se distinguen pertinencia y admisibilidad?}{Pertinencia: el hecho pertenece al proceso por haber sido afirmado. Admisibilidad: idoneidad y licitud del medio de prueba.}
\cqa{¿Qué diferencia hay entre negligencia y caducidad de la prueba?}{La negligencia es subjetiva, a pedido de parte y por incidente; la caducidad es objetiva, de oficio o a pedido y sin sustanciación.}
\cornellsummary{La prueba reconstruye las afirmaciones sobre hechos conducentes y controvertidos; la carga recae en quien afirma, los medios deben ser pertinentes y admisibles, el plazo es de 40 días y el juez valora por sana crítica. Negligencia y caducidad sancionan la inactividad probatoria.}
\end{cornell}
\chapter{La prueba documental y la prueba informativa}
\label{cap:documental}

Como todo contrato de la vida cotidiana necesita un «papel que lo respalde», el proceso necesita la prueba documental para demostrar los hechos. Este capítulo recorre desde el ticket de estacionamiento hasta la escritura pública, y desde el correo electrónico hasta la firma digital, y estudia el valor probatorio de cada categoría, su ofrecimiento e impugnación y la prueba informativa. Presupone los conceptos de carga, pertinencia y admisibilidad del capítulo~\ref{cap:prueba}.

\section{Concepto y clasificación del documento}

El Código Procesal prevé como medios de prueba tipificados la documental, la informativa, la confesional, la testimonial, la pericial y el reconocimiento judicial. Este capítulo trata los dos primeros.

\begin{definicion}{Documento (sentido lato)}
Todo elemento susceptible de representar una manifestación del pensamiento, con prescindencia de la forma en que esté plasmada o registrada.
\end{definicion}

Tradicionalmente se lo asociaba a un objeto corpóreo y escrito; hoy comprende registros informáticos en formato binario, sin soporte necesariamente material (la nube, un soporte magnético). El Código Civil y Comercial (2015) «aggiornó» la materia con la regulación de instrumentos públicos, privados y particulares y de los documentos digitales o electrónicos, mientras que el Código Procesal, pensado desde el viejo Código Civil, no se actualizó; por eso todo manual que no contemple la reforma de 2015 está desactualizado en prueba documental. El CCyC regula también la \textbf{correspondencia} (arts.~318 y 319) y los \textbf{libros de comercio}, tras la derogación del Código de Comercio.

\begin{table}[H]
\centering
\renewcommand{\arraystretch}{1.25}
\begin{tabularx}{\textwidth}{>{\RaggedRight\arraybackslash}p{0.3\textwidth}>{\RaggedRight\arraybackslash}X}
\toprule
\textbf{Soporte} & \textbf{Categorías} \\
\midrule
Papel & Instrumentos públicos e instrumentos particulares (con o sin firma). \\
Electrónico & Con firma digital, con firma electrónica y sin firma. \\
Otros regulados & Correspondencia y libros contables. \\
\bottomrule
\end{tabularx}
\end{table}

\section{Documentos en soporte papel}

\subsection{Instrumentos públicos}

El art.~289 CCyC comprende las escrituras públicas y sus copias o testimonios, los instrumentos extendidos por escribanos o funcionarios públicos con los requisitos legales y los títulos emitidos por el Estado nacional, provincial o la Ciudad Autónoma de Buenos Aires (o municipal, según otra formulación de las notas). % TODO: verificar consistencia entre fuentes originales: el texto del art. 289 transcripto menciona «escribanos o funcionarios públicos» y «Ciudad Autónoma de Buenos Aires», y difiere levemente de la enumeración «oficial público; estado nacional, provincial o municipal».

Hacen \textbf{plena fe} hasta que una decisión judicial los declare falsos. Según el art.~296 CCyC:

\begin{table}[H]
\centering
\renewcommand{\arraystretch}{1.25}
\begin{tabularx}{\textwidth}{>{\RaggedRight\arraybackslash}p{0.34\textwidth}>{\RaggedRight\arraybackslash}X>{\RaggedRight\arraybackslash}X}
\toprule
\textbf{Aspecto} & \textbf{Fe pública} & \textbf{Cómo se impugna} \\
\midrule
Realización del acto, fecha, lugar y hechos cumplidos por el oficial público o ante él & Plena fe, hasta que sea declarado falso en juicio civil o criminal. & Redargución de falsedad (incidente) o acción civil o penal. \\
\addlinespace
Contenido de las declaraciones de las partes (convenciones, pagos, reconocimientos, enunciaciones directamente relacionadas con el objeto) & No hace plena fe: admite prueba en contrario. & Cualquier prueba en el juicio principal. \\
\bottomrule
\end{tabularx}
\end{table}

Por ejemplo, el escribano da plena fe de que las partes manifestaron que antes se entregó una suma o la posesión del inmueble, pero no de que esa entrega fue real.

\subsection{Instrumentos privados}

Es el instrumento particular \textbf{firmado}. El CCyC incorporó a la firma datos ya presentes en la doctrina y jurisprudencia: no es necesariamente la escritura del nombre, sino los signos habituales que la persona usa como marca distintiva; se equipara la impresión dígito-pulgar; y se regula la \textbf{firma a ruego}, estampada por otra persona a pedido de quien no sabe o no puede escribir. En los actos bilaterales debe haber \textbf{tantos ejemplares como partes}.

Su valor probatorio es sólido \textbf{en la medida en que esté reconocido o autenticado}: según el art.~356 el demandado debe expedirse sobre los documentos del actor y el incumplimiento los tiene por reconocidos (incluidas las cartas atribuidas). Reconocido, su valor puede ser incluso mayor que el del instrumento público, pues solo se lo impugna por vicios del consentimiento al reconocerlo.

\begin{itemize}
\item \textbf{Cotejo y pericia caligráfica.} Si quien supuestamente firmó niega la firma, se cotejan firmas por \textbf{pericia caligráfica} (art.~390, que remite al art.~458 y ss.). El perito necesita documentos \textbf{indubitados} (que pertenecen indudablemente al puño y letra de la persona) y el documento \textbf{dubitado}. Si las partes no acuerdan cuáles son los indubitados, el juez puede ordenar un \textbf{cuerpo de escritura}: en presencia del perito, el secretario y las partes se dicta a la persona palabras que escribe y firma en una hoja en blanco. Si la persona falleció (caso frecuente en testamentos ológrafos) se recurre a registros con firmas indubitadas: cuentas bancarias, Registro Nacional de las Personas, escrituras públicas y registros policiales.
\item \textbf{Abuso de firma en blanco (art.~315 CCyC).} Para impugnar un instrumento por abuso de firma en blanco se requiere un \textbf{principio de prueba por escrito}, además de otras pruebas: no bastan testimonios solos. Un \textbf{contradocumento} puede ser sustancial, pero no alcanza por sí solo.
\item \textbf{Impresión dígito-pulgar.} Se la considera principio de prueba, impugnable por cualquier medio; como suele firmar así quien no sabe leer y escribir y puede ser inducido a firmar algo distinto, es preferible la firma a ruego.
\end{itemize}

\subsection{Instrumentos particulares sin firma}

Son los extendidos por escrito sin firma: el comprobante de una playa de estacionamiento con los tres últimos números de la patente, vouchers, tickets, cupones, boletos. Según el art.~319 CCyC, el juez aprecia su valor ponderando la congruencia entre lo sucedido y narrado, la precisión y claridad técnica, los usos y prácticas del tráfico, las relaciones precedentes y la confiabilidad de los soportes y procedimientos técnicos (códigos de barras, por ejemplo). % TODO: verificar consistencia entre fuentes originales: los apuntes atribuyen el texto del art. 319 a los instrumentos particulares y mencionan los arts. 318 y 319 al tratar la correspondencia.

\section{Documentos electrónicos}

\begin{table}[H]
\centering
\renewcommand{\arraystretch}{1.25}
\begin{tabularx}{\textwidth}{>{\RaggedRight\arraybackslash}p{0.24\textwidth}>{\RaggedRight\arraybackslash}p{0.22\textwidth}>{\RaggedRight\arraybackslash}X}
\toprule
\textbf{Documento} & \textbf{Equivalente en papel} & \textbf{Característica} \\
\midrule
Con firma digital & Instrumento público & Certificación estatal; plena fe de autenticidad. \\
Con firma electrónica & Instrumento privado & Sistema de claves; la carga de probar recae en quien lo presenta. \\
Sin firma & Instrumento particular & Videos, audios, registros; valoración judicial libre. \\
\bottomrule
\end{tabularx}
\end{table}

\begin{table}[H]
\centering
\renewcommand{\arraystretch}{1.25}
\begin{tabularx}{\textwidth}{>{\RaggedRight\arraybackslash}p{0.22\textwidth}>{\RaggedRight\arraybackslash}X>{\RaggedRight\arraybackslash}X}
\toprule
 & \textbf{Firma digital (Ley 25.506)} & \textbf{Firma electrónica} \\
\midrule
Certificación & Organismo estatal certificante; el firmante obtuvo su clave allí. & Sin entidad pública certificante; claves que permiten identificar por cruce de datos. \\
Plena fe & Sí, de la autenticidad del contenido y la identidad del emisor. & No; la parte que la presenta debe demostrar la seriedad del sistema. \\
Impugnación & Quien impugna prueba técnicamente (pericia informática) un defecto en la certificación, emisión o vigencia. & Cualquier medio para negar autenticidad. \\
Uso típico & Trámites estatales y documentos oficiales. & Homebanking, MercadoPago, compras con tarjeta y presentaciones en el sistema de gestión judicial. \\
\bottomrule
\end{tabularx}
\end{table}

La Ley 25.506 define la firma digital como el resultado de aplicar a un documento digital un procedimiento matemático que requiere información de exclusivo conocimiento del firmante, bajo su absoluto control, y que permite a terceros verificarla, identificar al firmante y detectar toda alteración posterior. Los expertos sostienen que tiene menor falibilidad que los instrumentos públicos. Si el sistema de firma electrónica es endeble, difícilmente pueda admitirse una firma que no garantice identidad ni una plausible expresión de voluntad. El sistema judicial usa firma electrónica: los abogados presentan escritos con ella y los jueces validan providencias con una clave alfanumérica individual.

Los documentos electrónicos sin firma (registros visuales, auditivos, videograbados) son hoy de acceso inmediato por los teléfonos inteligentes. Google Street View, por ejemplo, permite ver la evolución de un lugar por las distintas «pasadas» de la cámara (el Paseo del Bajo de la Ciudad ha cambiado significativamente en dos o tres años). Su validez depende de la calidad técnica y de otros elementos que permitan valorarlos congruentemente; en última instancia, lo valora el juez.

\section{Correspondencia y libros contables}

\subsection{Correspondencia}

En papel o electrónica, se rige por los arts.~318 y 319 CCyC: puede presentarla el destinatario, pero la confidencial requiere asentimiento del remitente; los terceros no pueden valerse de ella sin asentimiento del destinatario y, si es confidencial, del remitente. Reglas: (a) el destinatario puede aportarla si no es confidencial; (b) si lo es, necesita conformidad del remitente, pues si no se violaría la inviolabilidad de la correspondencia (art.~18 CN); (c) la de contenido comercial o negocial (tratativas contractuales) no presenta inconveniente; (d) si la presenta el remitente, el destinatario debe reconocer o negar la recepción, no la confección; (e) la de terceros se autentica por la testimonial de quienes se las atribuye. Reconocida vale como plena prueba, igual que el instrumento privado; si el tercero comparece como testigo y reconoce haberla confeccionado, vale como indicio según las circunstancias.

\subsection{Libros contables}

Son documentos privados de elaboración unilateral cuyo valor radica en llevarse con las formalidades legales y estar legalizados por la autoridad. Según el art.~330 CCyC, la contabilidad llevada en forma debe admitirse en juicio como medio de prueba, sus registros prueban contra quien la lleva y el juez puede apreciarla y exigir prueba supletoria o complementaria. Rige la \textbf{indivisibilidad}: no pueden valorarse ciertos asientos y descartarse otros.

\begin{table}[H]
\centering
\renewcommand{\arraystretch}{1.25}
\begin{tabularx}{\textwidth}{>{\RaggedRight\arraybackslash}p{0.5\textwidth}>{\RaggedRight\arraybackslash}X}
\toprule
\textbf{Situación} & \textbf{Valor probatorio} \\
\midrule
Ambas partes llevan contabilidad y los asientos perjudican a quien los lleva & Plena prueba en su contra. \\
Una parte presenta libros y la otra, obligada a llevarlos, no & Se valoran favorablemente para quien los exhibe. \\
Ambas presentan libros contradictorios & El juez prescinde de ellos y se ciñe al resto de la prueba. \\
Solo una parte lleva contabilidad (la otra no es comerciante) & Principio de prueba, evaluado con el resto. \\
\bottomrule
\end{tabularx}
\end{table}

Quien postula un argumento debe guardar coherencia con sus propios libros, pues lo que surja en contra queda desvirtuado ante el juez.

\section{Ofrecimiento e impugnación de la prueba documental}

\subsection{Oportunidad y documentos en poder de otros}

El momento prácticamente exclusivo es la demanda, la reconvención y la contestación de ambas (art.~333): se acompañan los documentos en poder de la parte; los demás se individualizan, indicando contenido y lugar, archivo, oficina pública o persona que los tiene. Excepción (art.~335): los de fecha posterior o los que se jura no haber conocido antes, con traslado a la otra parte para que cumpla el art.~356 inc.~1; conviene acompañar alguna prueba de que llegaron a poder de la parte después.

Si el documento está en poder de la contraparte se pide la intimación a presentarlo (art.~388); no se puede constreñir a presentar un documento cuya existencia no es segura, y si por otros elementos resulta manifiestamente verosímil su existencia y contenido, la negativa es una \textbf{presunción en contra}. El proceso civil no permite el allanamiento para buscarlo. Los terceros pueden ser requeridos pero pueden oponerse si les causa perjuicio; si presentan documentos tienen derecho a copias certificadas y a su pronta devolución. De todo documento presentado se corre traslado y la contraria debe expedirse sobre su autenticidad.

\subsection{Impugnación del instrumento público (art.~395)}

El mero desconocimiento no basta para tachar de falso un instrumento público. Se requiere: (1) impugnarlo en el plazo del traslado de la documentación y (2) promover el \textbf{incidente de redargución de falsedad} dentro de los 10 días hábiles judiciales desde la impugnación, bajo apercibimiento de tenerlo por reconocido. % TODO: verificar consistencia entre fuentes originales: la unidad sobre contestación de demanda indica 15 días desde que se tachó de falso el documento; esta unidad indica 10 días.
Son partes el oficial público que intervino y las partes del acto (o sus sucesores); tramita en cuerda separada y se resuelve junto con la sentencia definitiva. La tacha puede ejercerse también por acción civil o querella penal autónomas. Cuando se impugna el contenido de las declaraciones de las partes ante el oficial (no el acto, la fecha o el lugar), \textbf{no es necesaria la redargución}: cualquier prueba en el juicio principal puede desvirtuarlas.

\section{La prueba informativa}

Es un modo de aportar al proceso hechos o actos \textbf{registrados} en documentos, archivos o registros de terceros ajenos a las partes, sean entidades públicas o privadas.

\begin{importante}
La prueba informativa no obtiene opiniones ni versiones sobre hechos (eso es la prueba testimonial), sino datos registrados.
\end{importante}

Por ejemplo, un oficio al SAME para que informe si registró la atención brindada el 14 de marzo de 2020 por un accidente en Rivadavia al 5700, a quiénes atendió y adónde los trasladó.

\subsection{El oficio (arts.~398 y 400)}

Debe indicar a quién va dirigido; las actuaciones y el juzgado en que tramitan; el domicilio del juzgado; el nombre del juez y del secretario (y un teléfono); una descripción detallada de lo que se requiere informar; la transcripción de la orden del juez que lo autoriza; y las sanciones por incumplimiento. Según el art.~398, las entidades que no contesten en diez días son pasibles de sanciones conminatorias progresivas. Según el art.~400, los pedidos de informes pueden ser firmados directamente por el abogado, con la sola acreditación de su calidad de tal; los oficios o exhortos al extranjero se ajustan a los tratados o a la costumbre internacional. En cambio, según el art.~38, los dirigidos al Presidente de la Nación, ministros, secretarios de Estado u otros jueces o magistrados deben ser firmados por el \textbf{juez de la causa}, y se dejan en secretaría para su firma.

\begin{ejemplo}{Modelo de oficio}
«Tengo el agrado de dirigirme a Usted, en los autos caratulados Fulano c/ Mengano sobre daños y perjuicios, que tramitan ante el Juzgado Nacional de Primera Instancia en lo Civil N.º 18 a cargo del Dr.\ Fulanito, Secretaría N.º [tal] a cargo del Dr.\ Menganito, a efectos de requerirle tenga bien informar a este Juzgado acerca de la atención brindada a la Sra.\ [Artal] en tal fecha en Rivadavia al 5700, detallando el diagnóstico que presentó y el lugar al que fue trasladada. La providencia que ordena la medida dice: Buenos Aires, 15 de febrero del 2020: Líbrese el oficio solicitado por la parte actora.»
\end{ejemplo}

\subsection{Plazo y caducidad (art.~402)}

La entidad tiene \textbf{10 días hábiles judiciales} para contestar; el abogado acredita el diligenciamiento con copia sellada. % TODO: verificar consistencia entre fuentes originales: los apuntes cuentan el plazo «desde el día siguiente a la presentación del oficio» y, en el cuadro de caducidad, «desde la recepción del oficio».
Vencido el plazo sin respuesta, se tiene por desistida la prueba a la parte que la pidió, sin sustanciación, si dentro del quinto día no solicita la reiteración del oficio.

\begin{table}[H]
\centering
\renewcommand{\arraystretch}{1.25}
\begin{tabularx}{\textwidth}{>{\RaggedRight\arraybackslash}p{0.27\textwidth}>{\RaggedRight\arraybackslash}p{0.27\textwidth}>{\RaggedRight\arraybackslash}X}
\toprule
\textbf{Etapa} & \textbf{Plazo} & \textbf{Consecuencia} \\
\midrule
La entidad contesta & 10 días hábiles & Si no contesta, corre el plazo para reiterar. \\
La parte pide reiteración & 5 días desde el vencimiento & Si no lo hace, caduca la prueba informativa. \\
\bottomrule
\end{tabularx}
\end{table}

\begin{importante}
Al comenzar a ejercer, es fundamental estar atento a los plazos: en ciertas circunstancias puede ser muy grave para los clientes.
\end{importante}

\subsection{Modernización}

La administración pública ha comenzado a incorporar sistemas electrónicos de pedido de informes: el Banco de la Nación Argentina, el Banco Provincia de Buenos Aires, el Banco Central, el Registro Nacional de la Propiedad Automotor, el Registro de la Propiedad Inmueble, el Registro Nacional de las Personas, la Dirección Nacional de Migraciones, varios bancos privados y las comisarías de la Ciudad y la Policía Federal. La cuarentena incrementó el tráfico electrónico. Pero el oficio en papel sigue siendo la única vía para muchas pequeñas y medianas empresas y entidades privadas con poca comunicación con los tribunales.

\begin{cornell}{prueba documental e informativa}
\cqa{¿Qué es un documento y qué cambió en 2015?}{Todo elemento que representa una manifestación del pensamiento, sin importar el soporte. El CCyC regula instrumentos públicos, privados, particulares y electrónicos; el Código Procesal no se actualizó.}
\cqa{¿Qué parte del instrumento público hace plena fe (art.~296)?}{La realización del acto, fecha, lugar y hechos cumplidos por el oficial o ante él. Las declaraciones de las partes admiten prueba en contrario.}
\cqa{¿Cuándo vale el instrumento privado y cómo se coteja la firma negada?}{Cuando está reconocido o autenticado; se coteja por pericia caligráfica con documentos indubitados o un cuerpo de escritura.}
\cqa{¿Qué diferencia hay entre firma digital y electrónica?}{La digital tiene certificación estatal y plena fe (quien impugna prueba técnicamente); la electrónica no, y quien la presenta debe probar la seriedad del sistema.}
\cqa{¿Cuándo puede usarse la correspondencia y cómo se valoran los libros contables?}{La confidencial requiere asentimiento del remitente. Los libros prueban contra quien los lleva, se valoran íntegramente y pueden quedar descartados si ambos son contradictorios.}
\cqa{¿Cómo se ofrece prueba documental y qué pasa si está en poder de la contraria?}{Con la demanda y la contestación (art.~333), salvo documentos posteriores (art.~335); la contraria es intimada (art.~388) y su negativa, si hay verosimilitud, es presunción en contra.}
\cqa{¿Cómo se impugna un instrumento público?}{Con impugnación en el traslado y luego incidente de redargución de falsedad (art.~395); no es necesario para atacar las declaraciones de las partes.}
\cqa{¿Qué es la prueba informativa y cuándo caduca?}{Datos registrados en archivos de terceros, por oficio; 10 días para contestar y 5 más para pedir reiteración (art.~402), si no caduca.}
\cornellsummary{La prueba documental valora cada categoría según su fe pública, autenticación o certificación, y se ofrece y se impugna en oportunidades y formas estrictas. La prueba informativa trae al proceso datos registrados por terceros mediante oficios sujetos a plazos de contestación y caducidad.}
\end{cornell}

\backmatter
\chapter*{Bibliografía}
\addcontentsline{toc}{chapter}{Bibliografía}
\markboth{Bibliografía}{Bibliografía}

Material de estudio y obras mencionadas en las fuentes:

\begin{itemize}
\item Libro de la cátedra: \textit{Nociones Básicas del Derecho Procesal Civil} (en prensa).
\item Palacio.
\item Falcón.
\item Arazi.
\item Código Procesal Civil y Comercial de la Nación (edición con leyes complementarias).
\item Linares, Juan Francisco. \textit{La razonabilidad de las leyes}.
\item Entelman, Remo. \textit{Teoría del Conflicto}.
\item Calvo Soler, Raúl. \textit{Mapeo de Conflictos}.
\item Fisher, R. y Ury, W. Método de Harvard sobre negociación.
\item Rojas, Jorge A. Trabajos sobre arbitraje, en particular \textit{La pericia arbitral como otro modo de conclusión del proceso}, y «Réquiem para el arraigo» (\url{www.rojasomar.com.ar}; otra unidad cita \url{www.jorgearojas.com.ar}). % TODO: verificar consistencia entre fuentes originales: dos direcciones web distintas para el sitio de la cátedra.
\end{itemize}

\end{document}
